%&pdflatex
\documentclass [12pt]{article}
\usepackage{authblk}
\usepackage{silence}
\usepackage{amsfonts}
\usepackage{amsmath}
\usepackage{wasysym}
\usepackage{hhline}
\usepackage{cite}
\usepackage{tocloft}
\usepackage[plainpages=false,pdfpagelabels]{hyperref}
\usepackage[nottoc]{tocbibind}

\WarningFilter{latexfont}{Font shape `U/wasy/b/n'}
\WarningFilter{latexfont}{Font shape `OMS/cmtt/m/n'}
\WarningFilter{latexfont}{Font shape `OMS/cmtt/bx/n'}

\setlength{\topmargin}{-2cm}
\setlength{\oddsidemargin}{0cm}
\setlength{\textheight}{23cm}
\setlength{\textwidth}{16cm}
\setlength{\parindent}{0pt}
\setlength{\parskip}{1.2pt}

\makeatletter
\renewcommand{\tt}{\begingroup\catcode`_=12\def\_{\string_}\tt@}
\newcommand{\tt@}[1]{\texttt{#1}\endgroup}
\newcommand{\ttsubsec}[1]{\subsection {\texorpdfstring{\tt{#1}}{#1}}}
\newcommand{\tr}{\mathop{\mathrm{tr}}}
\makeatother

\title
    {
    Gordion: Design and Implementation Notes
    \\[8pt]
    \large Version 2
    }
\date	{\small\today\vspace*{-30pt}}
\author	{Laurence G.~Yaffe}
\affil{%
    \small
    Department of Physics, University of Washington, Seattle WA 98195, USA
    \vspace*{-10pt}
    }%
\pagenumbering{roman}
\begin	{document}
\maketitle
\tableofcontents
\newpage
\setlength{\parskip}{5pt plus 2pt minus 1pt}
\pagenumbering{arabic}

\section {Introduction}

\tt{Gordian} is a program for solving, numerically,
for low energy properties of
a variety gauge theories in the large $N$ limit using the
coherent state variational method.%
\footnote
    {%
    For background on this approach for solving large $N$ limits,
    see Refs.~\cite{YaffeRMP, CSVAI, CSVAII,Yaffe:2026ouz}.
    Appendix \ref{sec:errata} lists corrections for
    a number of misprints in these references.
    Prior work on this approach is also presented
    in the Princeton University Ph.D.~theses
    of F.~R.~Brown \cite{Brown:1987km},
    U.~Lindqwister \cite {Lindqwister:1988xc},
    T.~Dickens \cite{Dickens:1987ih},
    and W.~Somsky~\cite{Somsky:1989}.
    }
The name \tt{Gordian} is an allusion to the Greek legend
associated with Alexander the Great, who
is reputed to have sliced through a horrendously complicated
knot, instead of carefully untying it,
in Gordion (Latin: Gordium, Phrygian: Gordum),
the capital city of ancient Phrygia.
Dealing with arbitrarily complicated loops is at the
core of the coherent state variational method as
implemented in \tt{Gordian}.
It is hoped that this program helps cut through
the complexity of large $N$ gauge theories
{\large\smiley{}}.

\setlength{\parskip}{3pt}

\section {Overall structure}

The key steps comprising the coherent state variation algorithm
\cite {CSVAII} are:
\begin{enumerate}\advance\itemsep -4pt
\item
    Construction of a list of all observables of the chosen theory
    to be retained in the finite truncation set,
    and whose expectation values will be computed.
    Here ``observables'' mean Wilson loops, possibly containing one or
    two electric field insertion (in Hamiltonian theories),
    or fermion bilinears, possibly including one electric field insertion.
\item
    Selection of the finite set of coherence algebra generators, for a given theory,
    which will be used to generate deformations of a coherent state,
    and whose coefficients will serve as Riemann normal coordinates
    in the neighborhood of any given point
    in the large $N$ phase space.
    These are anti-Hermitian combinations of Wilson loops,
    possibly with one electric field insertion, or fermion bilinears.
\item
    Evaluation, symbolically, of the commutators of all selected observables
    and generators.  These define the geodesic equations
    on the large $N$ phase space and encode how observable expectation
    values vary as one deforms the state.
\item
    Evaluation, symbolically, of the commutators defining the first and second
    variations of the Hamiltonian (or Euclidean free energy) with respect
    to Riemann normal coordinates.
    Also, for Hamiltonian theories,
    evaluation of commutators of generators whose expectation values
    define the Lagrange bracket (the inverse Poisson bracket)
    of the large $N$ phase space.
\item
    Numerical minimization of the Hamiltonian (or Euclidean
    free energy) using Newton iteration, with each iterative step
    involving numerical evaluation of the symbolic expressions
    for the gradient and curvature of the Hamiltonian (or free energy),
    solution of the linear equations predicting the location of the minimum,
    and numerical integration of the geodesic equations describing
    the change in expectation values of all observables as one moves
    to the predicted minimum.
\item
    In Hamiltonian theories,
    numerical evaluation of the symbolic expressions for the curvature of the
    Hamiltonian and the Lagrange bracket, in any chosen symmetry channel,
    and computation of the resulting
    small oscillation frequencies around the minimum,
    yielding a determination of the
    low-lying glueball or meson spectrum.
\end{enumerate}
%
Each of these steps is carried out first for the $O(N^2)$ pure gauge dynamics,
for some sequence of gauge coupling values,
and then again for the sub-leading $O(N)$ fundamental representation
fermion dynamics,
for some sequence of fermion mass values.

\section {Program scope}

In early work  on the coherent state variation method in the
1980's \cite{CSVAII},
separate C programs were written to handle different aspects
of the overall procedure:
symbol definition (the encoding used to represent Wilson loops,
loops decorated with electric field insertions,
and matter field bilinears),
observable generation
(the recursive construction of the truncated
set of loops and bilinears used to encode
properties of a coherent state),
polynomial script generation
(the construction of the geodesic equations
for retained observables, as well the gradient
and curvature of the Hamiltonian or Euclidean free energy),
the actual minimization numerics
(iteratively performing Newton minimization, repeatedly
evaluating the gradient and curvature to predict the minimum location
and integrating the geodesic equations to move to the
predicted minimum),
and computing small oscillations around the
minimum to extract glueball or meson spectra.
Given the constraints on available memory at the time,
this partitioned approach with explicit passing of
information between stages via large disk files was essential.

For the current implementation, the design goals were to:
\begin{itemize}\advance\itemsep -4pt
\item
    Create a unified program that can handle all needed
    steps from beginning to end.
\item
    Design a program capable of handling U($N$) gauge theories%
    \footnote
	{%
	Note that the difference between U($N$) and
	SU($N$) gauge theories first appears at relative
	order $1/N^2$ in a large-$N$ expansion;
	there is no difference in their large $N$ limits
	as discussed above.
	}
    in multiple dimensions,
    in either Hamiltonian or Euclidean lattice formulations,
    with or without fundamental representation fermions.
    If fermions are present, it is assumed that the
    number of fermion flavors remains fixed as $N \to \infty$.
\item
    Test the efficacy of possible approximation schemes for
    non-retained observables.
\item
    Write highly time- and space-efficient code.
\end{itemize}
Regarding the last point,
even through compute capabilities have vastly increased
since the 1980's,
the computational complexity of numerically
minimizing the large $N$ Hamiltonian (or Euclidean free energy)
to a given accuracy inevitably grows exponentially with the
correlation length of the theory, since the number of Wilson loops
grows exponentially with any reasonable measure of their size.
Because the continuum limit entails diverging
correlation length,
for any given level of computational resources
a more efficient implementation will allow one to
reach larger correlation lengths and hence
better control the continuum limit.

Of most physical interest are the cases of
Hamiltonian lattice Yang-Mills or QCD
\cite{Kogut:1974ag}
on lattices of two or three space dimensions, or
Euclidean lattice Yang-Mills or QCD
\cite{Wilson:1974sk}
on three or four dimensional
space-time lattices.
Including an ability to handle theories on
lower-dimensional lattices, possibly compactified,
allows one to test the program
(and the method) on various exactly soluble models including
the Euclidean one-plaquette model
\cite{Gross:1980he},
the Hamiltonian one-plaquette model
\cite{Wadia:1980cp,Jevicki:1980zq,Neuberger:1980qh},
1+1 dimensional QCD in a Hamiltonian lattice formulation
\cite{Berruto:1999ir,Berruto:2002gn,Kiskis:2002gr,GarciaPerez:2016hph}
and
two dimensional Euclidean Yang-Mills on an infinite lattice
\cite{Gross:1980he}.
Permitting one or more dimensions to be periodically compactified
also enables one to study finite temperature equilibrium states
in either a Euclidean formulation or in a Hamiltonian
formulation with a compact spatial dimension, and to study
the spectrum of strings wrapping a compactified dimension.
%
Central to the approach is the fact that
unbroken translation symmetry permits one to restrict
the coherent state minimization to translationally invariant states,
thereby allowing one to work directly in infinite volume.

These considerations motivated the following
specific program design decisions and assumptions:
\begin{itemize}

\item
    The lattice dimension may be 1, 2, 3 or 4, with each dimension
    either infinite, or periodically compactified.

\item
    Zero, one or two fermion flavors are allowed,%
    \footnote
	{%
	When fermions are present,
	any compactified dimension length must be
	an even number of lattice spacings.
	Supporting multiple fermion flavors potentially allows one to
	study possible spontaneous flavor symmetry breaking,
	perhaps induced by the inclusion of four-fermi operators.
	}
    with each fermion flavor defined as having one conjugate pair of
    fermion operators per site per flavor.
    For Hamiltonian theories, this is a staggered fermion discretization.

\item
    All theories under consideration are invariant under cubic
    lattice symmetries (translations, permutations and reflections)
    together with charge conjugation (C) and time reversal.%
    \footnote
	{%
	The Yang-Mills $\theta$ angle is zero.
	For Hamiltonian theories,
	only translations by an even number of lattice sites are 
	required to be symmetries of the theory.
	}

\item
    None of these symmetries are spontaneously broken, and
    observables of interest are invariant under all symmetries.%
    \footnote
	{%
	Performing spectrum calculations at
	non-zero momentum is a possible future addition but
	would require extending the code to evaluate
	commutators of observables carrying non-zero momentum.
	}

\item
    Use modern C++, instead of some higher level programming language,
    thereby enabling more efficient memory use and explicit
    control over performance critical aspects
    such as copy versus move semantics when passing
    and returning data structures,
    while also having strong type checking, template meta-programming,
    highly useful standard container classes, and reliable highly
    optimizing compilers.

\end{itemize}
\medskip

\section {Implementation considerations}

\subsection {Symbol encoding}\label{sec:symbols}

    Observables are represented by strings of symbols in which every
    symbol either denotes
    a possible combination of gauge fields residing on a single
    lattice link $\ell$, or
    a fermion field residing on a lattice site.
    In Hamiltonian theories,
    for each lattice dimension, there are ten possible
    combinations of gauge fields which may be combined together
    on any given lattice link:
    \begin{itemize}\advance\itemsep -3pt
    \item A forward directed link variable, $U[\ell]$;
    \item A backward directed link variable, $U[\bar\ell] \equiv U[\ell]^\dagger$;
    \item An electric field operator, $E[\ell]$ (transforming at the starting site of the link);
    \item A reflected electric field, $U[\bar\ell] E[\ell] U[\ell]$
	(transforming at the link end site);
    \item A squared electric field, $E[\ell]^2$;
    \item A reflected electric field square,
	$U[\bar\ell] E[\ell]^2 \, U[\ell]$;
    \item An electric field inserted forward link, $E[\ell] U[\ell]$;
    \item An electric field inserted backward link, $U[\bar\ell] E[\ell]$;
    \item A squared electric field inserted forward link, $E[\ell]^2 U[\ell]$;
    \item A squared electric field inserted backward link, $U[\bar\ell] E[\ell]^2$.
    \end{itemize}
    Using distinct symbols to denote each of these possible combinations of
    gauge fields which may appear together on a single link
    enables symmetry transformations to be implemented by simple symbol
    substitutions which never change the length to the symbol string
    representing an observable.

    For fermions in Hamiltonian theories, distinct symbols are needed for 
    staggered and unstaggered fermion operators, while for Eulcidean theories
    distinct symbols are needed for fermion fields and the Grassmann derivatives
    of these fields.

    To treat theories on lattices of potentially up to four dimensions,
    an individual symbol (for Hamiltonian theories)
    must be able to represent any of the above ten different combinations
    of gauge fields in every possible lattice direction,
    implying 40 possible values for gauge degrees of freedom, plus
    four additional possibilities for each fermion flavor,
    as detailed below.
    Hence, a single byte is amply sufficient to encode the range of symbols needed
    for all theories in the intended program scope.

    The specific encoding employed, defined in file \tt{Symb.h},
    was designed so that simple bitwise operators,
    rather than look-up tables,
    suffice to extract the various properties encoded by each symbol value,
    such as the lattice axis or the above symbol types (for gauge symbols),
    or the fermion flavor or staggering.
    Conversion of a symbol into its reflection, charge conjugate,
    or staggered variant, are also simple bitwise operations.
    To make the canonicalization table indexing as simple as possible,
    the smallest symbol values ($0,{\cdots},7$) are used for links.
    For more detail, see the discussion in section \ref{sec:Symb}.

    The basic \tt{symb} type used to encode individual symbols 
    is defined as a type alias for \tt{signed char},
    not as a distinct class,
    even though the latter option might seem natural.
    Using \tt{symb} as a mere alias for a \tt{char}
    allows the symbol string class \tt{SymbStr} to simply inherit
    from \tt{std::string} and automatically acquire
    the short string optimizations used in
    implementations of \tt{std::string}.

    For printing and for user input of observables or generators,
    the following 
    mapping is used to relate printable characters to symbols:
    \begin{center}
    \begin{tabular}{cp{0.8\linewidth}}
	\it character & \it associated symbol
    \\ \tt{x}, \tt{y}, \tt{z}, \tt{w} & gauge links in directions 1, 2, 3 or 4, respectively;
    \\ \tt{X}, \tt{Y}, \tt{Z}, \tt{W} & inverse gauge links in directions 1, 2, 3 or 4;
    \\ \tt{A}, \tt{B}, \tt{C}, \tt{D} & electric fields in direction 1, 2, 3 or 4;
    \\ \tt{f}, \tt{F}, \tt{g}, \tt{G} & fermion, conjugate fermion, 
		staggered fermion or Grassmann derivative $\partial/\partial \tt{F}$
		(Hamiltonian or Euclidean), and staggered conjugate fermion
		or Grassmann derivative $\partial/\partial \tt{f}$, respectively,
		all for flavor 1;
    \\ \tt{h}, \tt{H}, \tt{i}, \tt{I} & corresponding flavor 2 fermion fields;
    \\ \tt{S}, \tt{s} & gauge or fermion entropy, respectively.
    \end{tabular}
    \end{center}
    So, for example, \tt{xyXY} represents a basic plaquette in the 1-2 plane,
    \tt{YBy} is a backwards-directed electric field in the 2-direction,
    \tt{Fxf} is a one-link fermion hopping term in the 1-direction,
    and \tt{Gf} is a staggered fermion mass term (for Hamiltonian theories).

\subsection {Observable storage}

C++ standard library strings (\tt{std::string}) are used to
store the symbol strings representing observables.
More precisely, the class \tt{Obs}, whose instances 
represent physical observables, inherits from \tt{std::string}.
This use of \tt{std::string}
enables the \tt{Obs} class to automatically
take advantage of the short string optimizations present
in implementations of \tt{std::string} (allowing sufficiently
short strings to be stored directly in the class object
without needing any free store allocation/deallocation operations).

For storing the entire set of generated observables,
one needs some form of bidirectional data structure which:
\begin{enumerate}\advance\itemsep -4pt
\item
    Enables highly efficient determination of whether a newly
    constructed observable is already present in the set;
\item
    Assigns a unique integer index to each observable in the set;
\item
    Enables easy access to an observable's symbol string given
    its index in the set.
\end{enumerate}
The first goal is best accomplished by using a hash table
(or \tt{std::unordered_map}), as the typical performance of
hash table lookups greatly surpasses that of alternatives like
balanced binary trees.
To support the other requirements, the observable list data structure
(termed an \tt{ObsList}) combines a hash table with a list of pointers
to the individual observable objects stored in the hash table.
This avoids duplicate storage of the symbol strings of all the
stored observables.
This approach relies on the fact that pointers to individual entries
in an \tt{unordered_map} are not invalidated even when a rehashing
takes place.

\subsection {Observable canonicalization}

Observable canonicalization refers to the
transformation of an observable's symbol string
into a unique physically equivalent representation.
This may involve application of any lattice symmetry,
possibly combined with charge conjugation and,
for gauge observables,
cyclically rotating the symbol string to a different
starting point.
With 96 elements of the lattice symmetry group (times $C$)
in three dimensions, or 768 in four dimensions,
typical observables will have
thousands of physically equivalent representations.
It is essential that the transformation of
observables into some unique canonical representation
be as efficient as possible.

The strategy employed treats loop observables
(those only containing gauge links)
separately from ``special'' observables containing
electric field insertions and/or fermions.
For loop observables,
the canonical form is defined as the first
representative when all physically equivalent
symbol strings for a loop are lexically ordered.
But actually constructing and sorting the set of 
all physically equivalent forms, every time a
loop is to be canonicalized, is hugely inefficient.
The following procedure is used instead.

A compile-time constant \tt{loopchunksize},
whose value depends on the lattice dimension,
serves to define ``short'' loops as those no longer
than \tt{loopchunksize}.
Short loops are used directly as keys into
a canonicalization cache, which is a hash
table mapping observables to the index
of their physically equivalent canonicalized
form in the master list of canonicalized observables.
For loop observables longer than \tt{loopchunksize},
a moving window of \tt{loopchunksize} is walked around
the loop, with each such ``chunk'' converted to an
index into a static loop canonicalization table,
each of whose entries specifies the subset of
symmetry transformations which would transform
this chunk into its preferred (lowest lexically ordered)
form, and a score, based on that lexical ordering, of
the resulting preferred form of the chunk.
As the chunk window is walked around the loop,
the lowest scoring chunk(s) identify the 
small subsets of symmetry transformations
and starting positions which
need to be compared to identify the preferred
form for the full loop.

For gauge observables containing electric field insertions,
the canonical form is defined as the lowest
lexically ordered form among those physically equivalent
representations which begin with an electric-field
inserted symbol.
Since observables will have at most two electric field insertions,
this greatly reduces the number of possible starting points
which must be compared.
To facilitate canonicalization and commutation of
gauge observables containing electric field insertions,
it is required throughout the program that the
symbol string for any gauge observable containing electric
field insertions begin with an electric field inserted symbol.
Similarly, symbol strings for fermion bilinears always
begin with a conjugate fermion symbol and end with a fermion.
So ``special'' observables always begin with a non-link symbol.

The resulting canonicalization procedure for such special observables
is similar to that for loops, except that only transformations which
lead to acceptable symbol strings starting with an appropriate non-link
symbol need be considered.
A different compile-time constant \tt{specchunksize} defines
the separation between ``short'' and ``long'' special observables.
The canonicalization cache is used for special observables whose
number of (non-fermion) symbols is no more than \tt{specchunksize}.
For longer observables, \tt{specchunksize} chunks
which either begin or end at the endpoints of a fermion bilinear,
or on electric field insertions in gauge loops,
are converted into indices into a static non-loop canonicalization
table, each of whose entries specifies the subset of
symmetry transformations which would transform
this chunk into its preferred form, and a score,
based on lexical ordering, of the resulting preferred form of the chunk.
The scores of only two such entries need be compared
for fermion bilinears, one for the initial chunk of the symbol string
and one for the final chunk (involving symmetries which
include charge conjugation).
At most four such entries need be compared for
gauge observables with two electric field insertions.

The values of \tt{loopchunksize} and \tt{specchunksize}
are chosen to balance opposing goals:
larger chunks lead to smaller sets of possible symmetries which
will need to be considered in fully canonicalizing an observable,
but at the cost of requiring larger pre-computed canonicalization
tables and a larger cache for ``short'' observables.
If the canonicalization tables are overly large, repeated frequent
lookups into those tables will seriously compromise processor cache
performance.
These considerations lead to the definition of
\tt{loopchunksize} as 9, 7 or 6 for lattices
of dimension 2, 3 or 3, respectively, which in turn
leads to loop canonicalization tables
with between 262K and 280K possible entries.
Similarly, depending on lattice dimension \tt{specchunksize}
is defined to be 4 or 5 in Hamiltonian theories, or 6, 7 or 8
in Euclidean theories (which lack electric-field inserted
observables),
leading to  non-loop canonicalization table sizes ranging from
41K up to 524K for the various theory classes..
In practice, this canonicalization procedure enables the
number of symmetry transformations performed per canonicalization
to average below one in 2D to three in 4D.

\advance\abovedisplayskip -2pt
\advance\belowdisplayskip -2pt
\advance\abovedisplayshortskip -2pt
\advance\belowdisplayshortskip -2pt

\subsection {Commutator evaluation}

The basic U($N$) Hamiltonian lattice theory commutation relations are:
\begin{subequations}\label{eq:YMcomms}
\begin{align}
    [ (E_\ell)_{ij},\, (U_{\ell'})_{kl}]
    &=
    \frac {1}{N} \> \delta^{\ell \ell'} \, \delta_{kj} \, (U_\ell)_{il} \,,
\\
    [ (E_\ell)_{ij},\, (E_{\ell'})_{kl}]
    &=
    \frac {1}{N} \> \delta^{\ell \ell'}
    \left(
	\delta_{kj} \, (E_\ell)_{il} - \delta_{il} \, (U_\ell)_{kj}
    \right),
\end{align}
\end{subequations}
for the gauge degrees of freedom living on links, together with the
anti-commutation relation for fermions residing on sites,
\begin{align}
    \{ (\psi_s)^\alpha_i, \,  (\psi^\dagger_{s'})^\beta_j \}
    =
    \delta_{s s'} \, \delta_{ij} \, \delta^{\alpha\beta} \,.
\end{align}
Here, $\ell$ labels links, $s$ labels sites,
$i$, $j$, $k$, $l$ are U($N$) gauge indices,
and $\alpha$, $\beta$ denote fermion flavors.
The Yang-Mills commutation relations (\ref{eq:YMcomms})
may be realized in a ``coordinate representation'' for which
\begin{align}
    (E_\ell)_{ij} \equiv \frac 1N \> (U_\ell)_{ik} \,\frac {\partial}{\partial (U_\ell)_{jk}} \,.
\label{eq:Erep}
\end{align}
For Euclidean theories, ``electric field insertions'' in
coherence group generators are defined to mean the first order differential
operators (\ref{eq:Erep}).

As discussed in Ref.~\cite{CSVAII}, the encoding of observables
as strings of symbols, one per lattice link (or site), enables
the evaluation of commutators of generators and observables
to be converted, modulo a few details,
into lexical substitution operations in which, for example,
the symbol string representing a coherence group generator
with one electric field insertion
is sequentially spliced (minus the electric field)
into the observable symbol string at a location
just in front of a symbol representing a forward-directed
link on which the electric field can act
(or just after a corresponding backwards-directed link).

The exact lexical procedure is complicated by the choice,
discussed above, to use a distinct symbol to represent
every possible combination of electric fields and link
operators which may reside on a common lattice link.
The precise substitution rules may be encoded
in a static, dimension-independent splicing table,
discussed below in Sec.~\ref{sec:Splice}.
The commutation routines
carry out the resulting splicing and dicing operations,
keeping track of the associated minus signs,
suitably canceling any links and inverse links
arising at every joining operation, and yielding
a polynomial in observables (with terms collected)
produced by each commutator.

As discussed in section 4.5 of Ref.~\cite{CSVAII}, additional
terms involving triple products of Wilson loops can arise
in the commutator of a double-$E$ inserted loop with
a single-$E$ loop, due to the implicit normal ordering
of these observables.
Similarly, products of Wilson loops can arise in the
commutator of fermion bilinears with electric-field inserted
bilinears, due to their normal ordering prescriptions.
The commutation routines also generate and
add these extra normal-ordering induced terms
to the resulting observable polynomials.
See Sec.~\ref{sec:Commute} for further details.

\subsection {Observable generation \& classification}

As discussed in Refs.~\cite{CSVAI,CSVAII},
the most effective observable truncation scheme identified
to date is based on a grading of observables defined by
their \emph{strong-coupling} order.
This is the order in
a lattice strong coupling expansion at which an error first appears
if the given observable's expectation value is excluded from
an otherwise correct strong coupling expansion of the
ground state energy (or equilibrium free energy).
As explained in section 5.1 of Ref.~\cite{CSVAII}, the
strong-coupling order is the sum of an observable's
\emph{creation order} plus its \emph{expectation order}.
The creation order is the minimal sum of
the number of commutations with single-link fermion hopping 
terms plus twice the number of commutations with
single plaquette generators needed to produce the given
observable, while
the expectation order is 
the minimal sum of the number of commutations with single-link
fermion hopping terms plus twice the number of commutations with
single plaquette generators needed to yield the identity operator,
or a single-site fermion bilinear, starting from the given observable.%
\footnote
    {%
    These descriptions are not quite precise, particularly if
    there are compactified dimensions; see Ref.~\cite{CSVAII} for details.
    The counting of plaquette operations with double the weight of
    fermion hopping operations reflects the structure of the
    standard Kogut-Susskind Hamiltonian \cite{Kogut:1974ag}.
    In $\hat H/\lambda$, the gauge kinetic energy is $O(1)$
    while the fermion kinetic energy (or hopping term) is $O(1/\lambda)$ and
    the gauge potential energy (or plaquette term) is $O(1/\lambda^2)$.
    
    Note that the definitions of observable creation, expectation,
    and strong-coupling orders used in this work are all twice those
    used in Ref.~\cite{CSVAII}, so as to avoid half-integral creation
    or expectation orders for fermion bilinears.
    }

The observable generation procedure uses
the approach described in section 5.2 of Ref.~\cite{CSVAII}.
The target strong-coupling order is progressively incremented
up to the desired final value.
For each strong-coupling order,
the possible creation order is progressively incremented
from half the strong-coupling order, up to the strong-coupling
order minus four, for gauge observables, or minus two, for fermion bilinears.
(Four less, for gauge observables, because the expectation order of
any multi-plaquette loop is at least four.
Two less, for fermion bilinears, because the expectation order of
any multi-link fermion bilinears is at least two.)
New observables are generated via the commutation of known
lower order observables with ``primary'' coherence group generators.
These are generators, such as a plaquette with a single electric field insertion,
which cannot be produced via commutation of lower order generators.

Pure gauge Wilson loops (termed \tt{Loop} observables)
with specified creation and strong-coupling orders
are first generated followed,
for Hamiltonian theories,
by the generation of loops with one electric field insertion
(termed \tt{Eloop}'s),
and then loops with two electric field insertions
(termed \tt{EEloop}'s).
Then during the fermion minimization stage, a similar procedure is used
to generate fermion bilinears without any electric field insertions
(termed \tt{Fermion} observables) followed,
for Hamiltonian theories, by the generation of
fermion bilinears with one
electric field insertion (termed \tt{Efermions}).

Proper sequencing of these various observable generation
passes is needed to enable correct 
expectation order classification of observables.
As discussed below, expectation order classification
involves factorizing loops at self-intersections
and identifying the sub-loops as previously constructed
observables.  For this procedure to work, prior to
generating observables of any specific creation and
strong-coupling order, it is essential
that all observables of lower creation order,
or lower strong-coupling order, are already known.%
\footnote
    {%
    This is why it is not possible to merge the
    observable generation and geodesic equation evaluation
    steps and simply store any observable newly appearing
    in the evaluation of a given geodesic equation.
    Moreover, it is not sufficient to merely
    generate observables with progressively increasing
    creation orders, irrespective of
    their strong-coupling orders, as this simplified
    procedure fails to
    handle some rare cases where a loop of a given
    creation order factorizes into a sub-loop of the
    same creation order but lower strong-coupling order.
    (As an example, the creation order 14, expectation order 10
    loop \tt{xxyXyxYxYXXyxxyXYXyxYXXY} requires for its
    correct classification prior
    knowledge of the creation order 14, expectation order 6
    loop \tt{xyxxyXYXyxYXXY}.)
    }

As discussed in Appendix \ref{sec:Symmetries},
time reversal symmetry enables one to partition observables into
those whose expectation values act as classical coordinates on the
$N\,{=}\,\infty$ phase space,
and those whose expectation values serve as conjugate momenta and
vanish on the time-reversal invariant classical configuration space.
Expectation values of Wilson loops, and of loops with two electric
field insertions, act as classical coordinates, while
expectation values of loops with one electric field insertion
serve as conjugate momenta.
To generate expressions for the Hamiltonian gradient,
curvature, and the geodesic equations used for minimization
of the Hamiltonian, there is no need for any single-$E$ loops,
all of whose expectation values will vanish, to be present
in the master list of canonicalized observables.
However, the classification procedure described below for double-$E$
loops relies on the prior generation (and classification and storage)
of lower-order single-$E$ loops.
Consequently, even though they will not be needed in the
minimization, the observable generation routines
produce, classify and store single-$E$ loops.
For any given maximal strong-coupling order, the total
number of such single-$E$ loops is only a small fraction of the
number of double-$E$ loops, so the cost of retaining these
ultimately unnecessary observables is not significant.
For fermion bilinears, changing (in a Hamiltonian theory)
between the staggered and unstaggered version of the bilinear flips
whether its expectation value acts
as a classical coordinate or momenta
in the large-$N$ dynamics.
The observable generation routines only construct and store
in the master observable list those fermion bilinears which
serve as classical coordinates.

While generating observables with a known creation order
is quite straightforward,
determining the expectation order of each resulting observable
is more difficult.
Correct evaluation of the expectation order is
needed to determine how an observable's strong-coupling order
compares with the specified maximal
strong-coupling order for retained observables.
A brute force approach, directly applying the above definition,
is very inefficient.
Several different strategies are used to minimize the computational effort
required for expectation order determination.
Depending on an observable's type, two different quick
approaches are first employed to produce a lower bound on the observable's
expectation order.
These bounds are sufficient to eliminate from further consideration
a large fraction of observables newly produced by the commutators
of the observable generation process.

For observables containing electric field insertions, simply removing
those electric field insertions can only decrease the expectation order.
If the resulting no-$E$ observable,
either a Wilson loop or a fermion bilinear, is present in the
set of known (and correctly classified) observables, then
the expectation order of the no-$E$ observable provides a lower bound
on the expectation order of the original observable with 
electric field insertion(s).
Alternatively, if the no-$E$ modification of the observable is not
found in the set of already known observables, this indicates that
the strong-coupling order of the original observable must be larger
than the intended limit.%
\footnote
    {%
    This inference relies on the fact that the creation order
    of the no-$E$ modification of an observable will be no larger
    than the creation order of the observable with electric field
    insertions.
    }
This determination is performed by the \tt{Obs::noEbound} routine.
Of course, this procedure requires that the generation, and classification,
of all no-$E$ observables of a given strong-coupling order 
be performed prior to generating and classifying observables
of the same order containing electric field insertions.

A second lower bound on the expectation order is computed for gauge observables
(with or without electric field insertions), based on the observation that
the expectation order of a single trace observable in an SU($N$) gauge theory may be
larger than the expectation order of the same observable in a $U(1)$
gauge theory, due to the additional cancellations which may occur on multiply-traversed
links in a $U(1)$ theory, but cannot be smaller.
So a $U(1)$ theory expectation order determination provides a lower bound on
the SU($N$) expectation order.
Such a $U(1)$ theory determination (above two dimensions) is still very
computationally costly.
But a lower bound on the $U(1)$ expectation order may be relatively quickly
computed by projecting the observable's path through the lattice
onto a two-dimensional lattice plane, calculating the net $U(1)$ flux on
each plaquette enclosed by this two dimensional projection, summing
the absolute values of such flux over each enclosed plaquette of the projection
and then, in three or more dimensions, summing over all different orthogonal
lattice planes.
This calculation is performed by the \tt{Obs::u1bound} routine.

One or both of these lower-bound determinations are applied to all
newly generated observables other than fermion bilinears (without $E$'s),
for which no quick lower bound determination is available.
If neither lower bound can rule out the possibility that a newly
generated observable lies within the desired truncation set,
then a proper determination of its expectation order is needed
(or at least a determination that its order is too high for
the observable to lie within the truncation set).

For observables without electric field insertions, this is done
by finding all self-intersections of the observable,
factoring the observable into sub-observables at every such
self-inter\-section, and finding the minimal sum of expectation
orders for all such factorizations.
This is carried out by the \tt{Obs::factorize} routine.%
\footnote
    {%
    For fermion bilinears, it is necessary to consider both
    factorizations which extract from the self-intersecting bilinear
    a pure gauge loop, as well as factorizations into products
    of fermion bilinears with additional fermion/conjugate fermion
    pairs inserted at the intersection site.
    This is needed to account for the fact that, depending on the
    size of the sub-loop contour, the minimal expectation order
    could arise from either area-law or perimeter-law contributions.
    }
For no-$E$ observables without any self-intersections, or for
which no factorization leads to already known observables,
the expectation order must equal the observable's creation order.

For observables with electric field insertions, the only valid
procedure known is to actually perform commutations which will reduce
the observable to some already known and classified observable.
As discussed in Ref.~\cite{CSVAII}, one may argue that a
minimal commutation sequence determining the
expectation order may begin with commutation by a primary
generator (a plaquette, Polyakov loop if there are
compactified dimensions, or fermion hopping term)
which eliminates an electric field insertion.
This procedure is carried out by the \tt{Obs::reduce} routine.%
\footnote
    {%
    One subtlety, not explicitly addressed in Ref.~\cite{CSVAII},
    involves the fact that commutation of a double-$E$ loop with a
    single-$E$ generator leads to both modified double-$E$ observables
    as well as to products of no-$E$ loops.
    The \tt{Obs::reduce} routine also examines
    these loop-splitting no-$E$ terms as part of its
    expectation order determination.
    }

Appendix \ref{sec:counts} contains tables listing the
resulting number of canonicalized observables of a given type and
classification order, on cubic lattices of dimension 2, 3 or 4.%
\footnote
    {%
    For \tt{Loop} observables, the counts shown in
    Appendix \ref{sec:counts} agree with those shown
    in tables 1--3 of Ref.~\cite{CSVAII} for strong-coupling
    orders below 24 (labeled as order 12 in Ref.~\cite{CSVAII},
    but there are discrepancies at order 24,
    the highest order reported in Ref.~\cite{CSVAII},
    with undercounts in the older reference.
    The reason for this discrepancy with old results
    is not known.
    For \tt{Eloop} and \tt{EEloop} observables, there
    are more extensive discrepancies with the old results.
    This, it is believed, reflects the subtlety mentioned
    in the previous footnote: no electric field terms
    can arise from the extra loop-splitting terms in
    [\tt{EEloop}, \tt{Eloop}] commutators
    but, so far as can be recalled, such terms
    were not considered in previous observable
    classification routines.
    \label{fn:counts}
    }

\subsection {Observable approximation}

The current code was designed to allow testing the hypothesis
that the accuracy of the state representation, based
on retaining a large but finite set of observable
expectation values,
can be improved by approximating expectation values
of non-retained observables by products of expectation
values of retained observables, instead of simply
discarding expectation values of all non-retained observables.
The basic idea is to approximate a non-retained
observable (whose strong-coupling order exceeds the
specified limit) by a product of lower order
observables which arise from factorizing the 
original observable at a self-intersection.
This is motivated by the observation that
nearly all lattice loops are self-intersecting,
and that a factorization-based approximation is
consistent when expectation values
are well-described by a combination of area-law
and perimeter-law terms, which should be the case
in both strong and weak coupling regimes.%
\footnote
    {%
    This simple factorization-based approximation scheme
    is certainly not exact, and does not account for the
    interaction between electric flux sheets ---
    something which is especially significant in two (Euclidean)
    dimensions where typical self-intersecting loops lead to highly
    ``folded'' flux sheets.
    But the accuracy of the approximation should, it is hoped,
    improve with increasing lattice dimension.
    }

Given a non-retained observable (generated by the commutation
routines and identified as such by the classification routines),
the \tt{Obs::approximate} routine attempts to generate
a factorization-based approximation.
The routine examines all self-intersections of an observable,
assigning each possible split a numeric score which
is used to select preferred possible factorizations when
there are multiple possibilities.
The particulars of the intersection scoring are chosen to favor
factorizations which split the observable into comparably-sized
pieces, as this should maximize the probability
that both sub-observables are already known and stored
in the master observable list.
The scoring also disfavors splits which are immediately
next to an electric field insertion, and weakly disfavors
splits which arise from multiple traversals of the same link.
The scoring algorithm is unquestionably ad-hoc,
and is a target for future study and improvement.

\subsection {Generator selection \& normalization}\label{sec:gennorm}

The coherence group $\mathcal G$ is generated by the Lie algebra consisting
of all anti-Hermitian \tt{Loop}, \tt{Eloop} and \tt{Fermion}  operators.
Elements of this coherence group, when acting on a suitable base state $|0\rangle$,
generate large $N$ gauge theory coherent states.
In a lattice gauge theory, the base state may be taken to be
the vacuum state (of the Kogut-Susskind lattice Hamiltonian)
in the limit of infinitely strong gauge coupling
and large fermion mass.%
\footnote
    {%
    The state $|0\rangle$ is annihilated by
    all electric field operators $\{ E_\ell \}$,
    all fermion field $\{ \psi_s \}$ on the even sublattice,
    and all conjugate fermion field $\{ \psi^\dagger_s \}$
    on the odd sublattice.
    }

The coherence group $\mathcal G$ has a non-trivial little group $\mathcal L$
which leaves this base state invariant.
In the gauge sector, the Lie algebra of this little group consists
of all \tt{Eloop} generators whose reduction to a \tt{Loop} vanishes.
The generator reduction is the result of the ``inner commutator'' where
the electric field insertion in the \tt{Eloop} generator is allowed to act
on the links of within the \tt{Eloop} generator itself.%
\footnote
    {%
    More formally,
    if $\mathrm{tr} (E_\ell \, U(C))$ denotes an operator
    in which a Wilson loop $U(C)$ has an electric field insertion
    on some link $\ell$,
    the generator reduction refers to
    the commutator $ [(E_\ell)_{ij} , \, U(C)_{ji} ]$.
    Examples of gauge generators with vanishing reductions, and hence
    elements of the little group Lie algebra, are
    \tt{AyXYx} and \tt{AxyXY - ByXYx}.
    }
In the fermion sector, all even length bilinears
are elements of the base state little group.
Consequently, representatives of the coset space
$\mathcal H \equiv \mathcal G / \mathcal L$
consist of exponentials of anti-Hermitian linear combionations
of Wilson loops, odd length fermion bilinears, plus
of so-called ``commutator generators'',
produced by commuting a \tt{Loop} generator with $\tr(E^2)$.%
\footnote
    {%
    This form of coset representatives holds at least within
    the domain of convergence of a Baker-Campbell-Hausdorff
    expansion of the product of elements of $\mathcal G$
    and $\mathcal L$.
    }

As one moves around in the large-$N$ classical phase
space (i.e., moves to a different large-$N$ coherent state)
this little group is conjugated.
The coset space $\mathcal H$ is not itself a group
because the little group $\mathcal L$ is not a normal subgroup.
However, coherence group elements $g \in \mathcal G$
may be factorized as
$
    g = h \times l
$
with $l \in \mathcal L$ and $h \in \mathcal H$,
and hence only elements of the coset space need be considered
when generating coherent states.
Consequently, when deforming a coherent state,
it should be sufficient to consider only those
deformations produced by 
\tt{Loop} generators,
\tt{Eloop} commutator generators.
and odd-length fermion bilinears.

As noted initially, truncating the infinite set of
coherence algebra generators to a finite computational
subset is an unavoidable step in this variational approach.
Every generator may be assigned a \textit{creation order} which
is, essentially, the minimal sum of the number of
commutations with single-link fermion hopping 
terms plus twice the number of commutations with
single plaquette generators needed to produce the given
generator.
This classification reflects the relative importance of a
given generator on the strong coupling dynamics of the theory.

However, it is clearly desirable to examine
\textit{ex post facto} assessments of a generator's importance
based on the relative size of its component in the gradient $dH$
(\ref{eq:gradcurv}) or the predicted minimum location
$\delta$ (\ref{eq:delta}).
But no such comparison of relative component magnitudes
is meaningful in the absence of some basis-independent
normalization of generators.
The lack of any natural metric on
the coherence Lie algebra makes this problematic.

A partial solution, for Hamiltonian theories, is to
regard the base state expectation of
the double variation of the gauge kinetic energy,
$\langle 0 | [[\tr(E^2), \cdot \, ], \cdot \,  ]|0\rangle$,
as defining a pseudo-metric.
This pseudo-metric is non-degenerate within the subspace
of generators in the coset space $\mathcal H$.
%
But there is a null-space involving coherence group generators
of the base state little group $\mathcal L$.
Such generators will have vanishing norm as defined by the
above pseudo-metric.

If all \tt{Eloop} generators retained in the generator truncation
have non-vanishing single-loop terms in their reduction
and different generators involve geometrically distinct loops,
then one may uniquely and physically normalize
generators by requiring the strong coupling expectation of their
double variation of $\tr(E^2)$ to give unity.
Generators involving distinct loops or odd-length fermion bilinears are
automatically orthogonal by this measure.
Such normalization is performed if the global boolean variable
\tt{gennorm} is set to true.

If one chooses to include generators with vanishing
(or purely multi-trace) reductions, or even length fermion bilinears,
then such generators will have
vanishing norm as defined by the above pseudo-metric.
Implementing a well-defined normalization for such generators
is unavoidably ad-hoc.
The current code simple sets to unity the sum of squares of
the coefficients of all terms in such generators.

For Euclidean theories, instead of the double variation of $\tr(E^2)$
one may regard the double variation of the gauge field entropy $S$,
evaluated in the strong coupling limit, as defining a pseudo-metric
for gauge generators (with the same caveats regarding its null space
as above).
But the fermionic contribution to the entropy cannot serve the same
purpose for fermion generators, as only 
fermion bilinears which
begin and end on the same site lead to non-vanishing entropy
variations in the strong coupling limit.
Instead, in Euclidean theories the double variation of the fermion mass term
may be used to define a metric on the space of fermion bilinear generators.
This is the procedure used when generator normalization is enabled.

\subsection {Polynomial representation}

In Euclidean theories, the geodesic equation for any
Wilson loop is a linear combination of loops,
while the variation of the entropy (needed to minimize
the free energy) is a quadratic polynomial in loops.
In Hamiltonian theories, terms cubic in Wilson loops
arise via normal-ordering in the commutator of a
generator with one electric field insertion and
a gauge observable with two electric field insertions.
Approximating the expectations of observables outside
the selected truncation set by a product of expectations
of observables in the retained set further can further
increase the order of monomials in observables which
need to be stored.

As discussed above, every retained observable is
assigned a unique index in the master list of
canonicalized observables, so an order-$k$ monomial in
observables is represented by a floating point coefficient
together with a $k$-tuple of observable indices.
Large calculations can involve many millions of polynomial geodesic
equations comprising altogether many billions of individual monomial terms.
Consequently, storage efficiency is of some concern,
needing to be balanced against programming convenience.

Several different representations for storing observable
polynomials are used.
During commutator evaluation
a fixed-size \tt{PolyTerm} class object
is used to represent monomials.
Each \tt{PolyTerm} contains a
\tt{std:array<unsigned int,PSIZ>} data member holding observable
list indices plus a
floating point coefficient,
with \tt{PSIZ} a compile-time constant which is
set to 2 for Euclidean theories and 4 for Hamiltonian theories.
As the trivial identity loop with a fixed expectation value of 1
is always assigned index 0, a \tt{PolyTerm} instance can
represent an observable monomial of any order up to \tt{PSIZ},
while keeping the code needed to manipulate 
observable polynomials as simple as possible.
The multiple monomials comprising a complete observable polynomial are 
stored either in a simple list (an \tt{ObsPoly} class object)
or, to efficiently collect common terms,
in a hash table (a \tt{PolyMap} class object).

Because most terms in the geodesic equations are linear monomials,
storing a \tt{PSIZ} length index list for every monomial is not
the most efficient use of memory.
To enable large scale calculations, the \tt{PolyMap} objects
generated by the commutation evaluation routines are converted
to a packed format (a \tt{Poly}) in which a single header,
recording the length of the data structure, is followed by
a string of monomial coefficients and indices, with negated indices
used to signal non-terminal non-vanishing indices in the set of
observable indices for a given monomial, eliminating the need to
store redundant vanishing indices.
This packed format is efficient for numerical evaluation
of the symbolic polynomial representation, as well as for
saving to and retrieval from external files.

\subsection {Numerics}

Programming the numerical part of the coherent state variational
algorithm is relatively straightforward compared to the symbolic stages.
A list of numerical expectation values of all the retained
observables is initialized with known starting values
For the $O(N^2)$ gauge stage of the minimization,
this is the infinite coupling limit in which
all gauge observables other than the identity loop have
vanishing expectation values.
For the $O(N)$ fermion minimization stage,
this is the large mass limit in which
the only non-vanishing fermion bilinears are those
in which the fermion and conjugate fermion reside on
the same site, with the bilinear expectation value
determined by the corresponding Wilson loop expectation.

As explained in appendix \ref{sec:Symmetries},
in Hamiltonian theories
all non-vanishing observable expectation values are real.
But, as detailed in appendix \ref{sec:Euclidean}, in Euclidean formulation
theories with fermions, the chosen form of fermion discretization makes
translation invariance manifest with a kinetic term of the same
form as in Hamiltonian theories, but imparts imaginary expectation
values to even-length fermion bilinears.
This factor of $i$ is not made explicit anywhere in the code,
except in the final (Mathematica format) numerical results files,
because such even-length bilinear expectation values are always
multiplied by a pure imaginary coefficient.
In the list of stored expectation values,
and in
the polynomials representing the free energy, its gradient
or curvature, or the geodesic equations for Euclidean theories,
every even-length bilinear means $-i$ times the actual bilinear
expectation value, and
real arithmetic suffices for evaluating all such polynomials.

After setting (or changing) the finite value of the
lattice coupling constant $\lambda$
or fermion mass $m$,
moving to the resulting non-trivial ground (or equilibrium)
state involves iterative Newton minimization.
In each such minimization step,
one first needs to evaluate the gradient ($dH$) and curvature ($ddH$) of the
Hamiltonian (or free energy).
These are a vector and matrix, respectively, of dimension
equal to the number of defined (vacuum representation) generators,%
\footnote
    {%
    More precisely, the number of time-reversal even 
    vacuum representation generators, as these are the
    operators whose exponentials generate motion within
    the large-$N$ classical configuration space.
    Gradient components corresponding to time-reversal
    odd generators vanish in all time-reversal invariant
    states, and the curvature matrix is block diagonal
    with separate time-reversal even and odd blocks.
    }
with components
\begin{align}
    (dH)_\alpha \equiv [H, e_\alpha] \,,\qquad
    (ddH)_{\alpha\beta} \equiv [ [ H, e_\alpha], e_\beta] \,,
\label{eq:gradcurv}
\end{align}
where $e_\alpha$ denotes 
a vacuum representation generator,
and each commutator is a previously evaluated polynomial in the
expectation values of observables.
The curvature matrix $ddH$ is real but not, in general,
precisely symmetric.
The antisymmetric part equals
a gradient component in the direction of
a commutator of generators,
$
    (ddH)_{[\alpha, \beta]}
    =
    [H, [e_\alpha, e_\beta]]
$.
Due to the truncation of generators to a finite set,
some antisymmetric portion of the curvature will remain even
at the conclusion of the minimization.
%

After evaluating the gradient and curvature,
one finds the predicted location $\delta$ of the
minimum by solving the linear equation
\begin{align}
    ddH \cdot \delta = -dH .
\label{eq:delta}
\end{align}
One then moves to the predicted minimum by
integrating the geodesic equations describing the change
in each (retained) observable's expectation value.
These geodesic equations comprise a large set of coupled,
polynomial ordinary differential equations (ODEs) of the form
\begin{align}
    \frac{d \mathcal O_i}{ds}
    =
    \sum_\alpha \> [\mathcal O_i, e_\alpha] \,\delta_\alpha \,,
\end{align}
where $i$ runs over all observables in the retained set with 
$\mathcal O_i$ denoting the corresponding observable's expectation value,
and each commutator is again a known polynomial in all the
expectation values $\{ \mathcal O_i \}$,
which are themselves functions of the affine parameter $s$
which is integrated from 0 to 1.
%
This iterative procedure is repeated until convergence,
at which point the excitation spectrum (in Hamiltonian
formulations) may be evaluated by solving the
small-oscillation generalized eigensystem
\begin{align}
    ddH \cdot \delta = -i\omega \, L \cdot \delta \,,
\label{eq:spec}
\end{align}
where $ddH$ is the Hamiltonian curvature
(in any chosen irreducible representation)
and $L$ is the corresponding
\href{https://en.wikipedia.org/wiki/Lagrange_bracket}
{Lagrange bracket matrix},
with components
\begin{align}
    L_{\alpha\beta} \equiv [e_\alpha, e_\beta] \,.
\end{align}

The curvature matrix is block-diagonal with one block,
$ddH_{++}$, involving pairs of time-reversal even generators
and one block, $ddH_{--}$, involving pairs of time-reversal
odd generators,
while the Lagrange bracket matrix is block off-diagonal
(and antisymmetric)
with one block, $L_{+-}$, invoving the
commutator of time reversal even and time reversal odd generators.
The small oscillation eigensystem in the first-order form (\ref{eq:spec})
may be reexpressed as a smaller second-order eigensystem
\begin{align}
    ddH_{++} \cdot \delta_+ = -\omega^2 \,
    L_{+-} \cdot (ddH_{--})^{-1} \cdot L_{-+} \cdot \delta_+ \,,
\end{align}
which is computationally more efficient.

This summary shows that three types of numerical support routines
are needed:
\begin{enumerate}\advance\itemsep -6pt
\item
    An efficient and accurate linear system solver for
    dense real non-symmetric matrices;
\item
    An accurate ODE solver for large systems of first order 
    different equations;
\item
    An accurate generalized eigensystem solver for dense
    non-symmetric matrices.
\end{enumerate}
The \href{https://arma.sourceforge.net}{Armadillo C++ library}
for linear algebra and scientific computing provides high
quality and fast routines for the linear algebra tasks 1 and 3.%
\footnote
    {%
    It is also possible to compile the code using the 
    \href{https://libeigen.gitlab.io/}{Eigen C++ library}
    for linear algebra.
    This capability was added
    for testing purposes, including access to 80-bit arithmetic
    on Intel processors, but use of the Armadillo library is
    generally preferred.
    }
For the ODE integrator, this code employs, by default, an adaptive
fourth-order embedded Runge-Kutta algorithm \cite{SS32}
which uses the same set of function evaluations to compute
two different approximations to the solution,
with a norm of the difference between them used to adjust
the integration step size.
Various alternative embedded Runge-Kutta schemes, of both
lower and higher order, are also available and may be easily selected.
For further details and references, see Sec.~\ref{sec:RK}.

\subsection {Multi-threading}

On multi-core systems it is, of course, desirable to take
advantage of the speed-ups which multi-threading can enable.
To do so, it must be possible to exploit parallelism without
incurring excessive cost from any needed synchronization or
locking operations.

The symbolic construction of geodesic
equations for all retained observables
is a major time-consuming step which can be very
easily and effectively parallelized.
The loop over all observables is 
replaced by a parallel \tt{for_each} loop over a much shorter
list of buckets of observables, with the evaluation of 
geodesic equations for each such bucket of observables automatically
distributed across multiple threads.%
\footnote
    {%
    Linking with the Intel$^{\text{\textregistered}}$ Threading Building
    Blocks template library \cite{IntelTBB} enables very easy
    multi-threading via parallel \tt{for\_each} loops.
    }
%
The only special considerations needed involve
the declaration of various statistics counters
as well as a few dynamically-sized buffers
used by the commutation and observable assessment routines for
construction of symbol strings, temporary sets of
observables, etc.
With single-threaded execution such buffers could be
declared \tt{static}, thereby minimizing the
needed free store allocation and deallocation operations.
With multi-threaded execution, such buffers must not
be declared \tt{static} (since static variables are
unique and shared across all threads).
Instead, there is a choice between automatic scope,
local to a given block, or \tt{thread_local} storage,
which persists like static variables but is local to each thread.
Explicit testing found that automatic block-scoped declaration
of buffers is preferable in the commutation routines,
while \tt{thread_local} storage turns out to be preferable
for the buffers in the observable analysis routines.
The various statistics counters which are updated 
during polynomial script generation 
are declared as \tt{alignas(128) std::atomic<ulong>}
so as to be thread-safe, with the additional alignment
specification serving to prevent cache line contention.

The time needed for construction of the symbolic expressions for the
gradient and curvature of the Hamiltonian (or free energy)
is generally a small fraction of the time needed for observable
generation or construction of the geodesic equations.
Consequently, the current code does not
bother with parallelization of these steps.

During the numerical minimization stage of a calculation, 
the evaluation of the derivative vector of all observables
(i.e., the numerical evaluation of the polynomial scripts
encoding the geodesic equations for each observable)
can also easily benefit
from parallelization, simply by using another parallel
\tt{for_each} loop over buckets of observables to distribute
the computation across multiple threads.

Effectively parallelizing
the observable generation phase is more challenging.
One must prevent race conditions or inconsistencies in access
to the growing observable list (or the canonicalization cache).
The simplest approach for doing so would be to require
acquisition of a \tt{mutex} lock on
the observable list every time a new observable is added.
Explicit testing found that the cost of such
\tt{mutex} locking on every insertion 
is unacceptably high, yielding performance worse than
single threaded execution.

In the second attempt at multi-threaded observable
generation,
on every pass of the observable generation process
the list of possible starting observables was
divided into ``buckets'' of observables,
each of which will be assigned to one of the
threads spawned by a parallel \tt{for_each} loop.
Each of the observable generation routines called by
these parallel \tt{for_each} loops
first made a private copy of the current
observable list, to which new observables
found by the given thread were appended
during the evaluation of the commutators
involving the assigned bucket of observables.
Only following all commutator evaluation
for a given bucket of observables
was a \tt{mutex} lock acquired to synchronize merging of the
augmented local observable list back into the master list.
This reduced the cost of thread synchronization to an
acceptably small amount and made parallel execution
worthwhile.
However, when attempting to generate large observable sets,
on machines with many cores, replicating the
progressively growing observable list
(actually a hash table) in every thread can consume
problematically large amounts of memory (exceeding available physical
memory and essentially stalling forward progress).

To avoid such excessive memory consumption,
the next approach used a \tt{thread_local} 
\tt{std::unordered_set} to hold observables newly generated
by a given thread, with \tt{mutex}-protected merging of these
``inboxes'' of observables at the end of each bucket-processing
observable generation routine, followed by (single threaded)
addition to the master observable list of all new observables
at the end of each observable generation pass.
This approach works, with reasonable memory consumption,
but has a downside that the delayed addition of new observables
to the master list also delays additions to the canonicalization
cache of short observables,
leading to much worse canonicalization cache performance
during the observable generation phase,
with excessively many repeated canonicalizations involving
all possible symmetry transformations of short observables.

The current approach remedies this last problem by using
single-threaded observable generation, with immediate
addition of new observables to the master list as soon
as they are identified, as long as the total number of
observables is below a threshold (set in \tt{Build.cpp}),
while switching to multi-threaded observable generation,
as just described, above this threshold.
This approach has allowed the largest observable sets listed
in appendix \ref{sec:counts} to be successfully generated
on desktop computers.%
\footnote
    {%
    The generation of the multi-hundred million observable sets
    shown in appendix \ref{sec:counts} used an
    Apple Mac Studio with 128Gb of memory.
    Generation of still larger observable sets,
    numbering in the billions, is limited
    both by the needed memory capacity and by
    the current use of unsigned 32-bit integers
    to store observable indices.
    It would, of course, be straightforward to
    switch to 64-bit integers for such indices.
    But this would substantially increase the size
    of polynomial scripts encoding geodesic equations 
    which, as discussed below,
    can already be problematically large.
    Hence this step has not been taken (yet).
    }

Using the Intel$^{\text{\textregistered}}$ TBB library
\cite{IntelTBB}, parallel \tt{for_each} loops will, by default,
spawn as many threads as there are available processor cores.
This seems to work fine for geodesic script generation
and evaluation, but
when running on machines with many cores,
observable generation
for some intermediate size observable sets
can experience significant contention between
threads due to the synchronization required
for inbox merging of different threads.
A global variable (\tt{global.maxthread})
may be set to a non-zero value to limit
the maximum number of concurrent threads
used in any multi-threaded portion
of the program, and
suitably adjusting the value of this limit,
depending on the number of processor cores and
maximum memory available, can significantly 
improve observable generation performance for certain
observable generation cases.

\subsection{Input and output files}

The program handles several types of input and output files:
\begin{description}
\item [Command input]~\\
    In addition to reading from the standard input stream \tt{std::cin},
    an ordinary text file of pre-prepared commands may be 
    input with the \tt{read} command.

\item [Program output]~\\
    Program output may be directed to an output file,
    instead appearing on the standard output stream \tt{std::cout},
    by using the \tt{write} command.
    (One may also redirect output back to \tt{std::cout}.)

\item [Sys info files]~\\
    All symbolic information, including the selected
    set of observables and generators, the observable
    polynomials representing the
    gradient and curvature of the Hamiltonian (or free energy),
    and the symbolic geodesic equations may be saved
    in a \textit{system information} (or \tt{.sys}) file.
    These binary sus-info files may also be
    loaded using the \tt{load} command (or by listing
    on the initial program command line).
    This allows one to perform the symbolic part of a calculation
    just once and subsequently reload and use that information
    multiple times during the numerical part of a calculation.
    System information files may be saved with the explicit
    \tt{save sys} command or, if \tt{autosave} is set to \tt{true},
    will be written automatically when geodesic equations are generated.

\item [Vev data files]~\\
    The list of observable expectation values (or \textit{vev's}),
    together with the values of adjustable coupling constants,
    may be stored in \textit{vacuum expectation} (or \tt{.vev}) data files.
    Such files may hold multiple data sets recording observable
    expectation values for different values of coupling constants.
    Information stored in such binary vev-data files may
    also be loaded
    using the \tt{load} command
    (or by listing on the program command line).
    The particular data set loaded may be user specified or, by default,
    is the last data set in a vev-data file.
    Vev data files may be saved with the explicit \tt{save vev}
    command or, if \tt{autosave} is set to \tt{true},
    will be written automatically after each successful minimization step.

\item [Results files]~\\
    When a sequence of minimizations are performed
    (by the \tt{do} \textit{coupling-name} \tt{...} command),
    human readable results are automatically written to a
    result file in a format which may be directly input
    to \textsl{Mathematica}.
    Such \textit{MMA} (or \tt{.m}) files record the values
    of the Hamiltonian or free energy, the expectation values
    of a selected set of observables (by default, those
    appearing in the Hamiltonian or free energy) and,
    for Hamiltonian theories, the values of particle masses,
    or small oscillation frequencies, in all symmetry channels
    with defined curvatures.

\end{description}

For theories with fermions,
separate sys-info, vev-data and MMA results files are generated
for the pure gauge parts of a calculation and for the fermion
portions of a calculation.
Gauge information is stored in files named to indicate the
theory type and the maximum orders of generators and observables,
with file extensions indicating the file type.
Sys-info and vev-data file names include an additional two
(upper case) characters formed from a hash code or check-sum
of the entire set of gauge observables;
this serves to distinguish which vev-data files are associated
with a given sys-info file.
If observable approximation is turned on, a further `a'
is appended to the file name.
As an example, gauge stage output files for 3+1 dimensional Hamiltonian
Yang-Mills (or QCD)
with order 4 generators and strong-coupling order 20 observables
(without observable approximation)
could be named
\tt{YM3h_4g20QZ.sys}, 
\tt{YM3h_4g20QZ.vev},  and
\tt{YM3h_4g20.m}.

Fermion information is stored in files named to indicate the 
theory type as well as the generator and observables orders for both
gauge and fermion stages, with separate hash codes for the gauge
observable set (which will match the associated \tt{YM*} data files)
and the fermion observable set.
As an example, fermion stage output files for 3+1 dimensional Hamiltonian
QCD with order 4 gauge generators and order 20 gauge observables
along with order 3 fermion generators and order 10 observables
(without observable approximation)
could be named
\tt{QCD3h_4g20QZ_3f10KW.sys}, 
\tt{QCD3h_4g20QZ_3f10KW.vev},  and
\tt{QCD3h_4g20_3f10.m}.

Sys-info and vev-data files are stored together in theory-specific subdirectories
(\tt{YM3h}, \tt{QCD3h}, etc.) of the main save directory (by default, \tt{./save}).
Results files are stored in theory-specific subdirectories of the
MMA results directory (by default, \tt{./MMA}).
These directory specifications may be changed with the
\tt{set savedir} and \tt{set MMAdir} commands.

\subsection {Disk swapping}

As discussed below,
with large (multi-million) observable sets and quite a few
generators, the size of the polynomial scripts encoding
geodesic equations 
can be hundreds of gigabytes or more, substantially
exceeding physical memory size in typical workstations.
For such calculations, it may be necessary to swap
the geodesic equations for individual
buckets of observables to and from disk as needed.
With the boolean \tt{global.geoswap} set to true, 
during the script generation phase
each thread handling a particular bucket
of observables writes the generated geodesic equations
for its bucket of observables to disk upon completion
and then frees the associated memory.
During the numerical evaluation of the geodesic equations,
each thread handling a particular bucket of observables
first reads the needed portion of the geodesic equations
in from disk, performs the numerical evaluation,
and then frees the associated in-memory portion of the 
geodesic equations.
This process is repeated during every derivative
evaluation required by the ODE integrator, on every
Newton minimization step.

The writing to disk during script generation must use
a \tt{mutex} lock to ensure tread-safe operation, but
for large calculations
this imposes a negligible penalty as the time needed to
generate the geodesic equations for a substantial bucket of
observables easily exceeds the time required to write the newly 
generated equations to disk.
But when numerically evaluating the geodesic equations,
the time needed to evaluate each equation in a given
bucket is small compared to the typical time needed to
fetch that bucket of equations from disk.
Multi-threaded execution is still beneficial, but when
running on desktop systems with a solid-state disk capable of
transferring 5--6\,Gb/sec, three or four active threads can
saturate that bandwidth.

\subsection{Resource estimates}

The average length of Wilson loops having
a strong-coupling order $M$ is very nearly the same as $M$.%
\footnote
    {%
    For example, in two dimensions the average length of
    order 28 Wilson loops is 29.4.
    In four dimensions, the average length of order 20
    Wilson loops is 20.0.
    }
The number $N(M)$ of physically distinct (canonicalized) loops
of strong-coupling order $M$ may be roughly estimated
as
\begin{equation}
    N(M) \approx \, (3/2 + d/2)^{M-2} \, M^{-1} \, / n_{\rm sym}(d) \,,
\end{equation}
where $d$ is the cubic lattice dimension and
$
    n_{\rm sym}(d) \equiv 2^{d+1} \, d!
$
is the dimension of the lattice symmetry group
(lattice point group times charge conjugation).
For loops with one or two electric field insertions, the corresponding
estimates acquire additional factors of $M$ or $M^2/2$, respectively.

The above expression for $N(M)$ is by no means an exact asymptotic form,
but for dimensions 3 and 4 agrees tolerably well (to within $\approx 10\%$) with
the explicit counts for larger orders shown in appendix \ref{sec:counts}.
In two dimensions, the above expression substantially overestimates
the true counts by a factor of 10--15.
But the depicted exponential growth in the number of loops with increasing
strong-coupling order agrees well with the explicit results in
appendix \ref{sec:counts}.%
\footnote
    {%
    Given that the average length of order $M$ loops is very close to $M$,
    one might think that the number of loops of strong-coupling order $M$
    would be close to the number of closed, non-backtracking loops of length $M$.
    This is not true, as the number of non-backtracking closed loops 
    grows much faster with length, proportional to $(2d-1)^M$.
    }

Given that, these days, even low-cost desktop computers have multi-Gb memories,
one might naively expect that memory capacity would not be a key limiting
factor for calculations involving only, say, a few million observables.
However, one must also consider the number of generators in use,
the typical number of individual terms in each such generator, and how
many terms may be generated in the commutator of an observable with
a single generator term.
Each generator (beyond a few lowest order ones
which are invariant under a subgroup of lattice symmetries)
will have at least $n_{\rm sym}(d)$
terms after projection onto the
vacuum symmetry representation.
For so-called commutator generators, the total number of terms
will be $n_{\rm sym}(d)$ times the number of links in the underlying loop.
When commuting with an observable,
each term in an \tt{Eloop} generator can act on every link
in the observable which lies in the same direction as the
electric field insertion in the generator term.
So an \tt{Eloop} generator with $n_{\rm sym}(d)$ terms
acting on a length $M$ observable will typically lead to about
$
    n_{\rm sym}(d) \times L \times \times M \times 1/d
$
distinct terms in the result of the commutator,
where $L$ is the length of the generator in question.
So for $K$ generators of average length $L$,
the complete set of polynomials which encode
the geodesic equations for all \tt{Loop} observables
up through order $M$ can contain upwards of
\begin{equation}
    K \times n_{\rm sym}(d) \times L \times M/d \times N(M)
    \approx
    K \times L \times (3/2+d/2)^{M-2} /d
\end{equation}
individual monomials, each of which is stored as a
\tt{PolyTerm} instance.
As indicated above, for Euclidean theories
each such \tt{PolyTerm} occupies 16 bytes
(an 8-byte coefficient and two 4-byte observable indices),
while for Hamiltonian theories each \tt{PolyTerm}
requires 24 bytes (for a coefficient and four observable indices).

For a 2D Euclidean theory,
the above estimate suggests that a calculation using
order 28 \tt{Loop}'s and, say, a dozen generators
of average length 6 might require upwards of
12\,Tb of storage.
The actual size is not nearly so large, because
the above approximation for $N(M)$ is a substantial
overestimate in two dimensions, and because collecting
common terms in the result of commutator evaluation
leads to substantially shorter expressions.
But the actual size of the polynomial script
for the described 2D calculation is still 300\,Gb!
For a 3D Hamiltonian calculation, [[...]]

This illustrates why polynomial storage, memory utilization,
and possible swapping of information to disk are all
critical aspects of program design.

\subsection{Data structures}

Following standard C++ practice, a variety of custom
classes are defined and used to represent observables,
operators, symmetries, representations, polynomial terms,
and numerous other objects as detailed below in section \ref{sec:files}.
Liberal use is made of standard C++ container classes, especially
\tt{std::vector} and \tt{std::unordered_hash}.
Letting the standard container classes take care of all memory
allocation and deallocation is highly beneficial.

An initial program version used nested 
standard containers for the multi-dimensional data structures
storing the symbolic polynomial scripts defining the
gradient and curvature of the Hamiltonian (or free energy)
and the geodesic equations.
For example, the curvature was initially 
defined as a
\tt{std::array <std::vector <std::vector <std::vector <ObsPoly>>>>}
where the successive indices used to access a particular element
select, respectively, the symmetry representation
(whose possible range is known at compile-time),
a particular term in the Hamiltonian (or free energy), and a pair of
defined coherence group generators in the given representation.
This works fine for smaller calculations where there is ample
physical memory available.
However, as detailed above, for larger calculations the storage
needed to hold all polynomial script information, particularly for
the geodesic equations, can become very large.
To permit good virtual memory and processor cache performance,
it is important to keep the
geodesic equation information for a given bucket of observables
grouped together in memory.
And for calculations where explicit swapping of such information
to and from disk is needed, a serialized representation of data which can
be efficiently written to or read in from disk is essential.
Therefore, the current code stores the entire 
set of polynomials defining the geodesic equations for
a given bucket of observables as a single
vector of packed serialized data in which an eight-byte header
separates and identifies each distinct polynomial.
Similarly, a single such data record is used for the
entire set of polynomials comprising the gradient,
and further such data records are used for the
curvature and Lagrange bracket in each symmetry representation.
These packed representations are efficient to process,
enable good processor cache and memory performance,
and can be directly written to or read from disk with no
conversion needed.
(The various data headers used in these data records
are described below in section \ref{sec:Data} and related functions
summarized in sections \ref{sec:Poly} and \ref{sec:Save}.)

\section {Programming assumptions \& conventions}

Basic programming environment assumptions are:
\begin{itemize}\advance\itemsep -3pt
\item
    Both \tt{signed char} and \tt{unsigned char}
    types are available.
    Objects in the standard C++ container class
    \tt{std::string<char>} are used to
    represent observables.
\item
    \tt{short} integers are 16 bits,
    \tt{int} integers are 32 bits, and
    \tt{long} integers are 64 bits.
\item
    \tt{float} floating-point numbers are 32 bits, while
    \tt{double} floating-point numbers are 64 bits.
\item
    The version of C++ available is at least ISO C++ 2020.
\item
    The Intel$^{\text{\textregistered}}$ Threading Building Blocks template library
    \cite{IntelTBB}
    is available to support multi-threaded execution on multi-core systems.
\item
    The \href{https://arma.sourceforge.net}{Armadillo C++ library
    for linear algebra and scientific computing} is available.
\end{itemize}

As discussed below, by setting a preprocessor variable
\tt{NUM32} the code may be compiled to use 32 bit integers
for observable indices and 32 bit floats for polynomial coefficients.
When \tt{NUM32} is not set, the default is to use 64 bit
variables for both observable indices and polynomial coefficients.
Compiling (and running) in \tt{NUM32} mode nearly halves the size of
system information files and memory used for polynomial storage,
but in this mode the maximum number of retained obsrevables
is limited to $2^{31} - 1$, and 63 is the maximum allowable
generator or observable order creation order.

Stylistic conventions employed in the code include:
\begin{itemize}\advance\itemsep -3pt
\item
    Variable names are \tt{lower-case};
\item
    Class names and source file names are \tt{Capitalized};
\item
    Preprocessor symbols are \tt{FULLY CAPITALIZED};
\item
    Brace initialization of variables is generally preferred,
    except for simple integers or when using parentheses
    avoids the need for a nuisance cast of integer types;
\item
    Opening and closing braces of code blocks are lined up vertically
    and indented relative to the preceding statement or expression.%
    \footnote
	{
	This is not the most common C or C++ convention, but it is the best choice
	for program clarity.
	}
    The few exceptions are generally for short, simple inline functions in header files.
\end{itemize}

\section {Compilation \& linking}

A single \tt{Makefile} handles compilation and linking for
versions of the program targeting any theory class.
The choice of theory class may be specified by including an
argument of the form \tt{THY=}\textit{theory} on the \tt{make} command line,
where \textit{theory} is one of the available preprocessor shorthand theory names
listed in Sec.~\ref{sec:Theory}.
One may also define \tt{THY} via an exported shell environmental
variable prior to running \tt{make}.
A successful compilation will result in an executable named
\tt{gordion-}\textit{theory}, with \tt{gordion} a symlink pointing
to that executable.
Alternatively, if \tt{THY} is not specified on the \tt{make} command line
(or defined as a shell environmental variable), then the target name of an existing
\tt{gordion} symlink is used as the default theory selection.
An error message stating the necessity of defining \tt{THY} is generated
if \tt{THY} is undefined and no default can be inferred from an existing
\tt{gordion} symlink.

It is assumed that the \tt{C++} compiler supports a \tt{-M}
flag to automatically generate dependency files.
The command \tt{make clean} removes all \tt{.o} object files, \tt{.d} dependency
files, and the target executable.
If the setting of \tt{THY} does not match that of an existing
\tt{gordion} symlink, then a \tt{make clean} is automatically performed
along with replacement of the \tt{gordion} symlink.

If \tt{DBG} is set (for example, by including \tt{DBG=1} on the command line
or via an environmental variable), then a debugging compilation
with a \tt{-g} compiler flag is performed.
Otherwise, an optimized parallel compilation
with \tt{-O3 -DPARALLEL} compiler flags is performed.

If \tt{NUM32=1} is include on the \tt{make} command line, then the
preprocessor variable \tt{NUM32} will be set and, as indicated above,
32 bit variables will be used for observable indices and polynomial coefficients.
Note that sys-info and vev-data files produced by \tt{NUM32} compilations are
incompatible with those from compilations using the default (64-bit) types
for observable expectation values and polynomial coefficients.

The environmental variable
\tt{CPATH} may need to be suitably defined
so that Armadillo and Threading Building Blocks (TBB) include files
are found by the compiler.
Likewise,
\tt{LIBRARY_PATH} may need to be suitably set so that
\tt{-ltbb} and \tt{-larmadillo} linker flags suffice for
\tt{ld} to access the Armadillo and TBB object libraries.


\section {Source files}
\label{sec:files}

The following descriptions of each source file aim
to provide a overview of the code at an intermediate level
of detail, highlighting the class structure,
summarizing the operation of each function,
and hopefully making it easier to examine and
understand the actual source code.

\ttsubsec {Assess.\{h,cpp\}}\label{sec:Anaalyze}

Header file \tt{Assess.h} includes \tt{Theory.h}
and defines:
\begin{itemize}
\item
    \tt{class Site}, holding information relevant to finding
    and examining self-intersections in observables.
    Data members are:
    \begin{itemize}
    \item
    \tt{ulong point}, recording a site's coordinates
    relative to the first symbol in the observable;
    \item
    \tt{short pos}, the site's position along the \tt{SymbStr};
    \item
    \tt{short nEs}, recording
    the number of preceding electric field insertions
    in the \tt{SymbStr};
    \item
    \tt{Symb in} and \tt{Symb out}, recording the \tt{symb}'s
    which precede and follow this location in the \tt{SymbStr}.
    \end{itemize}
    Member functions
    \tt{operator==} and
    \tt{operator<} are defining for testing equality
    or the relative ordering of coordinates with another \tt{Site}.

\item
    \tt{class Intersect}, which holds information relevant to
    choosing a pair of coinciding \tt{Site}'s at which to
    factor an observable into sub-observables.
    Data members are:
    \begin{itemize}
    \item
    \tt{float score}, discussed below;
    \item
    \tt{short pos1} and \tt{short pos2}, giving the
    two positions along the \tt{SymbStr} corresponding to
    an intersection at which the observable might be factored;
    \item
    \tt{short nEs},
    the number of intervening electric field insertions between
    \tt{pos1} and \tt{pos2}.
    \item
    \tt{static constexpr int maxtry}, which sets the
    maximum number of factorizations to try when
    attempting to approximate an observable.
    \end{itemize}
    The member function \tt{operator<} is defined
    for sorting of \tt{Intersect}'s.

\end{itemize}

File \tt{Assess.cpp} includes
\tt{Assess.h},
\tt{Global.h},
\tt{Commute.h},
\tt{Blab.h} and
\tt{Gripe.h},
and defines the following functions
(all of which are declared in \tt{Obs.h}):
\begin{itemize}
\item
    \tt{ObsList::assess}, whose single argument is an
    \tt{Obs} generated by the commutation routines.
    This function contains the logic
    used to decide whether an observable
    should be stored in the given \tt{ObsList},
    discarded, or approximated, if possible,
    by a product of observables already stored in the list.
    The appropriate disposition depends on the
    list in question, the phase of the program
    leading to the function call and, for the main
    list of canonicalized observables,
    whether the strong-coupling order of the observable
    (equal to the sum of its creation and expectation orders)
    falls within the limit set by \tt{global.maxord()}.
    While the commutation routines easily track
    the creation order of generated observables,
    determining the expectation of an observable
    is non-trivial.
    When needed, performing this classification is part
    of the task of this routine.
    This classification process, described below, can lead to
    recursive calls to the commutation routines, and hence to
    \tt{Obs::assess} as well, so coding allowing reentrant
    calls is essential.

    The basic steps are as follows:
    \begin{itemize}
    \item
	Depending on the setting of the
	\tt{canonicalize} data member of the
	\tt{ObsList}, the observable is either
	fully canonicalized, by calling \tt{Obs::canon},
	or, for gauge observables, cyclically rotated
	to a unique preferred starting location,
	by calling \tt{Obs::findstart}.
	(Note that these transformations are done
	in-place on the observable, which is passed by reference.)
    \item
	The hash table of the \tt{ObsList} is used
	to find the observable in the list, if already
	present.
	If so, a \tt{PolyTerm} is returned recording
	the list index of the observable
	plus a possible sign from canonicalization.
    \item
	If the observable is not found, then
	if the \tt{ObsList::frozen} function returns true,
	indicating further additions to the master observable
	list are disallowed, then the setting of the
	\tt{ObsList::approx} boolean controls
	whether the observable is discarded, signaled by the
	return of a vanishing \tt{PolyTerm}, or if the function
	\tt{Obs::approximate} is called to attempt
	construction of a suitable \tt{PolyTerm} approximation.
    \item
	When \tt{ObsList::frozen} returns false,
	if \tt{ObsList::classify} is set
	then the routine \tt{Obs::classify} is called to determine
	the expectation order of the observable.
	(The \tt{ObsList::classify} boolean is only set true for
	the main \tt{ObsList::obs} list.)
	If the classification function return value indicates that the
	observable's strong coupling order equals the target
	set by \tt{global.maxord()},
	then \tt{Obs::retain} is called to retain the observable
	for addition to the master observable list.
	Otherwise the observable is discarded.
    \item
	If \tt{frozen} and \tt{classify} are both false,
	then the observable is simply stored in the \tt{ObsList}.
    \end{itemize}
    In every case, a \tt{PolyTerm} is returned recording
    the list index of the observable, if found or stored,
    plus a possible sign from canonicalization.
    (For retained observables which are yet to be stored,
    the index in the returned \tt{PolyTerm} is just an
    arbitrary placeholder which will not be used.)
    A vanishing coefficient in the returned \tt{PolyTerm}
    signals that the observable is to be discarded.

    \item
	\tt{Obs::classify}, which attempts to determine
	the expectation order of an observable.
	For observables containing electric field insertions,
	this routine first calls \tt{Obs::noEbound},
	and for pure gauge observables
	\tt{Obs::u1bound} is called.
	These two routines, as described below, quickly
	generate lower-bounds on the expectation order of
	an observable which, very often, are sufficient to
	determine that the observable's strong-coupling order
	is larger than the limit set by \tt{global.maxord()},
	and hence full classification is not needed.
	But if the lower bound indicates that the observable
	might fall within the \tt{maxord} selected set, then
	either \tt{Obs::reduce}, for observables with electric
	fields, or \tt{Obs::factorize}, for observables without,
	are called to do a full classification.

    \item
	\tt{Obs::factorize}, which calls \tt{Obs::sitelist}
	to generate a sorted list of \tt{Site} objects
	corresponding to every location along the \tt{SymbStr}
	forming the observable.
	The \tt{factorize} routine then loops through
	all instances of coincident \tt{Site}'s.
	Each such self-intersection defines a possible split of
	the observable in which a sub-loop is excised.
	Both the sub-loop and the remainder of the observable
	(with sub-loop excised) are looked for in the list
	of known observables.
	If both are found, then the sum of their two
	expectation orders provides a possible expectation
	order for the initial observable.
	The observable's \tt{xorder} value is set to
	the minimum value from all such factorizations,
	provided this value is not larger than the
	observable's creation order (if set).
	For fermions, in addition to factorizations where
	a gauge sub-loop is extracted from the observable,
	it is also necessary to test three-way factorizations
	in which fermion/conjugate-fermion pairs are inserted
	at both cuts in the observable, transforming it into
	a product of three fermion bilinears.  This enables
	proper treatment of cases where the correct
	expectation order of the sub-loop comes not from
	fully tiling the sub-loop with plaquettes but
	instead involves fermion hopping contributions.
	(In other words, cases where perimeter-law effects
	compete or mix with area-law contributions.)

	This routine, as well as the following \tt{approximate}
	and \tt{noEbound} routines,
	uses a \tt{SymbStr} buffer in which to construct the terms
	emerging from a possible factorization or transformation,
	with each constructed \tt{SymbStr} sequence in the buffer
	later emplaced in an \tt{Obs}.
	Because these assessment routines can be called from different
	threads during the generation of geodesic equations
	(when multi-threading is enabled) this \tt{SymbStr} buffer,
	and its analogs in other assessment routines,
	cannot be declared \tt{static}.
	The \tt{SymbStr} buffer could be defined as a normal
	block-scoped variable, with space on the free store
	(for strings longer than the short string optimization
	size of \tt{std:string})
	dynamically allocated and deallocated by the
	\tt{std:string} routines on every pass through the block.
	Alternatively,
	this buffer could be declared \tt{thread_local}, which should
	minimize the number of needed free store allocation
	and deallocation operators, but at the cost of threading
	specific overhead associated with any thread local variable.
	Direct comparison (on the multi-core development systems used)
	found that for these assessment routines
	\tt{thread_local} storage of the needed \tt{SymbStr} buffers
	is preferable.
	Similarly, \tt{thread_local} storage is used for the
	\tt{vector<Site>} list returned by \tt{Obs::sitelist}
	(as well as the \tt{vector<Intersect>} used in the
	following \tt{approximate} routine).

    \item
	\tt{Obs::approximate} which, (similar to \tt{factorize})
	first calls \tt{Obs::sitelist} to generate a sorted list of \tt{Site} objects
	identifying every location along the \tt{SymbStr}
	forming the observable.
	The routine then loops through all instances
	of coincident \tt{Site}'s defining a
	possible split of the observable into sub-observables.
	Each such possible split is assigned a numeric score which
	is used to select preferred possible factorizations when
	there are multiple possibilities.
	This scoring is defined to prefer more even factorizations
	for pure gauge observables,
	increasing the likelihood that both factored sub-observables
	will already be present in the \tt{ObsList},
	while for fermion bilinears the scoring favors splits leading
	to larger gauge sub-loops with smaller fermion ``tails''
	which, for fermion bilinears shorter than the string-breaking
	scale, should also increase the likelihood that both
	sub-observables are already known.
	Smaller contributions to the score are designed to
	disfavor splits adjacent to electric field insertions,
	and at multiply-traversed links.

	After sorting the scored list of possible intersections,
	the most highly scored possible splits are sequentially tested
	to see if the two possible sub-observables are already known.
	To do so, two \tt{Obs} instances are constructed,
	one representing the loop inside the cuts and the other
	the remainder of the original observable.
	If both are already present in the \tt{ObsList},
	then a \tt{PolyTerm} with this factorization is returned.
	If no successful factorization is found after
	\tt{Intersect::maxtry} attempts, then the process aborts
	with failure indicated via a zero coefficient of
	the returned \tt{PolyTerm}.

    \item
	\tt{Obs::sitelist}, which is called by both
	\tt{Obs::factorize} and \tt{Obs::approximate}
	and prepares a \tt{vector<Site>} containing
	a sorted list of site coordinates (relative to an origin
	arbitrarily defined as the beginning of the first
	\tt{symb} in the \tt{SymbStr}) of every site
	traversed by the observable, together with the
	\tt{symb}'s leading into and out of each site,
	and the number of preceding electric field insertions.

    \item
	\tt{Obs::reduce}, which explicitly commutes the
	observable in question with primary \tt{Op}'s
	in order to identify which such operation leads
	to terms already present in the \tt{ObsList}
	and whose sum of expectation orders plus the
	order of the generator \tt{Op} is minimal.
	As discussed in Ref.~\cite{CSVAII}, properly
	executed, the result defines the expectation
	order of the initial observable.%
	\footnote
	    {%
	    There is a subtlety involving whether one
	    first collects terms resulting from the
	    commutation which, in translationally
	    invariant states, can lead to cancellations
	    for certain observables such as
	    [[\tt{AXAx}]] or \tt{AXXAxx}.
	    Such exact cancellations would not occur
	    in, for example, were the theory to have a varying
	    position-dependent coupling constant.
	    As such variations have no bearing on the
	    structure of the strong-coupling expansion,
	    the \tt{reduce} routine treats the separate
	    terms in these special-case cancellations as
	    establishing a possible expectation order.
	    }
	For \tt{Eloop} observables (i.e., loops containing
	one electric field insertion), commutation with
	a single plaquette (or a primary Wilson line)
	is tested, thereby leading to no-$E$ loops
	which will already be known (up to the
	specified \tt{maxord()} limit), so if no
	term in the commutation result leads to a known loop,
	the expectation order of the initial \tt{Eloop} must
	lie beyond the \tt{maxord()} limit.
	For \tt{Fermion} observables, commutation with
	single link fermion hopping \tt{Op}'s
	(those whose staggering is such that they map phase space
	coordinates to phase space coordinates) are tested
	with, again, at least one such commutator evaluation
	necessarily leading to a lower (expectation) order
	fermion bilinear which will already be known if the original
	bilinear is destined to be retained.
	For \tt{Efermion} observables (i.e., fermion bilinears with
	one electric field insertion), both plaquette
	(\& Wilson line) generators as well as single link
	fermion hopping generators are tested, while
	for \tt{EEloop} observables (i.e., loops with two electric field
	insertions), in addition to commutation
	with basic plaquettes (\& Wilson lines), evaluation
	of the pure-loop terms which can emerge from commutation
	with single-$E$ plaquettes (and single-$E$ inserted
	Wilson lines) is also performed.

    \item
	\tt{Obs::u1bound}, which for pure gauge loops, generates
	a lower bound on the expectation order the loop would have
	in a $U(1)$ pure gauge theory which, in turn, provides
	a lower bound on the expectation order in the U($N$) theory.
	The routine first finds a bounding-box for the loop.
	Then for each $[ij]$ coordinate plane, it ``projectively''
	computes the net $U(1)$ flux which span each stack of
	$[ij]$-oriented plaquettes (i.e., those $[ij]$-oriented
	plaquettes having the same $i$- and $j$-coordinate values,
	irrespective of their locations in any other coordinate directions).
	The sum of absolute values of these fluxes, summed over both
	all $[ij]$ stacks and then over all $[ij]$ pairs, gives
	the $U(1)$ lower bound.  In two dimensions, this would
	be the correct $U(1)$ expectation order, but in higher
	dimensions this is only a lower bound as a true minimal
	spanning surface may have opposite signs of flux on
	differing plaquettes within the same coordinate plane stack.

    \item
	\tt{Obs::noEbound}, which removes electric field insertions
	from the observable and then looks to see if the resulting
	\tt{Loop} or \tt{Fermion} is already known.  If so, the
	expectation order order of the no-$E$ observable provides
	a lower bound on the expectation order of the initial
	observable.
	If the no-$E$ observable is not found, this implies that
	the expectation order of the initial observable is
	larger than the limit used in generating the current set
	of no-$E$ observables.
	(As discussed in Ref.~\cite{CSVAII}, this presumes that
	observables are generated progressively in order of
	increasing number of electric field insertions.)

\end{itemize}

\ttsubsec {Blab.\{h,cpp\}}\label{sec:Blab}

Header file \tt{Blab.h} includes \tt{Gordion.h} and  defines:
\begin{itemize}
\item
    \tt{enum BLAB}, whose members name the various
    major \tt{.ccp} files, and whose integer conversions
    serve as indices into the static array \tt{Blab::blablevel}
    whose values control the verbosity of different routines.

\item
    \tt{class Blab}, with inline static data members
    \tt{blablevel} and \tt{blabmap}
    holding verbosity control information, and
    static functions \tt{setblab}
    and \tt{resetblab}
    for initializing, setting and resetting
    verbosity levels, respectively.
\end{itemize}

File \tt{Blab.cpp} includes
\tt{Blab.h} and
\tt{Gripe.h},
and defines the functions \tt{Blab::setblab} and
\tt{Blab::resetblab}.

\ttsubsec {Build.\{h,cpp\}}\label{sec:Build}

Header file \tt{Build.h} includes
\tt{Gordion.h} defines:
\begin{itemize}
\item
    \tt{namespace Build}, which declares the
    various \tt{mk_*} and \tt{do_*} functions 
    described below involved in the
    the construction of observables,
    symbolic expressions for
    the geodesic equations,
    and derivatives of the Hamiltonian or free energy.

\end{itemize}

File \tt{Build.cpp} includes
\tt{Build.h},
\tt{Canon.h},
\tt{Commute.h},
\tt{Numerics.h},
\tt{Rep.h},
\tt{Save.h},
\tt{Blab.h} and
\tt{Gripe.h}, and defines the following functions:
\begin{itemize}
\item
    \tt{Build::clearobs} which, depending on the stage
    argument, removes either all observables, or all fermion
    observables, from both \tt{ObsList::obs} and \tt{Canon.cache},
    and calls \tt{Global::clearpolys} to clear the associated script polynomials.
\item
    \tt{Build::mk_obs} which
    clears any previously build observables for the current minimization stage
    beyond the specified target strong-coupling order,
    initializes the master observable list
    for the current minimization stage if needed
    and then progressively increments
    the maximal strong-coupling order of observables up
    to a specified target, generating on each pass all
    observables of a given strong-coupling order.
    Prior to this loop, booleans
    \tt{ObsList::obs.freeze} and
    \tt{Canon::cache.freeze} are set to \tt{true}
    so that additions to the observable list and canonicalization cache
    may be made, and
    \tt{ObsList::obs.approx} is set to \tt{false}.
    Within the loop,
    when \tt{global.stage == Global::Gauge}
    the routine
    \tt{Build::mk_loops} is called, followed 
    (for Hamiltonian theories) by
    \tt{Build::mk_Eloops} and \tt{Build::mk_EEloops}.
    Alternatively, when \tt{global.stage == Global::Fermi},
    a temporary list of primary conjugate momenta \tt{Op}'s is
    constructed for later use by the multi-threaded
    routines \tt{do_Fermion_bckt}
    and \tt{do_Efermion_bckt}.
    Then \tt{Build::mk_fermions} is called,
    followed (for Hamiltonian theories)
    by \tt{Build::mk_Efermions}.
    Finally
    \tt{ObsList::do_fermiinit} is called
    to generate the fermion initialization map.
    Upon the completion of the the observable generation loop,
    if the observable list has lengthened
    then previously constructed polynomial scripts are cleared,
    and \tt{global.mk_bcktlist} and
    \tt{numerics.initialize} 
    are called to update the
    relevant observable bucket list
    (used for geodesic script generation), and
    initialize the observable expectation value list
    \tt{numerics.vev}.
    Lastly,
    the temporary \tt{Obsset}'s
    \tt{newobs} and \tt{ObsList::inbox} are
    destroyed,
    \tt{ObsList::obs.approx} is reset to the value
    of \tt{global.approx},
    and \tt{ObsList::obs.freeze}
    and \tt{Canon::cache.freeze}
    are set to \tt{true}.

\item
    \tt{Build::mk_loops}, which loops over target creation orders
    running from half the specified strong-coupling order limit
    up to four less than the strong-coupling order limit,
    then runs a (normally parallelized) \tt{for_each} loop over
    buckets of observables calling \tt{do_Loop_bckt} to perform the
    generation of \tt{Loop} observables starting from
    a given bucket of lower order observables.
    Following this \tt{for_each} loop,
    any newly generated observables of the specified order
    are inserted into the master observable list
    \tt{ObsList::obs} and any new short observables
    are added to the canonicalization cache.
    This routine, along with the other \tt{mk_*} observable generation
    routines, also updates the value of \tt{nobsG} or \tt{nobsF}
    in \tt{ObsList::obs} to keep track of the total number
    of gauge and fermion observables.

\item
    \tt{Build::mk_Eloops}, which builds \tt{Eloop}
    observables using a similar procedure 
    as \tt{mk_loops}, but looping over target creation orders
    running from half the specified strong-coupling order limit
    up to either two less than the strong-coupling order limit if
    \tt{maxord} is four, or else four less than \tt{maxord},
    and calling \tt{do_Eloop_bckt} to perform the observable generation.

\item
    \tt{Build::mk_EEloops}, which builds \tt{EEloop}
    observables using a similar procedure 
    as \tt{mk_loops}, but looping over target creation orders
    running from two less than half the specified strong-coupling order limit
    up to four less than the strong-coupling order limit,
    and calling \tt{do_EEloop_bckt} to perform the observable generation.

\item
    \tt{Build::mk_fermions}, which builds \tt{Fermion}
    observables using the analogous procedure
    as \tt{mk_loops}, but with
    target creation orders
    running from half the specified strong-coupling order limit
    up to two less than the strong-coupling order limit,
    and calling \tt{do_Fermion_bckt} to perform the
    observable generation.

\item
    \tt{Build::mk_Efermions}, which builds \tt{Efermion}
    observables using a similar procedure 
    as \tt{mk_fermions}, but looping over target creation orders
    running from half the specified strong-coupling order limit
    up to one less than the strong-coupling order limit,
    and calling \tt{do_Efermion_bckt} to perform the observable generation.

\item
    \tt{Build::mk_grad}, which loops over all terms
    in \tt{global.info().Hterms}, evaluating their
    gradient with respect to all $T$-even vacuum
    representation generators
    (of either gauge or fermi type, depending on the current
    minimization stage), by calling
    \tt{commute_poly} for each term and each such generator.
    Results are stored in \tt{global.info().gradient}.

\item
    \tt{Build::mk_curv},
    which loops over all terms
    in \tt{global.info().Hterms}, evaluating their
    second variations with all generators
    in the specified symmetry representation
    (of either gauge or fermi type, depending on the current
    minimization stage),
    by calling \tt{commute_poly} for each term and one generator,
    and then calling \tt{commute_poly} with that intermediate
    result and the second generator.
    Time reversal symmetry implies that the resulting curvature
    matrix is block-diagonal; this is taken into account.
    A temporary \tt{ObsList} is used to store the intermediate
    results of the first commutation.
    For the vacuum representation,
    terms resulting from the first commutation may be canonicalized,
    but not for non-vacuum representations.
    The resulting curvature matrix expressions
    are stored in
    \tt{global.info().curv[global.repnum]}.

\item
    \tt{Build::mk_lagr},
    which evaluates the commutator of pairs of generators
    in the specified symmetry representation
    (of either gauge or fermi type, depending on the current
    minimization stage),
    by calling \tt{commute_gen}.
    Time reversal symmetry implies that the resulting Lagrange
    matrix is block off-diagonal; this is taken into account.
    Results are stored in \tt{global.info().lagr[global.repnum]}.
    A warning is given if
    the list of observables does not contain those which result
    from the commutator of any two generators, based on a presumption
    that spectrum results are likely unreliable
    if observable truncation affects the Lagrange bracket itself.
    (In such cases, suspending highest order generators in
    spectrum calculations may be appropriate.)
    To avoid such truncation in the Lagrange bracket,
    the observable truncation order should be
    at least four times the maximal generator creation order.

\item
    \tt{Build::mk_geos}, which evaluates the
    symbolic geodesic equations for observables
    of either pure gauge or fermionic type,
    depending on the minimization stage.
    To enable efficient multi-threading,
    the sum over observables is partitioned into
    multiple buckets, with a \tt{for_each} loop
    over the list of buckets calling
    \tt{Build::mk_geo_bkt} for each bucket of observables.
    As noted above, the setting of the preprocessor
    symbol \tt{PARALLEL} controls whether a
    single-threaded or multi-threaded \tt{for_each} loop
    is used.

\item
    \tt{Build::do_Loop_bckt}, which
    loops over all observables in a specified bucket,
    selecting those which are \tt{Loop}'s,
    and then over all primary \tt{Eloop} \tt{Op}'s
    (which are plaquettes with an electric field insertion,
    figure eights with an electric field on the crossover link,
    or related Wilson lines with one electric field insertion).
    For each such \tt{Op,Obs} pair whose creation orders
    sum to the target value and whose order classification
    is such that the commutator may include observables of the
    desired strong-coupling order,
    \tt{do_commute} is called to evaluate the commutator.
    Within \tt{do_commute}, \tt{Obs::assess} is called
    to handle the classification by expectation order
    for each observable appearing in the generated terms,
    and storage in the \tt{thread_local} \tt{ObsList::inbox}
    if the resulting strong coupling order equals
    the target set by \tt{global.maxord()}.
    After completing all commutations for the designated bucket
    of observables,
    a \tt{mutex} lock is acquired, protecting the
    merging of new additions to \tt{Build::newobs} set of
    observables, which will be merged into the
    master \tt{ObsList::obs} at the end of the observable
    generation pass.

\item
    \tt{Build::do_Eloop_bckt}, which generates
    new \tt{Eloop} observables starting from a given bucket
    of initial observables, in a manner
    analogous to that of \tt{do_Loop_bckt}.
    New \tt{Eloop}'s, of the specified creation and
    strong-coupling order, are generated either by
    by commuting an existing
    \tt{Loop} with a gauge kinetic (\tt{EE}) operator,
    or by commuting an existing \tt{Eloop} with a
    primary \tt{Eloop} \tt{Op}.
    The same \tt{mutex} protected procedure is used to
    merge new \tt{Eloop}'s back into \tt{Build::newobs}.

\item
    \tt{Build::do_EEloop_bckt}, which generates
    new \tt{EEloop} observables starting from a given
    bucket of initial observables, in a 
    manner analogous to that of \tt{do_Loop_bckt}.
    New \tt{EEloop}'s, of the specified creation and
    strong-coupling orders, are generated by
    by commuting an existing
    \tt{EEloop} with a
    primary \tt{Eloop} \tt{Op}.
    The same \tt{mutex} protected procedure is used to
    merge new \tt{Eloop}'s back into \tt{Build::newobs}.

\item
    \tt{Build::do_Fermion_bckt}, which generates
    new \tt{Fermion} observables starting from a given
    bucket of initial observables, in a 
    manner analogous to that of \tt{do_Loop_bckt}.
    New \tt{Fermion}'s, of the specified creation
    and strong-coupling orders, are generated by
    commuting an existing
    \tt{Fermion} with either a primary \tt{Eloop} \tt{Op} or
    a primary conjugate momentum \tt{Fermion} \tt{Op} (with, for Hamiltonian
    theories, implies the correct staggering so that the resulting
    \tt{Fermion} observables are phase space coordinates,
    not conjugate momenta).
    The same \tt{mutex} protected procedure is used to
    merge new \tt{Eloop}'s back into \tt{Build::newobs}.

\item
    \tt{Build::do_Efermion_bckt}, which generates
    new \tt{Efermion} observables starting from a given
    bucket of initial observables.
    %in an analogous manner as \tt{do_Eloop_bckt}.
    New \tt{Efermion}'s, of the specified creation
    and strong-coupling orders, are generated by
    commuting an existing
    \tt{Fermion} with a gauge kinetic \tt{EE} operator,
    or by commuting an
    \tt{Efermion} with a primary \tt{Eloop} \tt{Op} or
    primary conjugate momentum \tt{Fermion} \tt{Op}.
    The same \tt{mutex} protected procedure is used to
    merge new \tt{Eloop}'s back into \tt{Build::newobs}.

\item
    \tt{Build::do_geo_bxkt} which loops over each observable
    in its allocated bucket, and then over all
    time-reversal even generators in the vacuum representation,
    calling \tt{commute_poly} to evaluate the corresponding
    commutators and storing the results in
    \tt{global.info().geos[bcktnum]}.
    Following this,
    \tt{check_xorder} is called to test for
    inconsistency in the assignment of each observable's
    expectation order,
    and the number of terms of various orders are counted
    in the newly evaluated geodesic equations and added to
    the relevant geodesic statistics counters in \tt{global.count()}.
    Finally, if \tt{global.autosave} is set true, then
    \tt{Save::write_geo_bckt} is called to write the geodesic
    equations for the bucket of observables to disk.

\item
    \tt{Build::check_xorder}, which examines each term of a
    geodesic equation,
    checking to see if the sum of the expectation orders
    of the observables in the geodesic equation term,
    plus the order of the generator producing this term,
    is ever less than the assigned expectation order of the initial
    observable.
    If so, this indicates an inconsistency in the classification
    of this observable during the observable generation phase,
    which is reported.

\item
    \tt{Build::do_geostats}, which calls
    \tt{do_geostat_bckt} in a \tt{for_each} loop
    to reevaluate the geodesic equation statistics,
    including the number of polynomial terms and their breakdown
    by monomial order.

\item
    \tt{Build::do_geostat_bckt}, which scans the
    symbolic geodesic equations for a specified bucket of
    observables, adding the appropriate counts to
    the \tt{ngeos}, \tt{geoterms} and \tt{geotermord}
    elements of \tt{global.count()}.

\end{itemize}
The file also defines the static variables:
\begin{itemize}
\item
    \tt{int cord}, used to pass the target creation order
    to the various \tt{do_*_bckt} observable generation routines;
\item
    \tt{Obsset newobs}, used to hold a growing set of newly generated
    observables before their insertion into the master observable
    list at the end of each observable generation pass;
\item
    \tt{nector<Op> primepOps}, used to hold the set of primary
    conjugate momentum \tt{Op}'s, for use by 
    \tt{do_Fermion_bckt} and \tt{do_Efermion_bckt};
\item
    \tt{std::mutex obsmutex}, used to synchronize additions to
    \tt{Obsset newobs} from different threads;
\item
    \tt{constexpr int single_thread\,=\,1024}, which sets a
    threshold for when to switch to multi-threaded generation
    of further observables.
\end{itemize}

\ttsubsec {Canon.\{h,cpp\}}\label{sec:Canon}

Header file \tt{Canon.h} includes
\tt{Index.h} and
\tt{Obs.h},
and defines:
\begin{itemize}\advance\itemsep -3pt
\item
    \tt{constexpr} values
    \tt{loopchunksize} and \tt{specchunksize},
    defining the chunk sizes used for canonicalizing loop
    and non-loop observables, respectively.

\item
    \tt{constexpr} functions
    \tt{looptblsize} and \tt{spectblsize},
    computing compile-time dimensions of the
    canonicalization tables.

\item
    \tt{constexpr} function \tt{ipow}, providing
    compile-time evaluation of integer powers.

\item
    type alias \tt{SymmSet}, defined as \tt{vector<uint>}
    and whose instances hold a list of indices
    of some subset of symmetry transformations contained
    in \tt{Symm::list}.

\item
    \tt{class Node}, used for entries in the
    canonicalization tables, with each \tt{Node}
    instance containing a chunk score and associated
    symmetry subset index.

\item
    \tt{class CanonCache}, which inherits
    from \tt{std::hash<string,int>},
    contains one data member
    \tt{bool freeze} signaling when cache
    additions are allowed (only during
    single-threaded operation),
    and defines the member functions:
    \begin{itemize}
    \item
	\tt{purge}, which removes cache entries whose
	observable index is above a specified limit;
    \item
	\tt{load}, which loads suitable short members
	of an \tt{Obsset} of soon-to-be-stored observables
	into the cache;
    \item
	\tt{reload}, which loads all suitable short
	observables from \tt{ObsList::obs} into the cache,
    \end{itemize}
    and includes a \tt{friend ostream\& operator<<} declaration
    for printing the cache.

\item
    \tt{class Canon}, which only contains
    static functions and data members.
    The static data members are:
    \begin{itemize}
    \item
    \tt{inline vector<Node> looptable},
    the canonicalization table for pure gauge loops;
    \item
    \tt{inline vector<Node> spectable},
    the corresponding table used for other observables;
    \item
    \tt{inline Index<SymmSet> symmset},
    an indexed list of possible symmetry subsets;
    \item
    \tt{inline CanonCache cache},
    used for observables whose lengths are
    less than the size of the symbol chunks
    used for indexing into the main canonicalization tables.
    This class inherits from \tt{hash<string,int>}
    and contains one additional boolean \tt{freeze}
    which signals when additions to the cache are disallowed.
    \item
    \tt{inline struct Stats stats},
    holding various canonicalization-related statistics.
    \end{itemize}
    %
    Static functions in the class are:
    \begin{itemize}
    \item
    \tt{speccode}, which computes the
    \tt{spectable} index corresponding to
    a special observable chunk;
    \item
    \tt{looptblinit} and \tt{spectblinit}, for initializing
    the canonicalization tables;
    \item
    \tt{loopchunk} and \tt{specchunk}, used in those
    initialization routines for converting a table index
    back into the corresponding observable chunk.
    \end{itemize}

\end{itemize}

File \tt{Canon.cpp} includes
\tt{Canon.h},
\tt{Symm.h},
\tt{Global.h},
\tt{Gripe.h} and
\tt{Blab.h},
and defines the functions:
\begin{itemize}

\item
    \tt{Obs::canon}, which transforms an observable
    into its canonicalized form (in-place), returning
    a possible associated sign.
    This function is a key routine whose efficiency is critical.
    The general strategy was described above.
    In a bit more detail:

    Loop observables (only containing gauge links with no
    electric field insertions) that are at most
    \tt{loopchunksize} in length 
    are converted into lookup keys for
    \tt{Canon::cache}.
    If not found (a cache miss), then all possible
    symmetry transforms, including cyclic rotations,
    of the loop are generated, and
    the canonical form based on lexical ordering is selected.
    If \tt{Canon::cache.freeze} is false, the canonical form is
    stored in the \tt{ObsList::obs} observable list,
    and \tt{Canon::cache} entries are created which map
    all symmetry-equivalent forms of the observable into
    the correct \tt{ObsList::obs} entry.
    If need be, prior to storing \tt{Obs::factorize}
    is called to determine (and set) the expectation
    order for the short loop.
    (The canonicalization cache is frozen,
    along with the main \tt{ObsList::obs} list of
    observables, after the observable generation stage.)
    For loop observables longer than \tt{loopchunksize},
    a moving ``window'' of \tt{loopchunksize} is walked around
    the loop, with the lowest scoring \tt{Canon::looptable}
    entries for all such chunks
    providing the small subsets of symmetry transformations
    needing to be tested to identify the preferred
    starting position and symmetry transformation for the
    unique canonicalized form of the full loop.
    In the special case of a one-dimensional lattice,
    the loop over all possible starting positions is
    needless, and is omitted.

    ``Special'' observables (fermion bilinears
    or gauge loops containing electric field insertions)
    are treated similarly, except that only
    transformations which lead to symbol strings
    starting with an appropriate non-link symbol
    need be considered.
    The \tt{Canon::cache} is used for
    observables of length \tt{specchunksize} or less
    (with \tt{Obs::factorize} or \tt{Obs::reduce} called, if necessary,
    to determine the observable's expectation order prior
    to storing in the cache),
    while for longer observables \tt{specchunksize} chunks
    which either begin or end on the endpoints of a fermion bilinear,
    or on electric field insertions in gauge loops,
    are converted into \tt{Canon::spectable} indices,
    with the lowest scoring entries for these chunks then providing
    the possible subsets of symmetry transformations
    and starting positions
    needing to be tested to identify the preferred
    canonicalized form.
    Because (for Hamiltonian theories) observables may
    contain two electric field insertions,
    the indexing of a chunk starting with one electric field
    insertion must allow for the possibility of a second
    such insertion within the chunk; this makes the indexing
    of chunks for \tt{Canon::spectable} different and more
    involved than for \tt{Canon::looptable}, and leads to
    smaller values for \tt{specchunksize} than \tt{loopchunksize}.
    In the spacial case of a Hamiltonian theory on a one
    dimensional lattice, Gauss' law implies that the position
    an electric field insertion may be freely moved along
    a gauge observable.
    This freedom is used to position both electric field
    insertions in \tt{EEloop}'s at the beginning of the
    observable.

\item
    \tt{Canon::looptblinit}, which initializes
    \tt{Canon::looptable}.
    For every valid sequence of link symbols (or ``chunk'')
    of length \tt{loopchunksize},
    all possible symmetry transforms of the chunk are generated
    and scored (based on lexical ordering of \tt{symb} strings)
    with the minimal score and associated
    symmetry subset (consisting of all symmetries which
    transform the initial chunk into the minimal score chunk)
    stored in the \tt{Node} of \tt{looptable} at the table index
    generated from the initial chunk.

\item
    \tt{Canon::spectblinit}, which initializes
    the special (i.e., non-loop) observable
    canonicalization table \tt{Canon::spectable}.
    For every valid \tt{specchunksize} sequence of symbols,
    necessarily starting with an initial special (a \tt{symb}
    representing either a fermion or a link with one or
    two electric field insertions),
    all possible symmetry transforms of the chunk are generated
    and scored, with the minimal score and associated
    symmetry subset (consisting of all symmetries which
    transform the initial chunk into the minimal score chunk)
    stored in the \tt{Node} of \tt{spectable} at the table index
    generated from initial chunk.

\item
    \tt{Canon::loopchunk} and \tt{Canon::specchunk},
    which are used inside the corresponding initialization functions
    to determine if a given canonicalization table index
    value corresponds to a valid symbol string.

\item
    \tt{Canon::speccode}, used inside \tt{Obs::canon}
    to compute the \tt{spectable} index corresponding
    to a possible starting chunk.

\item
    \tt{friend ostream\& operator<<}, for printing a \tt{CanonCache}.

\end{itemize}

\ttsubsec {Commute.\{h,cpp\}}\label{sec:Commute}

Header file \tt{Commute.h} includes
\tt{Poly.h} and
\tt{Gen.h}, and 
defines
\begin{itemize}
\item
    \tt{namespace Commute}, which
    declares the various functions
    in \tt{Commute.cpp} described next.
    This namespace also defines the trivial convenience
    functions:
    \begin{itemize}
    \item
	\tt{cordsum}, which sums the creation orders of an
	\tt{Op} and \tt{Obs}, but only if the \tt{Obs}
	creation order is set.
    \item
	\tt{commtype}, which returns the \tt{ObsType} resulting
	from the commutator of an \tt{Op} of any type with an
	\tt{EEloop} or \tt{Efermion}.
    \end{itemize}
\end{itemize}

File \tt{Commute.cpp} includes
\tt{Commute.h},
\tt{Global.h},
\tt{Splice.h},
\tt{Gripe.h}
\tt{Blab.h},
and contains the functions which
evaluate the commutator of generators with observables,
or of pairs of generators.
These functions (along with those in \tt{Canon.cpp})
form the most performance-critical heart of the program.
Functions defined in \tt{Commute.cpp} are:
\begin{itemize}
\item
    \tt{Commute::commute_poly}, which takes as arguments
    references to a coherence group generator (\tt{Gen}),
    an observable polynomial in either \tt{ObsPoly} or \tt{PolyMap} form,
    and an \tt{ObsPoly} for holding the result.
    The \tt{commute_poly} routine loops over the sequence
    of monomial terms in the supplied polynomials
    and calls \tt{commute_term} to do the indicated commutation.

\item
    \tt{Commute::commute_term}, which takes as arguments
    (references to)
    a \tt{Gen} and a \tt{PolyTerm}, whose commutator is to be
    evaluated, the \tt{ObsList} to which the \tt{PolyTerm} refers,
    and a \tt{PolyMap} to hold the result.
    The routine loops over the multiple factors to be acted upon
    in quadratic or higher order monomial terms.
    Within these loops, this function
    directly handles the special case
    (for Euclidean theories) of variations in entropy.
    For other cases, this routine then loops over the individual
    \tt{Op} terms in the generator and calls the further routine
    \tt{do_commite}.
    A second sub-function \tt{do_split} is also called
    to add the terms, potentially cubic in loops, which
    arise via operator ordering in the commutator of an \tt{Eloop}
    and an \tt{EEloop}.
    (See Ref.~\cite{CSVAII} for an explanation
    of how such terms arise via normal ordering.)

\item,
    \tt{Commute:do_commute}, which takes as arguments (references to)
    an \tt{Op}, \tt{Obs}, and \tt{PolyMap} for the result,
    plus a \tt{PolyTerm} and associated \tt{ObsList} representing
    further overall factors, not part of the commutation, to be 
    included in the result.
    Depending on the type of operator and observable,
    the routine calls the actual evaluation routines
    \tt{do_commuteA},
    \tt{do_commuteB},
    \tt{do_commuteC},
    \tt{do_commuteD} and
    \tt{do_commuteE},
    described below.

\item
    \tt{Commute::do_commuteA}, which is called when the
    generator \tt{Op} term is an \tt{Eloop}, and
    handles terms arising when that electric field operator
    differentiates a link (possibly containing an electric
    field insertion) in the
    observable which the \tt{Op} acts on.
    Loosely speaking, an \tt{Eloop} \tt{Op} splices
    the loop of the \tt{Op} into the observable at all
    possible locations where the electric field operator
    can act.
    The splice table, contained in \tt{Splice.h} and
    described below, encodes details of this splicing operation.
    The function \tt{SymbStr::join} is repeatedly called to
    handle the joining of one symbol string to another,
    appropriately taking care of canceling adjoining links
    and inverse links, or the combining of electric
    field insertions and associated links into the
    correct ``ligatures'' defined in the basic \tt{symb} encoding.
    \tt{ObsList::assess} is called for each resulting observable
    to determine, whenever possible, the correct expectation order
    for the observable and to decide (based on the
    values of \tt{global::maxord} and the the creation and expectation
    orders of the observable) whether the observable should be stored
    in the relevant observable list, approximated by products
    of lower order observables, or neglected.

    This routine, as well as the following commutation routines,
    uses a \tt{SymbStr} buffer in which to construct the various terms
    emerging from the commutator, with each constructed \tt{SymbStr}
    sequence in the buffer later emplaced in an \tt{Obs}.
    Because the commutation routines can be called from different
    threads during the generation of geodesic equations
    (when multi-threading is enabled) this \tt{SymbStr} buffer,
    and its analogs in other commutation routines,
    cannot be declared \tt{static}.
    The \tt{SymbStr} buffer may be defined as a normal
    block-scoped variable, with space on the free store
    (for strings longer than the short string optimization
    size of \tt{std:string})
    dynamically allocated and deallocated by the
    \tt{std:string} routines on every pass through the block.
    Alternatively,
    this buffer could be declared \tt{thread_local}, which should
    minimize the number of needed free store allocation
    and deallocation operators, but at the cost of threading
    specific overhead associated with thread local variables.
    In contrast to the situation found for the observable assessment
    routines, for the commutation routines
    direct comparison (on the multi-core development systems used)
    found that \tt{thread_local} storage of the needed \tt{SymbStr} buffers
    led to a performance penalty, not a speed-up,
    compared to the local scope alternative.
    This, presumably, reflects a
    higher ratio of buffer-related operations compared
    to allocation/deallocation overhead in the commutation routines.
    Hence, the \tt{SymbStr} buffers for these routines
    are all defined to have purely local scope.

\item
    \tt{Commute::do_commuteB}, which is called when the
    \tt{Obs} observable is an \tt{Eloop}, and handles
    terms arising when this electric field differentiates
    links (possibly with electric field insertions) in the generator \tt{Op}.
    \tt{ObsList::assess} is again called to determine
    the disposition of each generated term.

\item
    \tt{Commute::do_commuteC}, which is called when the
    \tt{Obs} observable is an \tt{EEloop}, and handles
    terms arising when the initial electric field operator of
    the \tt{Obs} differentiates links (possibly with an electric field)
    in the generator \tt{Op}.
    A complication in this routine is that the various splicings
    must be carried out in such a manner that the resulting symbol
    string, containing one or two electric field insertions,
    either ends up with an electric field inserted link as the
    initial symbol of the string if the \tt{Op} is an \tt{Eloop},
    but the initial symbol remains a fermion when the \tt{Op} is
    a fermion bilinear.

\item
    \tt{Commute::do_commuteD}, which is called when the \tt{Obs} is
    an \tt{EEloop} or \tt{Efermion},
    and handles terms arising when an electric field insertion
    not at the beginning of the \tt{Obs} differentiates
    a link in the generator \tt{Op}.

\item
    \tt{Commute::do_commuteE}, which is called when both the generator \tt{Op}
    term and the \tt{Obs} observable are both fermion bilinears, and handles
    terms arising when a fermion in one factor anti-commutes with a
    conjugate fermion in the other.
    This routine also handles the generation of the pure loop
    term which arises, via operator ordering, when both
    bilinears form closed loops.
    (Again, see Ref.~\cite{CSVAII} for details.)

\item
    \tt{Commute::do_split}, which constructs the 
    pure \tt{Loop} terms, potentially up to cubic order,
    which arise due to normal-ordering
    in commutators of an \tt{Eloop} with an \tt{EEloop}.
    (See Ref.~\cite{CSVAII} for an explanation.)
    The routine loops over the terms in the generator reduction,
    calls \tt{do_commuteC} to evaluate the
    commutator of the \tt{EEloop} with either of the \tt{Loop} factors
    in the reduction term,
    and then calls \tt{do_inner} to perform the ``inner'' commutation
    of that intermediate result.
    A temporary \tt{ObsList} is used
    to hold (and collect terms from) the results of the first
    commutation step prior to the \tt{do_inner} final reduction.
    As with the \tt{SymbStr} buffer, a block-scoped temporary
    \tt{ObsList} proves preferable to \tt{thread_local} storage.

\item
    \tt{Commute::do_inner}, which performs the reduction of an \tt{Eloop}
    observable to the pure loops, or quadratic polynomial of loops,
    which result when the electric field acts on links of the observable
    itself.
    In other words, if $\mathrm{tr} (E_\ell U(C))$ denotes an operator
    in which a Wilson loop $U(C)$ has an electric field insertion
    on some link $\ell$, \tt{do_inner} evaluates the commutator
    $\mathrm{tr} [(E_\ell)_{ij} , \, U(C)_{ji} ]$.

\item
    \tt{Commute::commute_gen}, which evaluates the commutator of two
    generators by looping over the \tt{Op} terms in each
    and calling \tt{op_commute} to do the work, and finally
    calling \tt{Poly::collect} to collect terms in the result.

\item
    \tt{Commute::op_commute}, which evaluates the commutator of two
    \tt{Op}'s, with code very similar to that in
    \tt{do_commuteA} and \tt{do_commuteE}.

\end{itemize}

\ttsubsec {Counter.h}\label{sec:Counter}

Header file \tt{Counter.h} includes
\tt{Theory.h}
and defines:
\begin{itemize}
\item
    \tt{static constexpr int cachelinelen = 128}, which
    is used in the following alignment specification.
    Its value is set equal to the
    cache line size on Apple silicon
    (which is twice that of Intel processors);

\item
    \tt{class Counter}, which is defined with
    \tt{alignas (cahcelinelen)} and holds a single
    data member \tt{std::atomic<ulong> x}.
    Definitions are included for the functions
    \tt{operator++},
    \tt{operator=},
    \tt{operator+=},
    a conversion operator to \tt{ulong},
    and \tt{friend ostream\& operator<<}.
    These enable a \tt{Counter} to be used much
    as if it were a simple \tt{ulong}.
    Instances of \tt{Counter} are used for
    statistics counters 
    which may be updated by multiple threads.

\item
    \tt{class CountArr<N>}, a templated class
    which inherits from \tt{array<Counter,N>}
    and adds member functions:
    \begin{itemize}
    \item
	\tt{operator=}, for either assigning a
	single value to all counters in the array,
	or copying an equal-sized \tt{ulong} array
	into the array of counters;
    \item
	\tt{operator+=}, for incrementing all the counters
	in the array by the values in an
	equal-sized \tt{ulong} array.
    \end{itemize}

\item
    \tt{struct Counters}, whose instances will
    hold a variety of statistics \tt{Counter}'s
    which are incremented by the canonicalization, observable
    assessment and polynomial script generation routines,
    and which are reported by the \textit{print stats} and
    \textit{print geostats} commands.
    A member function \tt{cleargeostats} is defined to
    zero the geodesic-related counters.

\end{itemize}

\ttsubsec {Coupling.\{h,cpp\}}\label{sec:Coupling}

Header file \tt{Coupling.h}
includes \tt{Obs.h} and defines:
\begin{itemize}
\item
    \tt{class Coupling}, which inherits from \tt{char8},
    holding the name of the adjustable coupling constant,
    and adds data members:
    \begin{itemize}
    \item
	\tt{doub value}, holding the numerical value of the coupling;
    \item
	\tt{int stage}, indicating of the coupling is used (first)
	in the gauge or fermion minimization stage.
    \end{itemize}

    The static data member
    \tt{vector<Coupling> list} holds the list of all
    defined adjustable coupling constants, while
    the static member function
    \tt{indx} takes a possible coupling name as argument
    (as either \tt{string} or \tt{char8})
    and returns the index into \tt{Coupling::list} of the specified
    coupling, if present, or else $-1$.
    A \tt{friend ostream\& operator<<} definition
    is included for printing the \tt{Coupling} name.

\item
    type alias \tt{Couplings = vector<Coupling>}.

\item
    \tt{struct Factor}, holding an index into
    \tt{Coupling::list} plus a power of the indicated
    coupling constant.

\item
    \tt{class Coeff}, inheriting from
    \tt{vector<Factor>} and holding a list of
    adjustable coupling and associated powers.
    One member function,
    \tt{doub operator()}, returns the
    numerical value of the product of powers of adjustable couplings.

\item
    \tt{class AdjTerm},
    whose instances
    represent polynomials in observables with a common
    overall adjustable coefficient.
    Class data members are:
    \begin{itemize}
    \item
	\tt{Coeff coeff}, the adjsutable coefficient for the term;
    \item
	\tt{ObsPoly poly}, the non-canonicalized form of the observable polynomial
	for this term;
    \item
	\tt{ObsPoly cpoly}, the canonicalized form of the observable polynomial.
    \end{itemize}
    A \tt{friend ostream\& operator<<} declaration
    is included for printing an \tt{AdjTerm} vector.
\end{itemize}

File \tt{Coupling.cpp} includes \tt{Global.h} and defines
\tt{ostream\& operator<<} for \tt{AdjTerm} vectors.
If the coupling exponent is non-zero, it outputs the coupling name
and associated exponent (if not unity), and then prints
the term's observable polynomial (parenthesized).

\ttsubsec {Data.h}\label{sec:Data}

Header file \tt{Data.h} includes
\tt{Gordion.h} and
\tt{Gripe.h},
and defines:
\begin{itemize}
\item
    \tt{struct ObsHdr}, with data members
    \tt{short corder},
    \tt{short xorder},
    \tt{uchar type},
    \tt{uchar pad} and
    \tt{ushort len},
    which will record metadata for an \tt{Obs},
    plus one member function \tt{data} is defined which returns
    a \tt{doub\&} reference representing the entire struct.
\item
    \tt{struct OpHdr}, with data members
    \tt{short order},
    \tt{uchar\,type},
    \tt{uchar prim},
    \tt{ushort pad} and
    \tt{ushort len},
    which will record metadata for an \tt{Op},
    plus the same \tt{data} member function.
\item
    \tt{struct GenRec}, with data members
    \tt{short rep},
    \tt{short order},
    \tt{uchar type},
    \tt{uchar T_odd} and
    \tt{ushort len},
    which will record metadata for a \tt{Gen},
    plus the same \tt{data} member function.
\item
    \tt{struct GradRec}, with data members
    \tt{ushort gen},
    \tt{ushort Hterm},
    \tt{ushort pad} and
    \tt{ushort len},
    which will record metadata for a gradient polynomial,
    plus the same \tt{data} member function.
\item
    \tt{struct CurvRec}, with data members
    \tt{ushort gen2},
    \tt{ushort gen1},
    \tt{ushort Hterm} and
    \tt{ushort len},
    which will record metadata for a curvature polynomial,
    plus the same \tt{data} member function.
\item
    \tt{struct LagrRec}, with data members
    \tt{ushort gen2},
    \tt{ushort gen1},
    \tt{ushort pad} and
    \tt{ushort len},
    which will record metadata for a Lagrange bracket polynomial,
    plus the same \tt{data} member function.
\item
    \tt{struct GeoRec}, with data members
    \tt{uint\,indx},
    \tt{ushort gen} and
    \tt{ushort\,len},
    which will record metadata for a geodesic equation polynomial,
    plus the same \tt{data} member function.
\item
    \tt{struct RecHdr}, with data members
    \tt{ushort info[3]} and
    \tt{ushort len},
    plus the same \tt{data} member function.
    Conversion operators
    to any of the above specialized \tt{*Hdr} types,
    and constructors which accept any of the above
    \tt{*Hdr} objects are also defined.
    A \tt{struct RecHdr} effectively functions as a
    union of the above specific record header types,
    but is not defined as such because unions cannot
    serve as base classes in C++.
\item
    \tt{enum class RecordID}, whose
    enumerated values,
    \tt{Null}, \tt{Obs}, \tt{Op}, \tt{Gen},
    \tt{Grad}, \tt{Curv}, \tt{Lagr}, \tt{Geos}, \tt{Vev},
    \tt{Coup} and \tt{Stat} serve to identify the various specific
    data record types
    (with \tt{Null} identifying an uninitialized \tt{RecIndx});
\item
    \tt{class RecIndx}, with data members:
    \begin{itemize}
    \item
    \tt{RecordID id}, giving the \tt{RecIndx} type;
    \item
    \tt{uchar nslice}, recording the number of
    slices for \tt{PolyCube} data (only used by
    \tt{Curv} data records);
    \item
    \tt{ushort ncol}, recording the number of
    columns for muilti-dimensional polynomial script data records
    (used by \tt{Grad}, \tt{Curv}, \tt{Lagr} and \tt{Geos}
    data records);
    \item
    an unnamed union making
    \tt{uint nrow} and \tt{uint items}
    equivalent synonyms, used to record either the
    total number of items for one-dimensional data records
    (\tt{Obs}, \tt{Op}, \tt{Gen}, \tt{Vev}, \tt{Coup}) or the number of
    rows for the multi-dimensional polynomial script data records;
    \item
    \tt{ulong reclen},
    recording the total length of the data record
    in units of \tt{sizeof (RecHdr)}
    (which is the same as \tt{sizeof (doub)});
    \item
    \tt{ulong filepos},
    holding the file offset of the data record
    when written to disk.
    \end{itemize}
    The static data array
    \tt{idname[]} maps \tt{RecordID} values to printable
    names, and
    one member function \tt{items} returns
    the total number of items for 
    multi-dimensional data records
    (equal to $\tt{nslice}  * \tt{ncol} * \tt{nrow}$).
\item
    \tt{class RecIndxArr<N>}, a templated class
    inheriting from \tt{array<RecIndx,N>}.
    The member function \tt{next} 
    returns a reference to the first uninitialized
    element of the array;
\item
    \tt{class DataRec}, which inherits from
    \tt{vector<RecHdr>} and adds the data member
    \tt{reference_wrapper<RecIndx> indexref}
    which associates the \tt{DataRec} with a
    \tt{RecIndx} entry in
    the \tt{RecIndxArr global.sysindex}.

    One constructor is defined, taking a
    \tt{RecordID} as argument.
    Member functions in class \tt{DataRec} are:
    \begin{itemize}
    \item
	\tt{entry}, which returns the reference
	bound to \tt{indexref};
    \item
	\tt{recptr}, which returns the actual
	data starting address (i.e., \tt{data()}) cast
	as a \tt{char*} or \tt{const char*} pointer;
    \item
	\tt{writerec}, which takes an \tt{fstream} argument,
	stores the data length and
	current file position reported by
	\tt{fstream::tellp} in the \tt{RecIndx} returned
	by \tt{entry()}, and writes
	the record data to the supplied file stream;
    \item
	\tt{readrec}, which takes an \tt{fstream} argument,
	resizes the \tt{DataRec} to the size of
	\tt{entry().reclen},
	seeks the input file to the position in
	\tt{entry().filepos}, and reads the stored data from the
	stream into resized data vector.
    \item
	\tt{clear}, which clears the \tt{vector<RecHdr>},
	sets \tt{items} to zero,
	and sets \tt{nrow} and \tt{nslice} to one.
    \item
	\tt{free}, which destroys the \tt{vector<RecHdr>},
	forcibly freeing its memory using the C++ swap-trick.
    \end{itemize}
\end{itemize}

\ttsubsec {Gen.\{h,cpp\}}\label{sec:Gen}

Header file \tt{Gen.h} includes
\tt{Obs.h} and
\tt{Op.h},
and defines:
\begin{itemize}
\item
    \tt{class Gen}, which inherits from \tt{OpSum}
    and encodes a linear combination of coherence group operators
    defined in the associated \tt{OpList}, which altogether define
    a symmetry-projected coherence group generator.
    The \tt{Gen} class adds a number of data members:
    \begin{itemize}
    \item
	\tt{short order}, giving the grading order of the generator
	(which must be identical to the orders of all the individual terms);
    \item
	\tt{OpType type}, giving the type of the generator (which must
	also be identical to the types of all the individual terms);
    \item
	\tt{bool T_odd}, which indicates whether or not the
	(anti-Hermitian) generator is
	time-reversal odd, and hence does not serve as a generator of deformations
	within the large $N$ classical configuration space;
    \item
	\tt{bool imag}, indicating if the generator
	(which is necessarily defined to be anti-Hermitian)
	has an imaginary coefficient.
    \item
	\tt{bool active}, indicating if the generator is active or suspended
	(default true).
    \item
	\tt{doub coeff}, an overall coefficient
	(allowing a common factor to be pulled out of the
	individual terms in the sum);
    \item
	\tt{ObsPoly reduction}, which for \tt{Eloop} generators
	stores the quadratic polynomial in loops which results
	when the electric field operator in each generator term
	is interpreted as acting on the generator itself,
	while for fermionic generators the reduction
	is the result of anticommuting the fermion operators
	at either end (if they reside at the same site).
	These generator reductions are used in evaluating
	the variation of the entropy (for Euclidean theories), and
	in the terms involving products of pure loops which can contribute
	to the commutator of an \tt{Eloop} with an \tt{EEloop}
	(in Hamiltonian theories).
    \end{itemize}
    The \tt{Gen} class has constructors which take a single \tt{Op}
    plus a \tt{Proj} into which the \tt{Op} is projected,
    or take an \tt{OpSum} linear combination of \tt{Op}'s plus a projector.
    Member functions are:
    \begin{itemize}
    \item
	\tt{collect}, for collecting terms referring to the same \tt{Op}
	within the \tt{OpSum};
    \item
	\tt{settype}, which sets the type, order and imaginary status of
	a \tt{Gen} based on information in a single \tt{Op};
    \item
	\tt{normalize}, which normalizes a generator in a particular
	representation,
	as described in section \ref{sec:gennorm}.
    \item
	\tt{inner_commute}, which evaluates the ``inner commutator''
	defining the \tt{ObsPoly reduction} of the generator;
    \item
	\tt{allzero}, which tests for identical vanishing of the \tt{Gen};
    \item
	\tt{valid}, \tt{is_Loop}, \tt{is_Eloop},
	\tt{is_Fermion} and \tt{isgauge},
	which are trivial convenience functions.
    \end{itemize}
    Static data members are:
    \begin{itemize}
    \item
	\tt{inline bool autoToddgens}, default \tt{true},
	indicating whether
	\tt{Eloop} generators are to be automatically
	created from each pure-loop generator, by
	evaluating the commutator of the loop with
	the gauge kinetic energy;
    \item
	\tt{bool gennorm}, indicating whether generators
	should be automatically normalized (as described
	in \ref{sec:gennorm}),
    \end{itemize}
    while static class functions are:
    \begin{itemize}
    \item
	\tt{project}, which projects an \tt{Op}, or \tt{OpSum} of \tt{Op}'s,
	onto all defined irreducible representations, storing the resulting
	symmetry projected generator in the \tt{gen} vector of the
	\tt{Rep} instance (provided that result is both non-zero and linearly
	independent from already present generators);
    \item
	\tt{isnew}, which tests for linear dependency;
    \item
	\tt{addgen}, which symmetry projects an
	\tt{OpSum} onto all defined representations
	and adds resulting new independent generators
	to the set of defined generators;
    \item
	\tt{suspend_group}, which suspends all generators of a specified order;
    \item
	\tt{activate_group}, which activates all generators of a specified order;
    \item
	\tt{suspend_gen}, which suspends an individual generator
	of the currently active representation;
    \item
	\tt{activate_gen}, which activates all, or a specified individual generator,
	of the currently active representation;
    \item
	\tt{normalize}, which normalizes all defined generators
	in all representations;
    \item
	\tt{geninit}, which initializes low-order generators.
    \end{itemize}
    A \tt{friend ostream\& operator<<} declaration
    is included for printing \tt{Gen}'s.
\end{itemize}

File \tt{Gen.cpp} includes header files
\tt{Gen.h},
\tt{Global.h},
\tt{Commute.h},
\tt{Rep.h},
\tt{Gripe.h} and
\tt{Blab.h}, and defines the
above mentioned functions.

The function \tt{Gen::inner_commute} loops over the
\tt{Op} terms of the generator and, for \tt{Eloop} generators,
calls \tt{do_inner} to evaluate the ``inner'' commutator,
while for \tt{Fermion} generators it tests to see if
the \tt{Op}, with fermions removed, forms a closed loop.
In both cases, as in \tt{commute_poly}, a \tt{PolyMap} is
used to collect common terms.

The initialization function \tt{Gen::geninit}
defines one and two plaquette generators
when \tt{global.stage == Global::Gauge}
plus, for any compactified dimension,
straight line Polyakov loops and
single plaquette-inserted Polyakov loops.
When fermions are present
and \tt{global.stage == Global::Fermi},
this function defines single link fermionic generators.
These initializations include the definition of all primary
generators used for observable generation.

\ttsubsec {Global.\{h,cpp\}}\label{sec:Global}

Header file \tt{Global.h} includes
\tt{Counter.h}
\tt{Coupling.h}
\tt{Data.h},
\tt{Op.h} and
\tt{Version.h},
and defines:
\begin{itemize}
\item
    Preprocessor symbols \tt{FOR_EACH},
    \tt{TASK_ARENA} and \tt{FINALIZE}.
    If the preprocessor symbol \tt{PARALLEL} is defined,
    then several \tt{tbb} header files
    are included and \tt{FOR_EACH} is defined as
    \tt{tbb:parallel_for_each},
    the function macro
    \tt{TASK_ARENA} is defined to initialize
    a \tt{tbb::task_arena} with a specified thread limit and execute
    provided code within this arena, and
    \tt{FINALIZE} is defined to execute \tt{tbb::finalize}.
    Otherwise,
    \tt{FOR_EACH} is defined as \tt{std::for_each},
    \tt{TASK_ARENA} simply expands to the provided code argument,
    and \tt{FINALIZE} does nothing.
    These definitions are used within \tt{Build.cpp}
    and \tt{Numerics.cpp}.
\item
    \tt{struct StageInfo}, with data members holding
    minimization stage specific information:
    \begin{itemize}
    \item
	\tt{short maxgen}, the maximum order of defined \tt{Gen}'s
	for the stage;
    \item
	\tt{short maxord}, the maximum strong
	coupling order of observables for the stage;
    \item
	\tt{uint MMAlimit}, specifying the number of initial observable
	expectation values, for the stage, to include in
	minimization numerics output;
    \item
	\tt{ObsSet MMAlist}, specifying a further set
	of observables whose expectation values will be
	included in minimization numerics output;
    \item
	\tt{OpList ops}, holding the list of defined
	operators for the stage;
    \item
	\tt{array<vector<Gen>,NREP> gens}, holding the list of
	defined generators, for the stage,
	in each symmetry representation;
    \item
	\tt{array<ushort,NREP> neven}, giving the number of
	time-reversal even generators, for the stage,
	defined in each symmetry representation;
    \item
	\tt{vector<AdjTerm> Hterms}, holding the definition
	of the Hamiltonian (or free energy) for the given stage,
	in the form of a sum of \tt{AdjTerm} instances.
    \item
	\tt{vector<uint3> bckt}, holding the list of
	tuples defining buckets of observables for the stage;
    \item
	\tt{Counters count}, holding the various
	thread-safe statistics counters defined in
	\tt{class Counters};
    \end{itemize}

\item
    \tt{static constexpr int MAXBCKT = 256}, 
    defining the maximum number of buckets of observables.
    This value is used for compile-time dimensioning of
    the \tt{geos} polynomial script arrays in \tt{global.stagedata}.
    (Note that any change to this compile-time constant
    will invalidate any previously generated system information files.)


\item
    \tt{struct SerialData}, whose instances will hold
    minimization stage specific data in a serialized form
    which may be moved directly to or from disk.
    Data members are:
    \begin{itemize}
    \item
	\tt{DataRec op}, which will temporarily hold a
	serialized version of the \tt{OpList}
	\tt{global.info().ops}
	during save or load operations;
    \item
	\tt{DataRec obs}, which will temporarily hold a
	serialized version of (the stage specific portion of)
	\tt{ObsList::obs} during save or load
	operations;
    \item
	\tt{DataRec gen}, which during
	save or load operations will hold a serialized
	version of a \tt{StageInfo::gens} array of
	generators;
    \item
	\tt{PolyRec grad}, which holds the polynomial
	scripts defining the gradient of the Hamiltonian
	(or free energy) for a given stage,
	given by commutator of each Hamiltonian term
	with each vacuum representation generator;
    \item
	\tt{PolyArr<NREP> curv}, which holds
	the polynomial scripts for the given minimization stage
	defining the curvature of the Hamiltonian
	(or free energy),
	given by the double commutators of each Hamiltonian term
	with all pairs of generators in
	each symmetry representation;
    \item
	\tt{PolyArr<NREP> lagr}, which holds the polynomial
	scripts defining the Lagrange bracket matrix,
	for each representation, given by commutators of all
	pairs of generators (of a given gauge or fermi type)
	in each symmetry representation;
    \item
	\tt{PolyArr<MAXBCKT> geos}, which holds
	the symbolic forms of the geodesic equations,
	for a given minimization stage,
	produced by commutators of each observable with
	every vacuum representation generator.
	The resulting polynomial scripts for each bucket
	of observables are stored in a separate \tt{PolyRec}
	data record;
    \item
	\tt{DataRec stat}, recording the size and possible
	on-disk location for the statistics counters in
	\tt{Coupling::list}.
    \end{itemize}
    The first \tt{DataRec} is used to save or restore
    all defined operators (for a given minimization stage),
    and likewise the following records for all stored observables,
    defined generators, the polynomial scripts for gradient,
    curvature, Lagrange bracket and geodesic equations,
    observable expectation values, coupling constants,
    and statistics counters.
    For the curvature and Lagrange bracket polynomials scripts,
    a separate data record is used for each symmetry representation,
    and due to the potential size of geodesic equation
    polynomial script, separate \tt{DataRec}'s are used
    for the geodesic equations of individual buckets of
    observables, as defined by the bucket lists stored
    in the corresponding \tt{StageInfo} structure.

\item
    type alias
    \tt{SysIndex = RecIndxArr<2\,*\,(5\,+\,2\,*\,NREP\,+\,MAXBCKT)>}.
    The number of elements equals
    the number of \tt{DataRec} or \tt{PolyRec} items in \tt{Global::stagedata}.

\item
    \tt{class Global}, which has a single instance
    \tt{global} holding much of the core symbolic information needed
    for applying the coherent state minimization algorithm.
    Data members are:
    \begin{itemize}
    \item
	\tt{enum stage}, defined to have two values:
	\tt{Gauge = 0} and \tt{Fermi = 1}.
	This \tt{enum} is effectively a boolean
	identifying whether the gauge or fermion stage of
	the variational algorithm is active;
    \item
	\tt{SysIndex sysindex}, which serves as an index
	to various data records possibly written to disk.
	Each entry is a \tt{RecIndx} recording the type of data
	in an associated \tt{DataRec}, the length and
	associated dimensions of that data, and a sys info file
	offset for the data record;
    \item
	\tt{SerialData stagedata[2]}, holding stage-specific
	data which is in a serialized form allowing direct
	writing to/from disk;
    \item
	\tt{StageInfo stageinfo[2]}, holding stage-specific
	information which need not be written to disk by
	a \tt{save} operation, or is not in a serialized form
	allowing direct writing to disk;
    \item
	\tt{short repnum}, which identifies the
	currently selected representation for curvature
	and spectrum calculations
	(default 0);
    \item
	\tt{short approx}, which if non-zero selects the
	approximation scheme by which
	observables whose strong-coupling order exceeds the
	\tt{global.maxord()} limit should attempt to be approximated
	via a product of lower order observables
	(default \tt{0});
    \item
	\tt{uint maxthread} which, if non-zero, sets a
	limit on the maximum number of threads spawned
	for parallel execution (default 0). 
	Suitably adjusting this value depending
	on the number of available processor
	cores and the maximum memory available
	can improve performance when generating
	very large observable sets;
    \item
	\tt{bool autosave}, indicating whether a sys info file
	should be automatically written during geodesic
	script construction
	(default \tt{false});
    \item
	\tt{bool geoswap}, specifying whether 
	buckets of geodesic equation polynomials
	are to be individually swapped to and from disk as needed
	(default \tt{false});
    \item
	\tt{bool massreinit}, specifying whether Euclidean
	fermion bilinear expecation values are automatically
	reinitialized when the fermion mass is changed
	(default \tt{true});
    \item
	\tt{bool symcurv}, specifying whether
	to use the symmetrized curvature when computing
	the small oscillation spectrum
	(default \tt{true});
    \item
	\tt{bool okfermivev}, signalling whether
	fermion expectation values have been initialized
	and are consistent with current gauge vevs;
    \item
	\tt{bool oknegeig}, specifying whether
	the presence of negative curvature eigenvalues
	(for sectors other than Euclidean fermions)
	will cause minimization to abort
	(default \tt{false});
    \item
	\tt{bool vevappend}, specifying if a pre-existing
	vev data file is opened in append mode
	(default \tt{false});
    \item
	\tt{bool MMAappend}, specifying if a pre-existing
	results output file is opened in append mode
	(default \tt{true});
    \item
	\tt{atombool interrupt}, recording whether an interrupt
	signal has been received;
    \item
	\tt{string MMAdir}, specifying the directory for
	numerics output data files (default \tt{./MMA/\textit{THY}}
	with \textit{THY} the value of the Makefile variable \tt{THY});
    \item
	\tt{string savedir}, specifying the directory 
	sys info and vev data files (default \tt{./save/\textit{THY}},
	with \textit{THY} the value of the Makefile variable \tt{THY});
    \item
	\tt{string sysfile}, specifying the sys info file name
	(default is auto-constructed by \tt{dfltfilename});
    \item
	\tt{string vevfile}, specifying the vev data file name
	(default is auto-constructed by \tt{dfltfilename});
    \item
	\tt{string MMAfile}, specifying the results output
	file name
	(default is auto-constructed by \tt{dfltfilename});
    \item
	\tt{ofstream sysstream}, holding the output stream
	for an open sys info file;
    \item
	\tt{ofstream vevstream}, holding the output stream
	for an open vev data file;
    \item
	\tt{ofstream MMAstream}, holding the output stream
	for an open numerics results file;
    \item
	\tt{Version version}, holding the running progrma's
	version number.
    \end{itemize}

    Member functions of class \tt{Global} defined for convenience are:
    \begin{itemize}
    \item
	\tt{info}, which return a reference to the
	\tt{stageinfo} structure for either the currently
	active stage, or a specified stage;
    \item
	\tt{data}, which return a reference to the
	\tt{stagedata} structure for either the currently
	active stage, or a specified stage;
    \item
	\tt{maxgen}, which returns a reference to
	\tt{info().maxgen}, holding the maximal
	\tt{Gen} order for the current stage;
    \item
	\tt{maxord}, which returns a reference to
	\tt{info().maxord}, holding the maximal
	strong-coupling order for retained observables
	in the current stage;
    \item
	\tt{count}, which returns a reference to
	\tt{info().count}, holding the statistics counters
	for the current stage;
    \item
	\tt{nobs}, which returns the number of retained
	observables for the current stage;
    \item
	\tt{fg}, which returns either an 'f' or 'g' character
	based on the current minimization stage (used in labeling
	of generators in printed output).

    \item
	\tt{dfltfilename}, which creates default
	data file names used for numerical output data files
	or save files, if not overridden by the user.
	The default file name is
	constructed from the short-form theory name,
	the maximum order of generators,
	the maximum strong-coupling
	order of the generated observables for the current stage,
	and the observable approximation setting,
	followed by the supplied extension 
	(which is \tt{.m} for numerical results files,
	\tt{.sys} for sys info files, 
	or \tt{.vev} for vev data files).

    \item
	\tt{bckt_pos}, which takes an observable number
	as argument and returns a tuple with the stage, bucket number,
	and position within the bucket at which the
	the specified observable is located in the bucket list.

    \item
	\tt{mk_bcktlist}, which prepares a
	\tt{vector<uint3>} bucket list, for a given minimization stage,
	with each element specifying bucket number,
	and starting and ending observable index values
	for the bucket.
	If the number of observables is less than
	\tt{1024 * MAXBCKTS} then
	the number of observables per bucket is set to 1024.
	Otherwise, the total number of observables are evenly
	divided among \tt{MAXBCKTS} buckets.
	(The last bucket need not be full.)
	This list is used in the parallelizing \tt{for_each}
	loops in the various \tt{Build::mk_*} routines as well as 
	\tt{Build::do_geostats} and
	\tt{Numerics::do_dvev},
	plus the \tt{Save} routines which read and write
	geodesic polynomial scripts..

    \item
	\tt{clearpolys} which clears all gradient, curvature,
	Lagrange bracket and geodesic equation polynomials for a
	specified minimization stage.

    \item
	\tt{stageinit}, which sets
	\tt{global.stage} to the indicated argument.

    \item
	\tt{sysfilename}, returning either a user-defined
	or default sys info file name;
    \item
	\tt{vevfilename}, returning either a user-defined
	or default vev data file name;
    \item
	\tt{MMAfilename}, returning either a user-defined
	or default MMA results file name.

    \end{itemize}

    If preprocessor symbol \tt{PARALLEL} is defined, then
    the class also includes a static data member
    \tt{tbb::task_scheduler_handle handle},
    which is used by the \tt{FINALIZE} macro
    invoked in \tt{Parse::quit} to wait termination
    of worker threads upon receipt of an interrupt signal.

\item
    constructor \tt{DataRec::Datarec (RecIndx Id)}, which is defined
    out-of-line in this header file because the constructor must
    bind the \tt{indexref} member of any \tt{Datarec} to an entry in
    \tt{global.sysindex};
\item
    the templated 
    constructor \tt{PolyArr<N>::PolyArr (RecIndx Id)}, which is defined
    out-of-line in \tt{Global.h} because the constructor must
    bind the \tt{indexref} member of each \tt{PolyRec} array element to an entry in 
    \tt{global.sysindex}.
\end{itemize}

File \tt{Global.cpp} includes \tt{Global.h},
and defines the member functions
\tt{dfltfilename},
\tt{bckt_pos},
\tt{mk_bcktlist} and
\tt{clearpolys}
described above.

\ttsubsec {Gordion.\{h,cpp\}}\label{sec:Gordion}

Header file \tt{Gordion.h}
contains \tt{using} declarations
for various standard container classes,
I/O classes and functions,
and basic types used through the code.
These are:
\begin{itemize}\advance\itemsep -9pt
\item \tt{std::pair},
\item \tt{std::array},
\item \tt{std::string},
\item \tt{std::vector},
\item \tt{std::complex},
\item \tt{std::set},
\item \tt{std::unordered_set},
\item \tt{std::unordered_map},
\item \tt{std::exception},
\item \tt{std::cout},
\item \tt{std::flush},
\item \tt{std::fstream},
\item \tt{std::istream},
\item \tt{std::ostream},
\item \tt{std::ifstream},
\item \tt{std::ofstream},
\item \tt{std::istringstream},
\item \tt{std::ostringstream},
\end{itemize}
along with:
\begin{itemize}\advance\itemsep -9pt
\item
    \tt{real} which is set to \tt{float} or \tt{double} if
    preprocessor synbol \tt{VEV32} is or is not defined,
    respectively,
\item
    \tt{cmplx = complex<double>},
\item
    \tt{doub = double},
\item
    \tt{symb = signed char},
\item
    \tt{uchar = uint8_t},
\item
    \tt{ushort = uint16_t},
\item
    \tt{uint = uint32_t},
\item
    \tt{ulong = uint64_t},
\item
    \tt{uint2 = std::pair<uint,uint>},
\item
    \tt{uint3 = std::array<uint,3>},
\item
    \tt{char8 = std::array<char,8>},
\item
    \tt{strview = std::string_view},
\item
    \tt{hash = std::unordered_map<Key,Value...>}.
\end{itemize}
The final \tt{hash} alias is a templated definition
applicable to any \tt{Key} and \tt{Value} types.
A further template definition makes
\tt{cast_to} a synonym for \tt{reinterpret_cast}.

This header file also defines:
\begin{itemize}\advance\itemsep -8pt
\item
    \tt{static string program}, which records the invoking program name,
\item
    \tt{static string cmdargs}, a usage message string,
\item
    \tt{static constexpr ulong ipow}, a compile-time integer power function.
\end{itemize}

Finally, as some older compilation environments
lack the \tt{std::format} function
(despite having the rest of C++17 and C++20 functionality),
this file tests for the definition of
\tt{__cpp_lib_format_uchar}
and conditionally includes either
a \tt{using} definition for \tt{std::format},
or an stand-alone templated replacement for \tt{std::format}.

File \tt{Gordion.cpp} includes
\tt{Init.h}.
\tt{Parse.h} and
\tt{Global.h},
and contains the routines:
\begin{itemize}
\item
    \tt{process_flags}, which handles command line flags.
    A \tt{-f} flag enables one to override
    or change an initial start-up file of commands which,
    by default, is \tt{./.gordion}.

\item
    \tt{main}, which calls \tt{process_flags}
    followed by
    \tt{initialize} for global variable initialization,
    \tt{signal} to install \tt{Parse:sig_catch} as
    a \tt{SIGINT} signal handler,
    optionally calls \tt{Save::load_save} to load
    command line arguments naming
    sys-info or vev-data files,
    optionally calls \tt{read_input} to
    read commands from a start-up file,
    and then calls \tt{read_input} to
    read from \tt{std::cin}.
\end{itemize}
This file also defines two static variables:
\begin{itemize}
\item
    \tt{dfltstart}, which records the program start time,
\item
    \tt{infile}, which holds the current input file name.
\end{itemize}

\ttsubsec {Gripe.h}\label{sec:Gripe}

Header file \tt{Gripe.h} defines:
\begin{itemize}
\item
    \tt{class BadInput}, which inherits from
    \tt{std::runtime_error} and whose instances
    signal invalid user input.
    \tt{BadInput} exceptions are caught and
    reported at
    the highest level of the command interpreter,
    aborting further processing of any open command
    input files;

\item
    \tt{class Abort}, which inherits from
    \tt{std::runtime_error} and whose instances
    signal failure of convergence in Newton
    minimization or ODE integration.
    \tt{Abort} exceptions are reported, but do
    not terminate further processing from any open
    command input files;

\item
    \tt{gripe}, a shorthand inline function
    which throws a \tt{BadInput} exception with
    specified error message;

\item
    \tt{abort}, a shorthand inline function
    which throws an \tt{Abort} exception with
    specified error message;

\item
    \tt{oops}, a shorthand inline function
    which throws a \tt{std::logic_error} exception with
    specified error message.
    Such exceptions indicate programming errors.

\end{itemize}

\ttsubsec {Index.\{h,cpp\}}\label{sec:Index}

Header file \tt{Index.h} includes
\tt{Gordion.h}
and defines:
\begin{itemize}
\item
    \tt{class Index<T>}, which inherits from
    \tt{vector<T>} and adds one data member:
    \begin{itemize}
    \item
    \tt{hash<string,uint>  map},
    which serves to map \tt{T} objects,
    if present in the \tt{vector<T>}, to their list index;
    \end{itemize}
    plus member functions:
    \begin{itemize}
    \item
    \tt{Index<T>::store}, which takes either a \tt{T\& item}
    argument, for types \tt{T} which themselves inherit from
    \tt{std::string}, or a \tt{string key} plus \tt{T\& item} pair,
    and adds the item, if not already present,
    to the list and associated (key, list index) pair to the map,
    returning the list index of the item;
    \end{itemize}
    Instances of this template class are used to store
    indexed lists of \tt{Op}'s, \tt{Symm}'s, \tt{Rep}'s
    and \tt{SymmSet}'s.

\end{itemize}

\ttsubsec {Init.\{h,cpp\}}\label{sec:Init}

Header file \tt{Init.h} merely declares the initialization function
\tt{initialize}.

File \tt{Init.cpp} includes the header files
\tt{Canon.h},
\tt{Gen.h},
\tt{Numerics.h},
\tt{Rep.h} and
\tt{Blab.h},
defines the two global data structures:
\begin{itemize}\itemsep -3pt
\item
    \tt{Global global},
\item
    \tt{Numerics numerics},
\end{itemize}
and brings together the definitions of static data
members of certain classes, in particular:
\begin{itemize}\advance\itemsep -6pt
\item
    \tt{ObsList::obs}, the list of canonicalized observables,
\item
    \tt{ObsList::base},
    the list of basic (non-canonicalized) observables,
\item
    \tt{ObsList::redu},
    the list of observables appearing in generator reductions,
\item
    \tt{Proj::list},
    the list of symmetry projectors,
\item
    \tt{Symm::list},
    the indexed list of symmetry transformations,
\item
    \tt{Rep::list},
    the indexed list of irreducible representation projector rows,
\item
    \tt{Coupling::list},
    the list of adjustable couplings,
\item
    \tt{Symm::trans2indx},
    the symmetry transformation to symmetry index map.
\end{itemize}

Finally, this file defines one function
\tt{initialize} which
calls various class-specific initialization routines, namely:
\begin{itemize}\advance\itemsep -6pt
\item
    \tt{Global::resetblab}, which sets initial verbosity levels;
\item
    \tt{Theory::theoryinit}, which loads the theory name
    in the default directory names for output files
    and prints the selected theory name and program version;
\item
    \tt{Symm::symminit}, which initializes lattice symmetry data;
\item
    \tt{Rep::repinit}, which initializes irreducible representations;
\item
    \tt{Canon::looptblinit}, which initializes the loop
    canonicalization table;
\item
    \tt{Canon::spectblinit},
    which initializes the non-loop canonicalization table;
\item
    \tt{ObsList::base.obsinit}, which initializes the basic
    observable list;
\item
    \tt{Theory::theorydefn}, which defines the Hamiltonian
    or free energy of the theory (selected by the value of
    the Makefile defined variable \tt{THY});
\item
    \tt{Gen::geninit}, which defines initial sets of generators.
\end{itemize}
The initial values of
\tt{global.stage} and
\tt{numerics.rk} are also assigned.
The ordering of some of these initializations is critical,
which is why the these static initializations
are bought together in this one function.

\ttsubsec {Linalg.h}\label{sec:Linalg}

Header file \tt{Linalg.h} tests for the preprocessor
symbols \tt{ARMA} or \tt{EIGEN} and defines a number of type aliases
and inline wrapper functions which allow either the
\tt{Armadillo} or \tt{Eigen}
linear algebra packages to be used.
For the normal case of Armadillo, the definitions are:
\begin{itemize}\advance\itemsep -6pt
\item
    \tt{using} declarations for 
    \tt{arma::abs},
    \tt{arma::eig_gen},
    \tt{arma::eig_pair},
    \tt{arma::pinv},
    \tt{arma::rank},
    \tt{arma::index_max},
    \tt{arma::sort},
    \tt{arma::sqrt} and
    \tt{arma::svd};

\item
    type alias \tt{Uvec = arma::Col<uword>},
    for submatrix view index vectors;
\item
    type alias \tt{Rvec = arma::Col<real>},
    for numerical observable expectation values and
    derivatives of those expectation values;
\item
    type alias \tt{Dvec = arma::Col<doub>},
    for gradient and other numerical vectors;
\item
    type alias \tt{Cvec = arma::Col<cmplx>},
    for spectrum eigenvalues;
\item
    type alias \tt{Dmtx = arma::Mat<doub>},
    for numerical values of the curvature, metric
    and Lagrange bracket matrices;
\item
    type alias \tt{Cmtx = arma::Mat<cmplx>},
    for spectrum eigenvectors;
\item
    \tt{ncol}, returning the number of columns of a matrix;
\item
    \tt{nrow}, returning the number of rows of a vector, matrix;
\item
    \tt{nelem}, returning the total number of elements of a vector or matrix;
\item
    \tt{memptr}, returning a pointer to data of a vector or matrix;
\item
    \tt{set_zero}, setting to zero all elements of a vector or matrix, possibly resized;
\item
    \tt{resize}, resizing a vector;
\item
    \tt{raw_print}, printing a vector or matrix to \tt{cout};
\item
    \tt{linsolve}, calling \tt{arma::solve} to solve a linear system;
\item
    \tt{infnorm}, calling \tt{arma::norm} to compute the $L^\infty$ norm of a supplied vector;
\item
    \tt{l2norm}, calling \tt{arma::norm} to compute the $L^2$ norm of a supplied vector;;
\item
    \tt{realpart}, calling \tt{arma::real} to return the real part of a supplied complex vector;
\item
    \tt{imagpart}, calling \tt{arma::imag} to return the imaginary part of a supplied vector;
\item
    \tt{has_nam}, calling \tt{arma::has_nam} on a supplied vector;
\item
    \tt{inv}, returning the inverse of a matrix;
\item
    \tt{transpose}, returning the transpose of a matrix;
\item
    \tt{aliasvec}, returning an \tt{Rvec} bound to pre-allocated memory;
\item
    \tt{sort_index}, returning a \tt{Uvec} holding the indices of sorted elements of a vector.
\end{itemize}
See the file for the alternative \tt{Eigen} type definitions and additional wrapper functions.
The capability to use \tt{Eigen} was added for testing purposes;
the \tt{Armadillo} library is generally preferred.

\ttsubsec {Numerics.\{h,cpp\}}\label{sec:Numerics}

Header file \tt{Numerics.h} includes
\tt{Global.h},
\tt{Poly.h} and
\tt{Ode.h},
and defines:
\begin{itemize}\advance\itemsep -6pt

\item
    \tt{struct Status},
    whose instances have data members:
    \begin{itemize}\advance\itemsep -2pt
    \item
	\tt{bool warn}, indicating if there are warnings to report;
    \item
	\tt{array<doub,3> negeig}, holding (up to 3) negative
	curvature eigenvalues;
    \item
	\tt{array<cmplx,3> cmplxeig}, holding (up to 3) complex
	curvature eigenvalues;
    \item
	\tt{doub maxloopv}, holding the largest \tt{Loop} expectation value;
    \item
	\tt{uint maxloopi}, holding the index of the largest
	\tt{Loop} expectation value;
    \end{itemize}

\item
    \tt{class Numerics},
    which has a single instance \tt{numerics}
    holding the information needed for numerical minimization.
    Data members of the class are:
    \begin{itemize}\advance\itemsep -2pt
    \item
	\tt{doub H}, holding the expectation value of the Hamiltonian
	(or free energy);
    \item
	\tt{Dvec gradient}, holding the gradient of the Hamiltonian
	(or free energy);
    \item
	\tt{Dmtx curvature}, holding the $T$-even curvature of the Hamiltonian
	(or free energy) in the currently selected representation:
    \item
	\tt{Dmtx metric}, holding the $T$-odd curvature of the Hamiltonian
	(or free energy) in the currently selected representation:
    \item
	\tt{Dmtx lagrange}, holding the Lagrange bracket
	matrix for the current irrep;
    \item
	\tt{Dvec delta}, holding the predicted location of the
	Hamiltonian (or free energy) minimum;
    \item
	\tt{Rvec vev}, for numerical
	expectation values of \tt{ObsList::obs} observables;
    \item
	\tt{Rvec dvev}, for derivatives of these
	expectation values in the direction \tt{delta};
    \item
	\tt{Cvec spectrum}, holding the small oscillation
	frequencies (from the latest spectrum evaluation);
    \item
	\tt{Cmtx modes}, holding the small oscillation
	eigenvectors (from the latest spectrum evaluation);
    \item
	\tt{Uvec Tevens}, holding the indices of active
	$T$-even generators in the currently selected representation;
    \item
	\tt{Uvec Todds}, holding the indices of active
	$T$-odd generators in the currently selected representation;
    \item
	\tt{short lastrep}, recording the representation
	index of the latest spectrum evaluation;
    \item
	\tt{Rvec vev_tmp}, a temporary vector used to
	combine updated gauge or fermion expectation values during
	geodesic equation integrations;
    \item
	\tt{const real* vev_buf}, pointing to
	an expectation value buffer;
    \item
	\tt{real* dvev_buf}, holding a pointer to
	an expectation value derivative buffer;
    \item
	\tt{doub svdcut}, defining the cutoff threshold for
	singular value omission in the curvature pseudo-inverse;
    \item
	\tt{doub dflttol}, defining a default value for the
	ODE and minimization tolerance;
    \item
	\tt{doub mintol}, defining the Newton minimization
	error tolerance;
    \item
	\tt{doub odetol}, defining the Runge-Kutta ODE
	error tolerance;
    \item
	\tt{uint odemax}, defining the maximum number of ODE
	integration steps;
    \item
	\tt{uint minmax}, defining the maximum number of
	Newton minimization steps;
    \item
	\tt{RKdef rk}, holding the selected Runge-Kutta method;
    \item
	\tt{struct \{ uint tries, uint fails \} stats}, holding
	minimization stats.
    \end{itemize}

    Member functions of class \tt{Numerics} are:
    \begin{itemize}
    \item
	\tt{do_minimize}, which runs a Newton minimization
	loop, calling \tt{do_step} to perform each iteration
	and terminating when the norm of \tt{numerics.delta},
	the predicted displacement from the minimum, falls below
	\tt{numerics.mintol} (or when the number of minimization steps
	exceeds \tt{numerics.minmax}).
	After the minimization, \tt{report_warns} is called to report
	any warnings.

    \item
	\tt{do_step}, which performs one Newton iteration step,
	first calling \tt{eval_delta} to find the predicted
	minimum location and then calling \tt{ode_int} to integrate the geodesic
	equations for all observables along the geodesic to the predicted minimum.

    \item
	\tt{do_flow}, which takes as arguments a coupling
	index, and an initial value, final value, and increment
	for the specified coupling.
	The indicated coupling's value is progressively changed from
	the given initial value to the final value in steps of the
	specified increment,
	with \tt{do_minimize} called on each step, followed by
	\tt{Numerics::save_data} to record the minimization results
	in the numerics output data file.

    \item
	\tt{eval_ham},
	which numerically evaluates the
	Hamiltonian or free energy by evaluating the
	canonicalized \tt{ObsPoly} defined in each
	\tt{AdjCoupling} term in
	\tt{global.info().Hterms},
	multipling by the indicated power of an
	adjustable coupling constant, and summing the resulting contributions.
	A boolean argument to the function controls whether the result
	is printed to \tt{cout}, or merely stored in \tt{numerics.H}
	for later reporting.

    \item
	\tt{eval_geos}, which
	calls \tt{eval_delta} to predict the minimum location $\delta$
	then calls \tt{Numerics::do_dvev} to evaluate numerically
	the symbolic geodesic equation, in the direction $\delta$,
	for each observable in \tt{ObsList::obs} of the relevant
	type (gauge or fermi), depending on \tt{global.stage},
	with results stored in \tt{numerics.dvev}.
	Finally, the routine optionally prints a specified subset of
	observable derivatives.

    \item
	\tt{eval_delta},
	which calls \tt{eval_grad} and \tt{eval_curv} to evaluate the
	gradient $dH$ and $T$-even curvature $ddH_{++}$ of the Hamiltonian
	(or free energy), and then solves
	the linear equation
	\begin{align}
	    ddH_{++} \cdot \delta = -dH ,
	\end{align}
	predicting the location $\delta$
	of the minimum of the Hamiltonian (or free energy).
	If \tt{numerics.svdcut} is zero then
	\tt{arma::solve} is used as the equation solver,
	otherwise \tt{arma::pinv} is called to compute
	a pseudo-inverse in which contributions from
	singular values less than \tt{numerics.svdcut} in
	the singular value decomposition of the curvature matrix
	are omitted.

    \item
	\tt{eval_grad}  which,
	zeros \tt{numerics.grad} and then
	loops over each \tt{AdjCoupling} term in
	\tt{global.info().Hterms},
	then over every $T$-even generator in the vacuum representation,
	evaluating numerically the corresponding \tt{ObsPoly} in the
	\tt{global.info().grad} entry which
	corresponds to the given $T$-even generator.
	Each such numerically evaluated \tt{ObsPoly} is multiplied
	by the indicated power of an adjustable coupling constant
	and added to the corresponding element of \tt{numerics.grad}.
	A boolean argument to the function controls whether the result
	is printed to \tt{cout}, or merely stored in \tt{numerics.grad}
	for later use.

    \item
	\tt{eval_curv} which zeros \tt{numerics.curv} and then
	loops over each \tt{AdjCoupling} term in
	\tt{global.info().Hterms},
	then over pairs of $T$-even generators in the current representation
	(and stage),
	evaluating the associated \tt{ObsPoly}
	element of \tt{global.info().curv[global.repnum]},
	multiplying by the indicated power of an adjustable coupling constant,
	and adding the result to the corresponding element of
	\tt{numerics.curv}.
	The resulting maxrix is stored in \tt{numerics.curv}.
	A second integer argument to the routine controls whether nothing,
	the curvature eigenvalues, or the curvature matrix plus eigenvalues
	are printed.

    \item
	\tt{eval_metr} which zeros \tt{numerics.metr} and then
	loops over each \tt{AdjCoupling} term in
	\tt{global.info().Hterms},
	then over pairs of $T$-odd generators in the current representation
	(and stage),
	evaluating the associated \tt{ObsPoly}
	element of \tt{global.info().curv[global.repnum]},
	multiplying by the indicated power of an adjustable coupling constant,
	and adding the result to the corresponding element of
	\tt{numerics.metr}.
	A second integer argument to the routine controls whether nothing,
	the metric matrix eigenvalues, or the metric matrix plus eigenvalues
	are printed.

    \item
	\tt{eval_lagr}, which zeros \tt{numerics.lagrange} and then
	loops over pairs of $T$-even/$T$-odd generators in the current
	representation (and stage),
	evaluating the associated \tt{ObsPoly} element of
	\tt{global.info().lagr[repnum]}
	and adding the result to the corresponding element of
	\tt{numerics.lagrange}.
	A boolean argument to the function controls whether the result
	is printed to \tt{cout}, or merely stored in \tt{numerics.lagrange}
	for use by \tt{eval_spectra}.

    \item
	\tt{eval_inuse}, which returns the index list of active generators
	for a specified representation.

    \item
	\tt{eval_spectra}, which calls \tt{eval_curv},
	\tt{eval_metr}
	and \tt{eval_lagr} to evaluate the curvature matrix
	$ddH_{++}$, metric matrix $ddH_{--}$ and Lagrange, 
	bracket matrix $L_{+-}$ for the specified representation,
	and then calls \tt{arma::eig_pair} to solve the
	generalized eigensystem:
	\begin{align}
	    ddH_{++} \cdot v
	    =
	    \omega^2 \; L_{+-} \cdot (ddH_{--})^{-1} \cdot (L_{+-})^T \cdot v \,.
	\end{align}
	The resulting small oscillation frequencies $\omega$ (times $\hbar$)
	give the excitation energies to translationally invariant
	excited states, or in other words, the mass spectrum
	for low-lying particles (either glueballs or mesons,
	depending on the minimization stage).
	If \tt{global.symcurv} is true, then the symmetrized
	curvature matrix is used in the computation (default).
	Eigenvalues are sorted into ascending order.
	Overloaded forms allow either a representation name
	or number as argument.
	A second boolean argument signals if the resulting
	eigenvalues are to be printed.

    \item
	\tt{initialize}, which
	initializes the numerics data structures
	for the indicated minimization stage.
	For the gauge stage,
	all elements of \tt{numerics.vev} are zeroed,
	(with the first element
	corresponding to the identity loop then reset to unity).
	For the fermion stage
	the last \tt{ObsList::obs.nobsF} elements of
	\tt{numerics.vev} are zeroed
	after which entries in
	\tt{ObsList::obs.fermiinit} are used to
	set the numerical value of each indicated closed \tt{Fermion}
	bilinear equal to the value of the corresponding \tt{Loop}
	times the uncorrelated fermion condensate
	(-1/2 for Hamiltonian theories, or minus the inverse fermion mass for
	Euclidean theories).

    \item
	\tt{write_data},
	which takes the value of a flowing coupling constant
	as argument and calls \tt{data_write} to write current values
	of the Hamiltonian or free energy,
	a subset (defined by
	\tt{global.info().MMAlimit})
	of observable expectation values of the relevant type,
	and, for Hamiltonian theories, the
	small oscillation spectra for all representations
	(for which the curvature and Lagrange bracket have
	been symbolically evaluated).
	Before writing data,
	if the output data stream is not already open,
	the output data file is opened in either
	truncation or append mode, depending on the value
	of \tt{dataappend}.

    \item
	\tt{status_rpt} which,
	for the gauge stage first calls \tt{check_loops},
	then calls \tt{curv_rpt} to report any
	problematic curvature eigenvalues,
	and then reports if any \tt{Loop} expectations exceed
	(in magnitude) 1.0.
	If there are eigenvalues
	more negative than \tt{-numerics.svdcut}
	(except for the fermion stage of Euclidean theories
	where extremization, not minimization, is relevant)
	and \tt{global.oknegeig} is not set, an exception is thrown
	which terminates the current command.

    \item
	\tt{check_loops}, which
	is called by \tt{status_rpt}
	(for gauge stage minimization)
	and stores the value and index of the \tt{Loop} observable 
	with maximal expectation value in \tt{numerics.warnings},
	setting \tt{numerics.warnings.warn} to true if the
	maximum exceeds 1.0.

    \item
	\tt{termvalue}, which returns the
	numerical value of a \tt{PolyTerm} using the list
	of expectation values stored in the \tt{vev} data member.
    \end{itemize}

    Static class members declared in class \tt{Numerics} are:
    \begin{itemize}
    \item
	\tt{Status status}, holding
	warning related information.

    \item
	\tt{termvalue}, a static overload which takes a \tt{PolyTerm}
	plus a list of numerical expectation values as arguments,
	and returns the
	numerical value of a \tt{PolyTerm} using the provided list.
	This version of \tt{termvalue} is used by \tt{do_dvev_bkt}.

    \item
	\tt{check_curv}, which is called by \tt{eval_delta}
	and \tt{eval_spectra},
	and examines the eigenvalues of the curvature matrix,
	storing in \tt{numerics.warnings}
	(for later reporting)
	those which are the most negative (if any) or
	having the largest imaginary parts (if any).

    \item
	\tt{curv_rpt}
	which, if \tt{numerics.status.warn} is set,
	reports negative or complex curvature eigenvalues.

    \item
	\tt{built_rep}, which checks if the symbolic forms of the curvature
	and Lagrange bracket are defined for a given representation.

    \item
	\tt{MMAform},
	which converts the standard output scientific notation for
	a \tt{double} to Mathematica's input form
	by first writing the number to a \tt{stringstream} 
	and then converting an
	\tt{e} exponent marker, if present, to
	\tt{\textasciicircum *}. So, for example,
	\tt{1.23e-04} becomes \tt{1.23\textasciicircum *-04}.

    \item
	\tt{do_dvev},
	which numerically evaluates the symbolic
	geodesic equations stored in \tt{global.geodesics},
	for all observables of the current stage (i.e., gauge or fermion),
	for the geodesic direction defined by \tt{numerics.delta},
	using a provided list of observable expectation values.
	Results are stored in a \tt{Numvec} whose reference is passed
	as an argument to \tt{eval_dvev}.
	This function is called by \tt{ode_int} for each
	sub-step comprising part of a single Runge-Kutta integration step.
	The full evaluation is broken up into separate buckets of
	observables, with a parallelizable \tt{for_each} loop 
	calling \tt{do_dvev_bkt} for each bucket.

    \item
	\tt{do_dvev_bkt}, which does the actual numerical evaluation
	of the geodesic equations for a bucket of observables.
	The list of observable expectation values to be used
	and the destination for results are passed to 
	\tt{do_dvev_bkt} via
	the pointers \tt{numerics.vev_buf} and \tt{numerics.dvev_buf},
	which are set by \tt{do_dvev}.

    \item
	\tt{err_norm},
	which evaluates the $L_2$ norm of a provided
	ODE error vector.
	This function is called by \tt{ode_int} for each
	Runge-Kutta integration step.

    \item
	\tt{data_write},
	which takes arguments
	of an output stream, a coupling value $v$,
	a variable name and corresponding variable value
	which (for different overloads) may be
	either a single \tt{doub},
	a \tt{Numvec} real-valued vector, or a
	\tt{Complxv} complex-valued vector.
	In each case, a line having the form of
	$$
	    \mathit{name}[\mathit{v}] = \mathit{value} \,;
	    \hspace*{2.6cm}
	    \vspace{-10pt}
	$$
	or
	$$
	    \mathit{name}[\mathit{v}] =
	    \{ \mathit{value}_1, \mathit{value}_2, \cdots \} \,;
	$$
	is written to the
	specified output stream.
	By design, this output is formatted in a manner which may
	be directly input to Mathematica.

    \end{itemize}
\end{itemize}

File \tt{Numerics.cpp} includes
\tt{Numerics.h},
\tt{Data.h},
\tt{Gen.h},
\tt{Global.h},
\tt{Rep.h},
\tt{Print.h},
\tt{Save.h},
\tt{Blab.h} and
\tt{Gripe.h}, and defines the above functions of class \tt{Numerics}.

\ttsubsec {Obs.\{h,cpp\}}\label{sec:Obs}

Header file \tt{Obs.h} includes
\tt{Symb.h},
\tt{Poly.h} and
\tt{Assess.h}, 
and defines:
\begin{itemize}
\item
    \tt{enum class ObsType}, with members
    \tt{Loop}, \tt{Eloop}, \tt{EEloop},
    \tt{Fermion}, \tt{Efermion} and \tt{Entropy},
    for identifying the different types of observables;

\item
    \tt{class Obs}, which inherits from \tt{SymbStr}.
    Instances of this class store the symbol string and related
    information representing an observable.
    Additional data members are:
    \begin{itemize}
    \item
    \tt{ObsType type}, giving the observable type;
    \item
    \tt{short corder}, recording the observable's
    creation order;
    \item
    \tt{short xorder}, recording the observable's
    expectation order;
    \item
    \tt{short midE}, which  records the location
    of a electric field insertion, if any, in a position other
    than the start of the \tt{SymbStr}.
    \end{itemize}

    Multiple constructors are defined, allowing an \tt{Obs} to be
    constructed from a printable string (whose characters are
    converted to their corresponding \tt{symb} codes, with
    \tt{symb} ligatures then appropriately joined),
    from a \tt{SymbStr} or some subrange of a \tt{SymbStr},
    all with creation and expectation orders also specified,
    or from an \tt{Op}.

    Member functions in \tt{class Obs} include:
    \begin{itemize}\advance\itemsep -2pt
    \item
    \tt{findstart}, which rotates a gauge observable
    to a unique preferred starting position when full canonicalization
    is not needed;
    \item
    \tt{middleE}, which finds and/or returns the location of a
    mid-observable electric field insertion;
    \item
    \tt{validate}, which examines the data in an \tt{Obs} instance
    for consistency (used for testing);
    \item
    \tt{trans}, which applies (in place) a symmetry transformation
    to the observable;
    \item
    \tt{order}, trivial convenience function returning the strong-coupling order
    (equal to the sum of creation and expectation orders);
    \item
    \tt{known_xord}, trivial convenience function signaling a valid \tt{xorder};
    \item
    \tt{bilinear}, convenience function signaling a \tt{Fermion} or \tt{Efermion} observable;
    \item
    \tt{EorEEloop}, convenience function signaling an \tt{Eloop} or \tt{EEloop} observable;
    \item
    \tt{midEtype}, convenience function signaling an \tt{EEloop} or \tt{Efermion} which may
    contain a \tt{midE} insertion;
    \item
    \tt{staggered}, convenience function signaling a staggered fermion bilinear
	(for Hamiltonian theories);
    \item
    \tt{is_gauge} and
    \tt{is_fermi},
	convenience functions signaling a pure gauge observable
	(including gauge entropy) or fermionic observable
	(including fermion contributions entropy and gauge kinetic energy),
	respectively;
    \item
    \tt{is_Loop}, \tt{is_Eloop}, \tt{is_EEloop}, \tt{is_Fermion}, \tt{is_Efermion}
    and \tt{is_Entropy},
    each signaling the indicated observable type;
    \item
    \tt{is_FermionO}, indicating if a fermion bilinear is odd length;
    \item
    \tt{imag}, indicating in Euclidean theories if a fermion bilinear has an imaginary
    expectation value;
    \item
    \tt{has_Es}, signaling an observable which has electric field insertions;
    \item
    \tt{is_coord}, signaling observables whose time reversal equals their
    Hermitian conjugate, and hence whose expectation values serve as
    coordinates on the classical configuration space;
    \item
    \tt{oddlen}, indicating fermion bilinears, possible with an electric field
    insertion, with an odd number of links.
    \end{itemize}

    Further \tt{class Obs} member functions are defined in
    \tt{Canon.cpp} and \tt{Assess.cpp}:
    \begin{itemize}\advance\itemsep -2pt
    \item
    \tt{canon}, for canonicalization of the observable;
    \item
    \tt{approximate}, which attempts to approximate an observable's
    expectation by a product of lower order observables;
    \item
    \tt{u1bound}, which constructs a lower bound on expectation order;
    \item
    \tt{noEbound}, which constructs a lower bound on expectation order;
    \item
    \tt{classify}, which determines an observable's expectation order;
    \item
    \tt{Ereduce}, which commutes an observable containing $E$'s with primary \tt{Op}'s
    to determine expectation order;
    \item
    \tt{factorize}, which factors an observable not containing $E$'s at
    self-intersections  to determine expectation order;
    \item
    \tt{sitelist}, which constructs a list of coordinates of all sites
    traversed by an observable.
    \end{itemize}

    An inline static boolean \tt{Obs::check} controls whether \tt{validate} is called
    on every created \tt{Obs}, and the \tt{static constexpr} array
    \tt{Obs::type_name} stores printable names of the different \tt{ObsType} values.

    The static function \tt{Obs::obstype}, called from
    \tt{Obs::Obs(string\&)} constructor, examines a string
    and determines which type of \tt{Obs}, if any,
    could have this printable representation;

    The static \tt{std::map<string,ObsType> obstypes} serves to map observable
    types as printable strings into the corresponding \tt{ObsType} \tt{enum} value.
    (This is used in \tt{Print::print_obs_select}.

    A \tt{friend ostream\& operator<<} declaration
    is included for printing \tt{Obs}'s.

\item
    \tt{struct Obshash}, which defines member function
    \tt{operator()(const SymbStr\&)}
    as a hash function
    for \tt{Obs}'s;

\item
    \tt{struct Obs_eq}, which defines the member function
    \tt{operator()(const SymbStr\&, const SymbStr\&)}
    as an equality testing function
    for \tt{Obs}'s;

\item
    type alias \tt{Obsset = std::unordered_set<Obs,Obshash,Obs_eq>};

\item
    type alias \tt{Obsmap = hash<const Obs,uint,Obshash,Obs_eq>};

\item
    \tt{class ObsList}, which inherits from \tt{vector<const Obs*>},
    and whose instances are used to store lists of observables.
    Additional data members are:
    \begin{itemize}
    \item
	\tt{Obsmap map}, a private data member holding a
	hash table of \tt{Obs}'s in the list.
    \item
	\tt{string name}, giving a name to the observable list;
    \item
	\tt{bool canonicalize}, indicating whether observables
	are to be canonicalized prior to storing in the \tt{ObsList};
    \item
	\tt{bool classify}, indicating whether observables
	are to be classified by expectation order prior to
	storing in the \tt{ObsList};
    \item
	\tt{bool approx}, indicating whether classified
	observables whose strong-coupling order exceeds the limit
	set by \tt{global.maxord()} should attempt to be approximated
	via a product of lower order observables;
    \item
	\tt{uint nobsG}, giving the count of contained gauge observables;
    \item
	\tt{uint nobsF}, giving the count of contained fermion observables;
    \end{itemize}

    Member functions include:
    \begin{itemize}
    \item
	\tt{operator()}, which takes a \tt{uint} list index as argument and returns
	a reference to the corresponding \tt{Obs} in the list;
    \item
	\tt{find}, which finds an observable with the specified \tt{SymbStr} if
	present in the list, returning the observable's index
	or \tt{UINT_MAX} for failure.
	If the \tt{ObsList} instance is \tt{ObsList::obs} and
	\tt{ObsList::inbox} is non-empty, then
	the inbox is also searched with a return of \tt{UINT_MAX-1}
	signaling the observable is in the inbox, but not yet
	inserted into the master \tt{ObsList};
    \item
	\tt{purge}, which removes entries from the \tt{ObsList} whose
	indices are above a specified value;
    \item
	\tt{empty}, which removes all entries (except the identity
	observable) from the list;
    \item
	\tt{insert}, which inserts the elements of an \tt{Obsset} into the
	\tt{ObsList};
    \item
	\tt{frozen}, which tests if the \tt{ObsList} instance is
	\tt{ObsList::obs} and if so returns the thread-specific value
	of \tt{ObsList::freeze};
    \item
	\tt{freezeif}, which tests if the \tt{ObsList} instance is
	\tt{ObsList::obs} and if so sets \tt{ObsList::freeze} to true,
	returning its previous value;
    \item
	\tt{refreezeif}, which tests if the \tt{ObsList} instance is
	\tt{ObsList::obs} and if so sets \tt{ObsList::freeze} to the
	supplied value;
    \item
	\tt{nobs}, which returns the number of stored observables
	for the specified minimization stage;
    \item
	\tt{neq}, which tests for inequality with a provided \tt{ObsList}
	reference argument, indicating a need to re-index observables from a
	different list;
    \item
	\tt{do_fermiinit}, which
	populates \tt{ObsList::fermiinit}.
	Each element of this vector is a pair
	of unsigned integers
	which serve to map the \tt{ObsList::obs} index
	of a \tt{Fermion} observable having coincident endpoints
	to the index of the
	corresponding \tt{Loop} observable.
    \item
	\tt{obsinit}, which initializes a list with the basic
	observables needed for the indicated stage,
	i.e., those appearing in the Hamiltonian or free energy.
	For the \tt{Gauge} stage, these are
	single plaquettes,
	Polyakov loops (if dimensions are compactified),
	and squares of electric field operators (for Hamiltonian theories),
	in all possible lattice orientations,
	along with entropy (for Euclidean theories).
	For the \tt{Fermi} stage, these are
	fermion mass and single link hopping terms,
	along with the fermion contribution to entropy (for Euclidean theories)
	or gauge kinetic energy (for Hamiltonian theories);
    \item
	\tt{print} for printing to a given output stream either the entire list,
	or a single indexed observable within the list.
    \end{itemize}

    Further member functions taking one or two \tt{Obs} as arguments are:
    \begin{itemize}
    \item
	\tt{store}, which adds its \tt{Obs} argument in the list,
	if not already present, and returns the index of the stored \tt{Obs};
    \item
	\tt{catalog}, which first canonicalizes or rotates the supplied \tt{Obs}
	argument(s) as appropriate,
	adds the resulting observable(s) to the list (if not already present)
	and returns a \tt{PolyTerm} which records the list indices
	of the observable(s).
	Throws an exception if the list is frozen but
	the supplied observable(s) are not present in the list;
    \item
	\tt{is_known}, which canonicalizes or rotates its r-value argument(s) as appropriate,
	then tests for the existence of the resulting observable(s) in the list,
	and returns a \tt{PolyTerm} which records the list indices if present;
    \item
	\tt{assess}, which is defined in \tt{Assess.cpp}
	and encodes the logic determining whether an observable
	generated by the commutation routines is to be stored, approximated via
	factorization, or discarded.
    \end{itemize}

    \tt{class ObsList} also defines static data members:
    \begin{itemize}
    \item
	\tt{ObsList ObsList::obs}, which is the master list of canonicalized
	observables;
    \item
	\tt{ObsList ObsList::base}, which holds non-canonicalized observables
	appearing in the definition of the Hamiltonian or free energy;
    \item
	\tt{ObsList ObsList::redu}, which holds non-canonicalized observables
	appearing in the ``inner commutation'' reduction of
	\tt{Eloop} and \tt{Fermion} generators;
    \item
	\tt{inline vector<uint2> ObsList::fermiinit}, containing
	index pairs which map list indices of
	those fermion bilinears whose fermion and conjugate fermion
	fields reside on the same site
	to the corresponding pure gauge loop
	(needed for initializing fermion bilinear expectation values).
    \item
	\tt{thread_local Obsset inbox}, which serves to hold 
	observables found by the \tt{do_*_bckt} observable generation routines
	prior to their insertion into the master \tt{ObsList::obs} observable
	list at the end of the observable generation pass;
    \item
	\tt{thread_local bool freeze}, signaling when newly generated
	observables, in a specific thread, are to be retained (or not),
    \end{itemize}
    and the static function:
    \begin{itemize}
    \item
	\tt{retain}, which adds a newly constructed and classified
	observable, of the current target order, to a thread's \tt{inbox}.
    \end{itemize}

\item
    \tt{class ObsStats}, which inherits from
    \tt{array<vector<vector<uint>>,...>},
    to hold an array of matrices which will collect
    statistics of observables classified according to their
    types, creation orders, and expectation orders.
    The class includes data members \tt{maxc} and \tt{maxx},
    for recording the maximum creation and expectation orders,
    respectively, appearing in the statistics matrices,
    \tt{avglen}, recording the average length of observables,
    and \tt{maxloop}, recording the index of the \tt{Loop}
    with largest expectation value.
    The inline member function \tt{get}, with two overloads,
    returns the number of observables with given creation
    and expectation orders, and optionally specified type.
    One constructor is defined, which takes an \tt{ObsList} as argument
    and steps through the observables accumulating the type-specific
    statistics.

\end{itemize}

File \tt{Obs.cpp} includes
\tt{Obs.h},
\tt{Global.h},
\tt{Numerics.h},
\tt{Op.h},
\tt{Symm.h},
\tt{Stats.h},
\tt{Gripe.h} and
\tt{Blab.h}, and defines
above mentioned member functions,
other than those indicated as residing
in \tt{Canon.cpp} or \tt{Assess.cpp},
along with the \tt{thread_local} variables
\tt{ObsList::inbox} and \tt{ObsList::freeze}.

\ttsubsec {Ode.\{h,cpp\}}\label{sec:Ode}

Header file \tt{Ode.h} includes
\tt{Gordion.h},
\tt{Linalg.h} and
\tt{RK.h},
and defines:
\begin{itemize}
\item
    type alias \tt{OdeRhs = void (*)(doub,\,const Numvec\&,\,Numvec\&)},
    for pointers to an ODE derivative function;
\item
    type alias \tt{OdeNorm = doub (*)(const\,Numvec\&,\,const\,Numvec\&,\,const\,Numvec\&)},
    for pointers to an ODE error norm function;
\item
    \tt{class Ode}, which inherits from \tt{class RKdef} and adds data members:
    \begin{itemize}
    \item
	\tt{OdeRhs rhs}, binding a derivative function to the ODE integrator;
    \item
	\tt{OdeNorm norm}, binding a norm function to the integrator;
    \item
	\tt{double tolerance}, giving the specified error tolerance;
    \item
	\tt{double minstepsize}, giving a minimum step size (default 0);
    \item
	\tt{double maxstepsize}, giving a maximum step size;
    \item
	\tt{double initstepsize}, giving an initial step size;
    \item
	\tt{uint steps}, recording the number of steps taken by the
	\tt{Ode} integrator instance;
    \item
	\tt{uint rejects}, recording the number of rejected steps attempted;
    \item
	\tt{uint maxstep}, giving an maximum number of ODE integration steps;
    \item
	\tt{static struct Stats stats}, holding cumulative ODE integration statistics.
    \item
	\tt{static constexpr int dfltmax = 1500}, the default value for
	\tt{maxstep}.
    \end{itemize}
    The class defines a single member function:
    \begin{itemize}
    \item
	\tt{integrate}, which takes as arguments initial and final values
	of the integration variable, and a \tt{Numvec} of initial data,
	and attempts the specified integration.
    \end{itemize}

    File \tt{Ode.cpp} includes \tt{Ode.h} and defines the
    \tt{Ode::integrate} member function.
    Static \tt{Numvec} buffers are used for derivative evaluations,
    trial function values, and the error estimate vector.
    The adaptive step size scheme is essentially that of
    Ref.~\cite{NumRecipe}, generalized to RK schemes of variable order.
    Both FSAL (``First Same As Last'') and non-FSAL schemes are supported.

\end{itemize}

\ttsubsec {Op.\{h,cpp\}}\label{sec:Op}

Header file \tt{Op.h} includes
\tt{Obs.h} and
\tt{Index.h},
and defines:
\begin{itemize}
\item
    type alias \tt{OpTerm = Term<uint>},
    used when representing a specific indexed member of
    \tt{Op::list} times a numeric coefficient.

\item
    \tt{class OpSum}, which inherits from
    \tt{vector<OpTerm>}, add a single data member,
    a \tt{reference_wrapper} to the associated \tt{OpList},
    and declares the member functions:
    \begin{itemize}
    \item
	\tt{oplist}, returning a reference to the
	associated \tt{OpList},
    \item
	\tt{collect}, for collecting terms with a common
	list index;
    \item
	\tt{flipT}, which returns an \tt{OpSum} containing
	entries in which all fermion \tt{Op}'s have flipped
	staggering;
    \item
	\tt{loop_dt}, which returns an \tt{OpSum} containing
	entries in which every
	\tt{Loop} \tt{Op} in the calling \tt{OpSum} has been
	converned into a sum of \tt{Eloop} \tt{Op}'s each containing
	one electric field insertion on some link of the loop.
	Specifically, given the lattice representation of some Wilson loop
	$U(C)$ around a contour $C$, \tt{loop_dt} returns the sum of
	\tt{Op}'s corresponding to the commutator
	$[\textstyle \frac 12 \mathrm{tr} (E^2), \mathrm{tr} \, U(C)]$.
    \end{itemize}
    Static versions of \tt{loop_dt} allow one to supply an individual
    \tt{Op} or \tt{OpTerm}.

    An \tt{friend ostream\& operator<<} declaration for printing \tt{Op}'s 
    is also included.

\item
    \tt{enum class OpType}, with members
    \tt{Loop}, \tt{Eloop}, \tt{Fermion} and \tt{Invalid},
    for identifying the different types of operators
    appearing in coherence group generators.

\item
    \tt{class Op}, which inherits from \tt{SymbStr}
    and whose instances store the form of individual terms
    appearing in coherence group generators.
    The class adds data members:
    \begin{itemize}
    \item
	\tt{OpType type}, indicating the type of the \tt{Op},
    \item
	\tt{bool primary}, indicating whether the \tt{Op} is primary,
    \item
	\tt{short order}, giving the grading order of the \tt{Op}.
    \end{itemize}
    Member functions are:
    \begin{itemize}
    \item
	\tt{findstart}, for cyclically rotating \tt{Loop} and \tt{Eloop}
	instances to their unique preferred starting position;
    \item
	\tt{oddlen}, indicating if the \tt{Op} has an odd length;
    \item
	\tt{is_Loop}, \tt{is_Eloop} and \tt{is_Fermion},
	trivial convenience functions
	each signaling the indicated operator type;
    \item
	\tt{is_FermionO}, indicating if an \tt{Op} is an odd-length fermion bilinear;
    \item
	\tt{staggered},
	identifying staggered fermion \tt{Op}'s (for Hamiltonian theories);
    \item
	\tt{is_coord}, identifying those \tt{Op}'s whose
	time reversal equals their Hermitian conjugate,
	and hence serve as coordinates on the classical
	configuration space.
    \item
	\tt{validate}, a private function called in the constructors
	for testing purposes;
    \end{itemize}
    One conversion operator is defined,
    \tt{Op(const Obs\&)},
    for converting an \tt{Obs}
    (of suitable type) into an \tt{Op}.

    The class includes the following
    static data members and static functions:
    \begin{itemize}
    \item
	\tt{inline bool check}, controlling whether \tt{validate}
	is called on every created \tt{Op};
    \item
	\tt{optype}, which is called from the
	\tt{Op::Op(string\&)} constructor and examines a string
	to determine which type of \tt{Op}, if any,
	could have this printable representation;
    \end{itemize}
    An \tt{friend ostream\& operator<<} declaration for printing \tt{Op}'s 
    is also included.

\item
    \tt{class OpList}, inheriting from \tt{Index<Op>} and whose
    instances hold indexed lists of operators.
    The clas s defines member functions:
    \begin{itemize}
    \item
	\tt{store}, which stores an \tt{Op}, if new, in the list;
    \item
	\tt{setprimary}, which determines which of the \tt{Op}'s
	in the list are primary --- that is,
	are not produced from the commutator of any pair of
	lower order \tt{Op}'s whose grading orders sum to the
	order of a given \tt{Op};
    \item
	\tt{purge}, while is called from \tt{Op::setprimary}
	and \tt{Test::jacobi} and removes all \tt{Op}'s from
	the \tt{OpList} beyond a specified limit;
    \item
	\tt{clear}, which removes all entries from the \tt{OpList};

    \item
	\tt{print}, for printing
	either an individual \tt{Op} or the complete \tt{Op::list}.

    \end{itemize}
\end{itemize}

File \tt{Op.cpp} includes
\tt{Op.h},
\tt{Commute.cpp},
\tt{Global.h} and
\tt{Gripe.h}, and defines the above mentioned class functions.

\ttsubsec {Parse.\{h,cpp\}}\label{sec:Parse}

Header file \tt{Parse.h} includes
\tt{Gordion.h}
and defines:
\begin{itemize}
\item
    \tt{namespace Parse}, which declares the
    command line input parsing functions described below.
    This namespace also defines the inline functions:
    \begin{itemize}
    \item
	\tt{eos}, which tests if an \tt{istringstream}
	is at the end of the string;
    \item
	\tt{parse_args}, a templated function which
	extracts any number of arguments, of types
	inferred from the argument list, from
	a provided \tt{istringstream};
    \item
	\tt{isword}, which compares an input word against a
	specified command word.  To be considered a valid match,
	the input word need merely match some non-empty
	initial portion of a command word, thereby allowing
	commands to be radically abbreviated;
    \item
	\tt{isstar}, which tests if the next word is a single \tt{*}.
    \end{itemize}
\end{itemize}

File \tt{Parse.cpp} includes
headers \tt{Parse.h},
\tt{Build.h},
\tt{Commute.h},
\tt{Numerics.h},
\tt{Print.h},
\tt{Rep.h},
\tt{Save.h},
\tt{Test.h},
\tt{Blab.h} and
\tt{Gripe.h},
and defines the functions performing
the main input command loop, subsequent parsing of input lines.
Specifically:
\begin{itemize}
\item
    \tt{Parse::read_input}, which runs a loop which
    optionally prompts,
    reads one line from the current input stream,
    optionally echos that line to \tt{std::cout},
    and then calls \tt{parse} to process the input line.
    Versions taking both l-value and r-value references
    to the input \tt{std::istream} are defined.

\item
    \tt{Parse::sig_catch}, which is called upon
    receipt of a \tt{SIG_INT} interrupt signal.
    and sets \tt{global.interrupt} to true
    and, if \tt{Parse::read_input} is awaiting
    input, directly calls \tt{Parse::quit}.

\item
    \tt{Parse::quit} which, for parallel execution,
    first calls \tt{tbb:finalize} to wait for all
    worker threads to complete and then,
    if \tt{timing} is set to \tt{true} and the elapsed time is
    greater than 10 seconds,
    prints the total (wall clock) time elapsed,
    and finally calls \tt{exit} to terminate the program.

\item
    \tt{Parse::parse_line}, which
    discards any initial
    white space or colon (:) characters
    and then breaks a given input line into
    separate commands separated by semicolons.
    For each resulting command,
    \tt{parse_cmd} is called after first
    discarding any trailing white space.
    Ignoring initial colons, which are used
    as a command prompt, allows one to copy
    an entire prior command line,
    including the initial prompt,
    and paste it in to be re-evaluated.

\item
    \tt{Parse::parse_cmd}, which reads an initial word from
    the provided command,
    compares it using \tt{isword}
    to the various command category keywords,
    and if matched calls one of the
    following keyword-specific parsing functions.
    (See appendix \ref{sec:commands} for a listing of the
    recognized commands.)

\item
    \tt{Parse::parse_read}, which interprets a
    \textit{read filename} command
    to process further command input from a specified file.

\item
    \tt{Parse::parse_write}, which interprets a
    \textit{write filename} command to direct
    further output to a specified file.
    If the provided argument is a single `-',
    then subsequent output is redirected back to \tt{std::cout}.

\item
    \tt{Parse::parse_load}, which interprets a
    \textit{load [filename]} command
    to load a binary save file (either a sys info
    or a vev data file).

\item
    \tt{Parse::parse_save}, which interprets a
    \textit{save [filename]} command to write
    either a binary sys info file containing information on
    all defined observables, operators, and generators,
    the polynomial scripts defining the
    gradient, curvature, Lagrange bracket
    and geodesic equations,
    or else write or append to a binary vev data file records
    recording numerical observable expectation values
    along with adjustable coupling constant values.

\item
    \tt{Parse::parse_set}, which interprets and carries out various
    \textit{set} commands allowing one to select the active
    irreducible symmetry representation,
    select the minimization stage (\textit{gauge} or \textit{fermi}),
    specify the numerical value of a symbolic coupling constant
    in the Hamiltonian (or free energy),
    set the ode and minimization error tolerances,
    select a Runge-Kutta method,
    specify the numerical output or save data file names
    and associated directories,,
    set the limit on observable expectation values to be
    written to the output data file,
    set various global booleans including
    \tt{dots},
    \tt{echo},
    \tt{timing},
    \tt{approx},
    \tt{autosave},
    \tt{Gen::autoToddgens}
    \tt{Obs::check},
    or set verbosity levels for functions in a particular
    source file.
    One place where abbreviated matching is not
    allowed is when entering the file name (without extension)
    in a verbosity setting command.

\item
    \tt{Parse::parse_reset}, which interprets and carries out
    \tt{reset} commands allowing one to return to their
    initial values
    the active irreducible symmetry representation,
    the minimization stage,
    the max thread limit,
    all verbosity levels,
    the ode or minimization error tolerances,
    the Runge-Kutta method,
    the output data file, and
    the numerics data structures (including
    observable expectation values).

\item
    \tt{Parse::parse_build}, which interprets
    \textit{build} commands
    and calls the \tt{mk_}\textit{stuff} functions in namespace \tt{Build}
    to construct observables with a specified maximum strong-coupling order,
    or construct the symbolic expressions for the geodesic equations,
    gradient and curvature of the Hamiltonian (or free energy), or
    Lagrange bracket (for Hamiltonian theories).

\item
    \tt{Parse::parse_eval}, which interprets
    \textit{evaluate} commands and calls the
    various \tt{eval_}\textit{stuff} functions in
    \tt{Numerics.cpp}
    to numerically evaluate
    the Hamiltonian or free energy,
    its gradient or curvature,
    the Lagrange bracket matrix,
    the predicted minimum location,
    or the small oscillation spectrum.

\item
    \tt{Parse::parse_do}, which interprets
    a \textit{do} command and, depending on the supplied argument(s),
    calls the function \tt{numerics.do_step}
    to perform one Newton minimization step,
    \tt{numerics.do_minimize} to perform iterative Newton minimization,
    or \tt{numerics.do_flow} to progressively increment a given
    coupling constant from an initial value to final value in specified steps,
    carrying out Newton minimization at each step and saving results
    to the numerics data file.

\item
    \tt{Parse::parse_gen}, which interprets a
    \textit{generator} command to add a new generator,
    suspend a previously defined generator or generator group,
    or activate a suspended generator or generator group.

\item
    \tt{Parse::parse_add}, which interprets an
    \textit {add generator} command to define
    additional generators.
    This deprecated function may be removed at some point.

\item
    \tt{Parse::parse_print}, which interprets a wide variety
    of \textit{print} commands, calling the various
    printing functions in \tt{Print.cpp} detailed below.

\item
    \tt{Parse::parse_test}, which interprets a \textit{test} command and,
    depending on the supplied argument,
    calls the function
    \tt{Test::jacobi} to test that results
    of the commutation routines satisfy the Jacobi identity,
    or \tt{Test::irreps} to
    test the orthogonality and completeness of the
    defined irreducible representation projectors.

\item
    \tt{Parse::parse_call}, which interprets various
    \textit {call} commands which allow one to
    enter arguments and directly call
    various \tt{Symm} and \tt{Obs}-related
    functions:
    \begin{itemize}
    \item
	\tt{canon}, to canonicalize a given observable;
    \item
	\tt{compose}, to compose two symmetry transformations;
    \item
	\tt{transform}, to apply a symmetry transformation to a given observable;
    \item
	\tt{u1bound}, to evaluate the U(1) flux bound on an observable's expectation order;
    \item
	\tt{noEbound}, to evaluate the electric field removal bound on an
	observable's expectation order;
    \item
	\tt{classify}, to classify an observable's expectation order;
    \item
	\tt{commute}, to evaluate the commutator of an indicated generator and observable;
    \item
	\tt{do_inner}, to evaluate the ``inner commutator'' of an \tt{Eloop}
	observable.
    \end{itemize}
    This functionality is present for debugging purposes, not normal use.
\item
    \tt{Parse::parse_help}, which prints a command summary
    in response to either \textit{help} or \textit{?}.

\end{itemize}
This file also defines and uses the static variables:
\begin{itemize}\advance\itemsep -5pt
\item
    \tt{bool echo}, which controls whether command input lines are echoed to
    \tt{std::cout};
\item
    \tt{bool timing}, which controls whether to report the time taken by
    each command (if greater than 10 milliseconds).
\item
    \tt{bool awaiting}, which lets the interrupt handler know
    when the program is awaiting input;
\item
    \tt{std::chrono::timepoint starttime}, which records
    the program start time.
\item
    \tt{ostream myout}, which holds a copy of the initial
    \tt{std::cout} output buffer;
\end{itemize}

\ttsubsec {Poly.\{h,cpp\}}\label{sec:Poly}

Header file \tt{Poly.h} includes
\tt{Data.h},
\tt{Term.h} and
\tt{Theory.h},
and defines:
\begin{itemize}
\item
    \tt{class Polyindx}, which inherits from
    \tt{std::array<uint,PSIZ>}.
    \tt{Polyindx} objects are used to store tuples
    of indices of observables contained in an \tt{ObsList}.
    As the first element of every \tt{ObsList} is always
    the identity, a \tt{Polyindx} can be used to represent
    linear, quadratic, or higher order monomials in observables,
    up to the limit set by \tt{PSIZ}.
    Member functions are:
    \begin{itemize}
    \item
	\tt{mysort}, which sorts an index array into decreasing order;
    \item
	\tt{validate}, which verifies proper index sorting (used for debugging);
    \item
	\tt{order}, which counts the number of non-zero indices
	and returns the resulting order of the monomial term
	(used for polynomial statistics).
    \end{itemize}
    Constructors are defined which allow one to specify
    0, 1, or 2 observable indices.
    With multiple indices, \tt{mysort} is called
    to place the contained indices in descending order.
    An \tt{friend ostream\& operator<<} for printing \tt{Polyindx} objects
    is also declared.

\item
    \tt{class PolyTerm}, which inherits from
    \tt{Term<real,Polyindx>} and represents an observable
    monomial, with \tt{real} coefficient, of any order up to \tt{PSIZ}.
    The \tt{PolyTerm} class adds member functions:
    \begin{itemize}
    \item
	\tt{operator[]} for subscripting the underlying \tt{Polyindx};
    \item
	\tt{operator*} for properly combining two \tt{PolyTerm} objects
	(provided the result is of no higher order than \tt{PSIZ}),
	or scaling the coefficient of a \tt{PolyTerm} by a specified factor;
    \item
	\tt{reindex} for re-indexing the observables referenced in a
	\tt{PolyTerm} to a different \tt{ObsList};
    \item
	\tt{order}, returning the monomial order;
    \item
	\tt{validate}, which verifies correct sorting of the \tt{Polyindx}
	(used for debugging);
    \item
	\tt{print}, which takes \tt{ostream} and \tt{ObsList} references
	as arguments, and prints the encoded monomial term to the stream.
    \end{itemize}

\item
    type alias \tt{Polyvec = vector<PolyTerm>}, used to hold lists of
    \tt{PolyTerm}'s.

\item
    \tt{struct Polyhash}, which declares a hashing function
    for use in \tt{class PolyMap}.

\item
    \tt{class PolyMap}, which inherits from
    \tt{hash<Polyindx,doub,Polyhash>} and whose instances
    encode polynomials up to order \tt{PSIZ} in a hash table
    so that common terms may be automatically combined.
    Each \tt{PolyMap} instance contains a
    reference (actually \tt{reference_wrapper})
    to the associated \tt{ObsList} of the indexed observables.
    The class includes the additional member functions:
    \begin{itemize}
    \item
	\tt{add},
	for adding a \tt{PolyTerm} to the \tt{PolyMap},
	either by creating a new entry
	in the hash table if the \tt{Polyindx} in the \tt{PolyTerm}
	has not previously been entered, or by simply changing the
	\tt{doub} value of an existing map entry with the same
	\tt{Polyindx} key;
    \item
	\tt{add_gen},
	for converting the \tt{OpSum} of a \tt{Gen} to
	\tt{PolyTerm}'s and adding them to the \tt{PolyMap};
    \item
	\tt{negate}, which negates the coefficients of all
	stored elements in the \tt{PolyMap};
    \item
	\tt{purge}, which removes any map entries with vanishing coefficients;
    \item
	\tt{allzero}, which purges the map and then tests for no
	remaining entries.
    \end{itemize}
    An \tt{friend ostream\& operator<<} declaration is included
    for printing \tt{PolyMap} instances.

\item
    \tt{class ObsPoly}, which inherits from
    \tt{PolyVec} and whose instances encode
    polynomials up to order \tt{PSIZ} in observables which are
    stored in a common \tt{ObsList}.
    Each \tt{ObsPoly} instance contains a reference the
    associated \tt{ObsList}.
    Member functions in the class are:
    \begin{itemize}\advance\itemsep -1pt
    \item
	\tt{push_map}, which adds all terms in a \tt{PolyMap} to
	an \tt{ObsPoly};
    \item
	\tt{sort}, which sorts the entries in an \tt{ObsPoly}
	based on the relative ordering of the index tuples of the
    individual terms;
    \item
	\tt{negate}, for negating an \tt{ObsPoly};
    \item
	\tt{allzero}, which tests for a vanishing \tt{ObsPoly};
    \item
	\tt{obslist}, which returns a reference to the
	associated \tt{ObsList}.
    \end{itemize}
    Constructors allow creating an empty \tt{ObsPoly}
    given just the \tt{ObsList} to which it is linked,
    or creating an \tt{ObsPoly} from the index of an
    already indexed observable in the specified \tt{ObsList}.
    A \tt{friend ostream\& operator<<} declaration
    is included for printing \tt{ObsPoly} objects.

\item
    \tt{struct Poly}, which inherits from
    \tt{RecHdr} and serves as the header separating strings
    of \tt{PolyTerm}'s in packed polynomial script records.
    The \tt{struct} adds the member functions:
    \begin{itemize}
    \item
	\tt{begin()}, which returns a \tt{const PolyTerm*}
	pointer to the \tt{PolyTerm} which immediately follows
	the \tt{Poly} in a polynomial script record;
    \item
	\tt{end()}, which returns a \tt{const PolyTerm*}
	pointing to the end of the set of
	\tt{PolyTerm}'s,
	whose number is recorded in the \tt{RecHdr},
	which follow the \tt{Poly}.
    \end{itemize}
    These definitions enable range-based \tt{for} loops
    to be used with a \tt{Poly}.
    A \tt{static} \tt{constexpr} \tt{int}
    \tt{ptermsize = sizeof\,(PolyTerm)\,/\,sizeof\,(Poly)}
    is also defined for convenience.
    A \tt{friend ostream\& operator<<} declaration
    is included for printing an observable polynomial
    encoded by the \tt{Poly} header plus following
    \tt{PolyTerm}'s.

\item
    \tt{class PolyRec}, which inherits from
    \tt{DataRec} and adds the data member:
    \begin{itemize}
    \item
    \tt{mutable vector<ulong> offset},
	which is initialized on demand by the polynomial
	printing function indexing specific polynomials
	in the \tt{PolyRec}
	and will record the
	offsets within the associated \tt{PolyRec} of individual 
	polynomials in the record.
    \end{itemize}
    The class also defines the further member functions and nested class:
    \begin{itemize}
    \item
	\tt{add}, which takes a \tt{RecHdr} and a \tt{PolyMap}
	as arguments and appends the encoded polynomial, with its
	associated \tt{Poly} header, to the \tt{PolyRec};
    \item
	\tt{clear}, which calls \tt{DataRec::clear} and also
	clears \tt{offset};
    \item
	\tt{offloaded}, which signals whether the data record
	has been offloaded to disk by comparing the
	\tt{entry().reclen} with the current vector size;
    \item
	\tt{operator()}, which takes one, two, or three integer arguments,
	interpreted as the row, column and slice indices of a rectangular
	1, 2 or 3 dimensional block of polynomials (with dimensions specified
	by the \tt{nrow}, \tt{ncol} and \tt{nslice} members of the 
	\tt{RecIndx} associated with the \tt{DataRec})
	and returns a reference to the indexed polynomial.
	The first use of \tt{operator()} for a given \tt{PolyBlk} triggers
	a call to \tt{slice_n_dice} to initialize the needed \tt{offset} array.
    \item
	\tt{slice_n_dice}, which
	populates the \tt{offset} array with the
	offsets (in \tt{RecHdr} units) of each individual
	\tt{Poly} starting a new polynomial
	in the \tt{PolyRec};
    \item
	\tt{class PolyIter}, an iterator
	for polynomials within a \tt{PolyRec}, with the data member
	and member functions:
	\begin{itemize}
	\item
	    \tt{const Poly* ptr}, which will point to a \tt{Poly};
	\item
	    \tt{operator*}, which returns a reference to the \tt{Poly}
	    to which \tt{ptr} points;
	\item
	    \tt{operator++}, which increments \tt{ptr} by the size
	    of the packed polynomial to which \tt{ptr} points.
	\item
	    \tt{operator!=}, which compares the \tt{ptr} member with
	    another \tt{PolyIter};
	\end{itemize}
    \item
	\tt{begin()}, which returns a \tt{PolyIter} pointing
	to the beginning of the \tt{PolyRec};
    \item
	\tt{end()}, which returns a \tt{PolyIter} pointing to
	the end of the \tt{PolyRec}.
    \end{itemize}
    The definitions of \tt{PolyIter}, \tt{begin()}, and \tt{end()}
    enable range-based \tt{for} loops to be used with a \tt{PolyRec}.

\item
    a templated type alias making
    \tt{PolyArr<N>} equivalent to \tt{array<PolyRec,N>},
    with a constructor taking a \tt{RecordID}.

\end{itemize}

File \tt{Poly.cpp} includes
\tt{Poly.h} and
\tt{Gen.h},
and defines the above functions 
including
\tt{Polyhash::operator()},
\tt{PolyTerm::operator*},
\tt{PolyMap::add},
the \tt{ObsPoly} constructors and member functions,
and the \tt{operator<<} functions for \tt{ObsPoly},
\tt{PolyMap} and \tt{Poly} objects.

\ttsubsec {Print.\{h,cpp\}}\label{sec:Print}

Header file \tt{Print.h} includes
\tt{Gordion.h}
and defines:
\begin{itemize}
\item
    \tt {namespace Print}, which includes the functions:
    \begin{itemize}
    \item
	\tt{is_one}, testing if a supplied \tt{doub} equals 1 to
	within numerical precision;
    \item
	\tt{is_int}, testing if a supplied \tt{doub} equals an
	integer to within numerical precision;
    \item
	plus all the various other printing functions described next.
    \end{itemize}
\end{itemize}

File \tt{Print.cpp} includes
\tt{Print.h},
\tt{Canon.h},
\tt{Gen.h},
\tt{Numerics.h},
\tt{Rep.h},
\tt{Save.h} and
\tt{Blab.h},
and defines a long list of printing functions
(all members of \tt{namespace Print}):
\begin{itemize}
\item
    \tt{Print::coeffprt}, for nicely printing coefficients of
    \tt{Poly}'s or other linear combinations;
\item
    \tt{Print::print_obs}, which prints one or more observables in
    the current canonicalized observable list \tt{ObsList::obs}.
    Differing overloads will print
    one observable (by giving its number in the list or
    typing its symbol string),
    a numeric range $[\mathrm{first},\mathrm{last}]$
    of observables, or all observables;

\item
    \tt{Print::print_obs_select}, which prints all observables
    in \tt{ObsList::obs} having specified creation and expectation orders
    and optionally a specified observable type;

\item
    \tt{Print::print_base}, which prints the contents of the
    non-canonicalized basic observable list \tt{ObsList::base};

\item
    \tt{Print::print_op}, which has several overloads to print all,
    a range, or one operator from the list of defined operators,
    \tt{Op::list};

\item
    \tt{Print::print_primary}, which prints the subset of primary \tt{Op}'s;

\item
    \tt{Print::print_fermiinit}, which prints \tt{ObsList::fermiinit},
    holding the mapping between ``closed'' fermion bilinears and Wilson loops,
    used for fermion expectation value initialization;

\item
    \tt{Print::print_gen}, which will print one, or all, defined generators
    in the currently selected representation and stage,
    \tt{global.info().gens[rep]}.
    For each generator,
    the output shows the current representation, the generator index
    prefixed with ``g'' for gauge or ``f'' for fermion generators
    and followed by a ``*'' suffix if the generator is $T$-odd,
    and then the actual generator definition.

\item
    \tt{Print::print_rep}, which prints the projection operators for 
    all, or a specified individual irreducible representation row in \tt{Rep::list}.
    Or, with no argument, reports the names of all defined \tt{Rep}'s
    and the currently active \tt{Rep};

\item
    \tt{Print::print_symm}, which will print the explicit symmetry transformations
    for either all or a particular specified symmetry in \tt{Symm::list};

\item
    \tt{Print::print_symmsets}, which prints the list of all symmetry subsets
    used in canonicalization, \tt{Canon::symmset};

\item
    \tt{Print::print_couplings}, which prints names and values of all adjustable
    coupling constants stored in \tt{Coupling::list};

\item
    \tt{Print::print_hamiltonian}, which prints (for Hamiltonian theories)
    the current active Hamiltonian, namely \tt{global.info().Hterms};

\item
    \tt{Print::print_freeenergy}, which prints (for Euclidean theories)
    the current active free energy, namely \tt{global.info().Hterms};

\item
    \tt{Print::print_grad}, with overloads which will print all,
    or a selected portion
    %(specified by either one or two indices),
    of the symbolic gradient matrix,
    \tt{global.info().grad};

\item
    \tt{Print::print_curv}, with overloads which will print all,
    or a selected portion (specified by either one, two, or three indices),
    of the symbolic curvature matrix for the currently active representation
    and stage, \tt{global.data().curv[rep]}.

\item
    \tt{Print::print_lagrange}, with overloads which will print all,
    or a selected portion (specified by one or two indices),
    of the Lagrange matrix 
    for the currently active representation and stage,
    \tt{global.data().lagr[rep]};

\item
    \tt{Print::print_spectrum}, which prints the
    small oscillation frequencies stored in
    \tt{numerics.spectrum} and previously evaluated by an
    \tt{evaluate spectrum} or \tt{do} \textit{coupling} command.

\item
    \tt{Print::print_mode}, with overloads which print all, or
    a specified small oscillation eigenvector stored in
    \tt{numerics.modes} and previously evaluated by an
    \tt{evaluate spectrum} or \tt{do} \textit{coupling} command.

\item
    \tt{Print::print_geodesic}, with overloads which will print all,
    or a selected portion (specified by one or two indices),
    of the symbolic geodesic equations stored in
    \tt{global.data().geos};

\item
    \tt{Print::print_cache}, which prints all entries in \tt{Canon::cache};

\item
    \tt{Print::print_sysindex}, which prints a table of \tt{RecIndx}'s
    in \tt{global.sysindex};

\item
    \tt{Print::print_vevindex}, which prints a list of
    coupling constant values and expectation value sets
    contained in a currently open vev data file;

\item
    \tt{Print::print_theory}, which prints the class of theory
    for the running compilation,
    along with lattice size if any dimensions are compactified;

\item
    \tt{Print::print_version}, which prints the version number
    of the running program.

\item
    \tt{Print::print_rkmethods}, which prints the names of known
    Runge-Kutta integration methods in \tt{REdef::list},
    and reports the currently selected method;

\item
    \tt{Print::print_stats}, which prints a variety of statistics
    including the number of defined observables, operators,
    ODE integrations and minimization steps,
    the number of commutator evaluations,
    ODE integrations and minimization steps,
    canonicalizations, and symmetry transformations performed,
    canonicalization cache performance, etc.;

\item
    \tt{Print::print_state}, which prints various global control variables,
    including theory class, active representation and minimization stage,
    maximum observable creation and strong-coupling orders, etc.,
    as well as current verbosity levels.

\item
    \tt{Print::print_blab}, which prints the list of verbosity levels,
    \tt{Blab::blablevel};

\item
    \tt{Print::print_bcktlist}, which prints the contents of
    \tt{global.info().bckt},
    the observable bucket list for the current stage;

\item
    \tt{Print::print_obsstats}, which prints a table of counts of defined
    observables classified according to type, creation order, and expectation order;

\item
    \tt{Print::print_geostats}, which prints the number of
    non-vanishing geodesic equations,
    the number of terms in those equations,
    and the breakdown by monomial order of those terms.

\end{itemize}

\ttsubsec {Rep.\{h,cpp\}}\label{sec:Rep}

Header file \tt{Rep.h} includes
\tt{Symm.h}
and defines:
\begin{itemize}
\item
    \tt{class Proj}, which inherits from \tt{vector<doub>}.
    Every \tt{Proj} object defines a projection operator for
    an irreducible symmetry group representation, with
    the base \tt{vector<double>} providing the list of coefficients
    with which each symmetry transformation in
    \tt{Symm::list} is weighted and then summed to
    form the relevant projector.
    These projection operators are related to the
    representation matrices via
    \begin{equation}
	P^\lambda_{mn} = 
	\frac {d_\lambda}{|G|}
	\sum_{g \in G} \> D^\lambda(g)^*_{mn} \, g \,,
    \end{equation}
    where $G$ denotes the symmetry group
    (with dimension $|G|$),
    $\lambda$ labels an irreducible representation
    (with dimension $d_\lambda$),
    and $D^\lambda(g) \equiv \| D^\lambda(g)_{mn}\|$
    defines the associated matrix representation.

    The \tt{Proj} class adds data members:
    \begin{itemize}
    \item
	\tt{string name}, holding the
	representation name.
	Conventional names for
	crystalography point groups are used,
	augmented by a \tt{p} or \tt{m} denoting
	a representation even or odd, respectively,
	under charge conjugation.
	For multi-dimensional representations
	two further characters,
	taken either from the sequence \tt{abcd} or
	from \tt{xyzw}, are appended to the
	name to indicate the projector indices;
    \item
	\tt{doub denom},
	a denominator normalizing the
	sum of symmetry transforms;
    \end{itemize}

    Member function declarations are:
    \begin{itemize}
    \item
	\tt{operator()} for composing projectors;
    \item
	\tt{operator+},
	\tt{operator-},
	\tt{operator+=} and
	\tt{operator-=}, used in \tt{Rep::repinit}
	for constructing representation projectors;
    \item
	\tt{allzero} for testing complete vanishing
	(used for testing orthonormality),
    \item
	\tt{C_even}, signaling a charge-conjugation
	even representation;
    \item
	\tt{indices} which, for multi-dimensional irreps,
	returns the conversion of the two final characters
	in the projector name into their corresponding numerical
	index values;
    \item
	\tt{rowname} which, for multi-dimensional irreps,
	returns the representation row name by stripping off
	the last character from the projector name;
    \item
	\tt{repname} which, for multi-dimensional irreps,
	returns the representation name by stripping off
	the last two characters from the projector name.
    \end{itemize}
    The static data member
    \tt{vector<Proj> list} holds the list
    of all defined projectors.

    A \tt{friend ostream\& operator<<} declaration
    is included for printing \tt{Proj} objects.
\item
    \tt{class Rep}, which inherits from
    \tt{vector<Proj>}
    and holds the projectors which correspond to a
    given row of an irreducible symmetry representation.

    The \tt{Rep} class adds data members:
    \begin{itemize}
    \item
	\tt{string name}, holding the representation row name.
	For multi-dimensional irreps this is the same
	as the names of the projectors held by the \tt{Rep}
	with their final character stripped off;
    \item
	\tt{bool C_even}, which indicates whether the
	reprsentation is $C$ even or odd.
    \end{itemize}

    The static data member
    \tt{Index<Rep> list} holds the indexed list
    of all defined projectors.
    The fully symmetric vacuum representation is always first in the list.

    Static functions in class \tt{Rep} are:
    \begin{itemize}
    \item
	\tt{known}, which maps representation names to the
	index of the associated \tt{Rep} instance in \tt{Rep::list}
	(and throws an exception if not found);
    \item
	\tt{repinit}, which initializes \tt{Rep::list}.
    \end{itemize}
    A \tt{friend ostream\& operator<<} declaration
    is included for printing \tt{Rep} objects.

\end{itemize}

File \tt{Rep.cpp} includes \tt{Rep.h} and \tt{Gripe.h},
and contains definitions of the above
\tt{Proj} and \tt{Rep} member functions,
and the \tt{ostream\& operator<<} definitions for
both \tt{Proj} and \tt{Rep} objects.

Within \tt{Rep::repinit}, the various \tt{constexpr}
functions defined in \tt{Theory.h}
for distinguishing different lattice types
(like \tt{theory.isR2()} and \tt{theory.isR3()})
are used to select pre-programmed definitions of 
projections onto individual components of irreducible representations
of the lattice symmetry point group times charge conjugation
for theories living on $\mathbb R$, $\mathbb R^2$, $\mathbb R^3$,
or $\mathbb R^2 \times S^1$.
For $\mathbb R^4$, only the
charge conjugation even and odd representations
which are fully symmetric under all lattice
permutations and reflections are currently coded.
The various representation labels begin with
standard cubic crystallography representation names
(those of $C_{4v}$ for theories on $\mathbb R^2$,
$O_h$ for theories on $\mathbb R^3$, and
$D_{4h}$ for theories on $\mathbb R^2 \times S^1$)
followed by a \tt{p} or \tt{m} indicating even
or odd behavior, respectively, under charge conjugation
followed, for multi-dimensional representations, by two
further letters drawn from either \tt{abcd} or \tt{xyzw}
indicating the indices of a particular projector
within the representation.
%
Thanks to the operator overloads for multiplying and
adding representation projectors, the resulting definitions of
various representation projections are fairly concise.
As an example,
\tt{(s("E")\,-\,s("C"))} \tt{*}
\tt{(s("E")\,+\,s("Rx"))} \tt{*}
\tt{(s("E")\,+\,s("Ry"))} \tt{*}
\tt{(s("E")\,-\,s("Pxy"))}
serves to enter the definition of the two-dimensional
\tt{B1m} representation,
which is odd under charge conjugation and $x\leftrightarrow y$
permutations and even under $x\leftrightarrow -x$ and
$y\leftrightarrow -y$ reflections.
(Here, \tt{s} is a pointer to the \tt{Rep::known} function.)

\ttsubsec {RK.\{h,cpp\}}\label{sec:RK}

Header file \tt{RK.h} includes
\tt{Gordion.h}
and defines:
\begin{itemize}
\item
\tt{class RKdef}, whose instances define
an embedded explicit Runge-Kutta integration
method.%
\footnote
    {%
    ``Embedded'' RK methods refer to Runge-Kutta
    schemes in which a common set of derivative evaluations
    is used to derive two different approximations to the
    function being integrated, with the difference between
    the approximations providing an error estimate which
    may be used for adaptive step size control.
    }
Data members of the class are:
\begin{itemize}
\item
    \tt{const int nstage}, giving the number of stages
    (\# derivative evals per step);
\item
    \tt{const int order}, giving the order of the
    higher-order approximation;
\item
    \tt{const double* c}, pointing to the 
    list of RK nodes;
\item
    \tt{const double* const *a}, pointing to the
    matrix of RK coefficients;
\item
    \tt{const double* b}, pointing to the list
    of (higher order) RK weights;
\item
    \tt{const double* bh}, pointing to the list
    of (lower order) RK weights;
\item
    \tt{const double* db}, pointing to the list
    of error vectors weights, equaling $bh - b$;
\item
    \tt{char* name}, pointing to the method name;
\item
    \tt{bool fsal}, indicating if the scheme
    is ``First Same as Last (FSAL),'' meaning that the
    last derivative evaluation in a step may
    be reused as the first derivative evaluation
    needed in the next step.
\item
    \tt{static const std::array<RKdef,...> list}, holding
    the complete set of defined Runge-Kutta methods.
\end{itemize}
\end{itemize}

File \tt{RK.cpp} includes \tt{RK.h} and provides
the definitions of particular Runge-Kutta schemes,
all of which are coded as \tt{constexpr} definitions.
The currently defined RK methods are the following
(see the indicated references for details):
\begin{itemize}\advance\itemsep -4pt
\item
    \tt{bs32:~Bogacki/Shampine 3(2)} 
    \cite{BS32};
\item
    \tt{ss32:~Sofroniou/Spaletta 3(2)}
    \cite{SS32};
\item
    \tt{ss43:~Sofroniou/Spaletta 4(3)}
    \cite{SS32};
\item
    \tt{dp54:~Dormand/ Prince 5(4)}
    \cite{DP54};
\item
    \tt{bs54:~Bogacki/Shampine 5(4)}
    \cite{BS54};
\item
    \tt{jb54:~John Butcher 5(4)}
    \cite{JB54};
\item
    \tt{jv65:~Jim Verner 6(5)}
    \cite{JV65}.
\end{itemize}

\ttsubsec {Save.\{h,cpp\}}\label{sec:Save}

Header file \tt{Save.h} includes
\tt{Gordion.h},
\tt{Coupling.h}
\tt{Data.h} and
\tt{Version.h},
and defines
\begin{itemize}
\item
    \tt{namespace Save} which 
    declares the various functions described next
    for writing data to disk or reloading saved data back from disk,
    and defines:
    \begin{itemize}
    \item
	\tt{struct FileHdr} whose instances will constitute a 16-byte 
	header written to the beginning of sys info and vev data files.
	The \tt{FileHdr} struct contains:
	containing:
	\begin{itemize}
	\item
	\tt{char8 name}, for holding a theory name;
	\item
	\tt{Version version}, for holding the program version;
	\item
	\tt{uint ncoup}, for holding the number of adjustable coupling constants;
	\item
	\tt{uint nvev}, for holding the number of observable expectation values,
	\end{itemize}
	and defines the convenience functions:
	\begin{itemize}
	\item
	    \tt{vevsize}, returning \tt{nvev * sizeof (doub)}, the byte-length
	    of each list of expectation values;
	\item
	    \tt{coupsize}, returning \tt{ncoup * sizeof (Coupling)}, the byte-length
	    of the adjustable coupling list;
	\item
	    \tt{cvsetsize}, returning \tt{vevsize() + coupsize()};
	\item
	    \tt{is_sysfile}, returning \tt{true} if \tt{ncoup} and \tt{nvev} vanish;
	\item
	    \tt{is_vevfile}, returning \tt{true} if \tt{ncoup} or \tt{nvev} are
	    non-vanishing.
	\end{itemize}
    \end{itemize}
    The namespace \tt{Save} also defines the static variables:
    \begin{itemize}
    \item
	\tt{FileHdr filehdr}, holding the file header of an open save file and
	used by the various \tt{Save} commands;
    \item
	\tt{string syspath}, holding the full pathname of an open sys info file;
    \item
	\tt{std::mutex savemutex}, used to synchronize writing of a sys info file
	by multiple threads.
    \end{itemize}
\end{itemize}

File \tt{Save.cpp}
includes the header files
\tt{Save.h},
\tt{Canon.h},
\tt{Gen.h},
\tt{Numerics.h},
\tt{Rep.h},
\tt{Blab.h} and
\tt{Gripe.h}, and defines:
\begin{itemize}
\item
    \tt{static constexpr int rhdrsiz}, as shorthand
    for \tt{sizeof (RecHdr)};
\item
    \tt{static\,symb\,symbbuf[rhdrsiz]}, a \tt{RecHdr}-sized
    \tt{symb} array defined with \tt{alignas (rhdrsiz)} and
    used by \tt{Save::write_obs} and \tt{Save::write_op}
    to transfer trailing portions of \tt{Obs} or \tt{Op}
    symbol strings to a \tt{DataRec} record;

\item
    \tt{Save::save_sys}, which takes as argument a sys info file name,
    strips off any directory portion (redefining \tt{global.savedir}),
    attempts to open the file for writing in binary mode
    and, on success, calls the following \tt{Save::write_*} routines
    starting with \tt{write_header} and ending with \tt{write_sysindex}.
    If the theory contains fermions, separate sets of the
    intervening data records
    are written for each minimization stage;

\item
    \tt{Save::save_vev}, which takes as argument a vev data file name,
    strips off any directory portion (redefining \tt{global.savedir}),
    attempts to open the file for I/O in binary append mode
    and, on success, calls
    \tt{read_header} for a pre-existing file to validate the
    file header,
    or \tt{write_header} for a newly created vev data file,
    and then calls \tt{Save::write_coup} and \tt{write_vev} to append
    a coupling+vev data set to the file.

\item
    \tt{Save::write_header},
    which writes a \tt{FileHdr} instance containing 
    \tt{theory.name}, \tt{global.version}, and supplied values of
    the number of coupling constants and number of vev's
    to the open sys info or vev data file;

\item
    \tt{Save::write_op}, which steps through the \tt{Op}'s of 
    \tt{global.info(stage).ops}, copying their metadata
    to a \tt{OpHdr} and then appending that \tt{OpHdr} followed
    by the \tt{SymbStr} defining the \tt{Op} to
    \tt{global.data().op}.
    The current file position
    and number of stored \tt{Op}'s are recorded in
    \tt{global.data().op.entry()},
    the resulting serialized list of operators is written
    to the sys info file, and then this serialized data record is cleared
    (with its memory forcibly released via the C++ swap trick);

\item
    \tt{Save::write_obs}, which loops over all \tt{Obs}'s
    for the current stage in \tt{ObsList::obs},
    copying their metadata to a \tt{ObsHdr}
    struct and then appending that \tt{ObsHdr} followed by the \tt{SymbStr}
    defining the \tt{Obs} to \tt{global.data().obs}.
    After recording the current file position and
    number of stored observables in
    \tt{global.data().obs.entry()} and
    writing the resulting serialized list of observables 
    to the sys info file, this serialized data record is cleared
    (with its memory forcibly released);

\item
    \tt{Save::write_gen}, which loops over all symmetry representations
    and then over the \tt{Gens}'s 
    in \tt{global.info().gens[rep]}, copying each
    \tt{Gen}'s metadata to a \tt{GenHdr} and then appending
    that \tt{GenHdr} followed by the \tt{OpTerm}'s defining
    the generator to \tt{global.data().gen}.
    After recording the current file position and
    number of stored \tt{Gen}'s in
    \tt{global.data().gen.entry()} and
    writing the resulting serialized list of generators 
    to the sys info file, this serialized data record is freed;

\item
    \tt{Save::write_grad}, which
    writes \tt{global.data().grad} to the sys info file,
    after first recording the current sys info file position in
    \tt{global.data().grad.entry()};

\item
    \tt{Save::write_curv}
    which, for each symmetry representation,
    writes the polynomial scripts in
    \tt{global.data().curv[rep]}
    to the sys info file
    (after first recording
    the current file position in
    \tt{global.data().curv[rep].entry()});

\item
    \tt{Save::write_lagr}
    which, for each symmetry representation,
    writes the polynomial scripts in
    \tt{global.data().lagr[rep]},
    to the sys info file
    (after first recording
    the current file position in
    \tt{global.data().lagr[rep].entry()});

\item
    \tt{Save::write_geos}
    which, for each observable bucket,
    writes the polynomial scripts in \tt{global.data().geos[bckt]}
    to the sys info file (after first recording the current file
    position in \tt{global.data().geos[bckt].entry()});

\item
    \tt{Save::write_geo_bckt},
    which acquires a mutex lock on the sys info file
    and then writes
    \tt{global.data().geos[bckt]},
    for a specified bucket of observables,
    to the sys info file, overwriting the
    initial sys info file index before appending
    an updated file index and releasing the lock.
    If \tt{global.geoswap} is \tt{true},
    the geodesic bucket data record is then freed;

\item
    \tt{Save::write_stat}, which updates
    \tt{global.data().stat.entry()}, and then writes
    the statistics counters in
    \tt{global.info().count} to the sys info file;

\item
    \tt{Save::rewrite_stat}, which seeks to the
    file position of the \tt{Stat} record specified in
    \tt{global.data().stat.entry()}, and then writes
    to that file location in the sys info file
    the current values of the statistics counters in
    \tt{global.info().count};

\item
    \tt{Save::write_sysindex}, which appends
    the contents of \tt{global.sysindex}
    to the sys info file;

\item
    \tt{Save::write_coup}, which 
    appends the contents of \tt{Coupling::list} to the
    vev data file;

\item
    \tt{Save::write_vevs}, which 
    appends the contents of \tt{numerics.vev} to the vev data file;

\item
    \tt{Save::load_save}, which takes as arguments a save file name
    and data set index (default $-1$),
    strips off any directory portion in the supplied file name
    (redefining \tt{global.savedir}),
    attempts to open the file for reading in binary mode
    and on success calls \tt{read_header} to determine
    if the supplied file is a sys info or vev data file.
    For sys info files, the following \tt{Save::read_*} routines
    are called starting with \tt{read_sysindex}
    and ending with \tt{read_stat}.
    For vev data files,
    \tt{load_vev} is called 
    with the supplied data set index passed as argument;

\item
    \tt{Save::load_vev}, which takes a data set index
    as argument and calls
    \tt{read_coup} and \tt{read_vev} to load the
    specified data set from an open vev data file.
    A set index of $-1$ indicates the final data set in the file.

\item
    \tt{Save::read_header}, which reads the
    \tt{FileHdr} stored at the beginning of
    sys info and vev data files,
    and compares
    the contained program version to \tt{global.version}
    and the contained theory name to \tt{theory.name},
    (raising an exception if either are incompatible with the
    running program),
    and for vev data files validating the values of
    \tt{ncoup} and \tt{nvev} in the file header;

\item
    \tt{Save::read_sysindex}, which reads \tt{sizeof global.sysindex} bytes from
    the end of a sys info file into \tt{global.sysindex};

\item
    \tt{Save::read_op}, which reads
    into \tt{global.data().op}
    the serialized \tt{Op} data record
    for the specified stage
    from the sys info file (after seeking to the file position recorded in
    the associated \tt{RecIndx} read in by \tt{read_index}),
    then steps through the data record reconstructing each
    \tt{Op} and storing it in \tt{global.info(stage).ops}, and lastly
    frees the data record;

\item
    \tt{Save::read_obs}, which reads
    into \tt{global.data().obs}
    the observable data record
    for the specified stage
    from the sys info file (after seeking to the file position recorded in
    the associated \tt{RecIndx} read in by \tt{Save::read_index}),
    then steps through this record reconstructing each
    \tt{Obs} and storing it in \tt{ObsList::obs}.
    The routine also
    sets \tt{global.info().nobs} and
    \tt{global.info().maxord} based on the reconstructed observables,
    calls \tt{global.mk_bcktlist} to (re)initialize the observable
    bucket list, and frees the data record;

\item
    \tt{Save::read_gen}, which reads
    into \tt{global.data().gen}
    the serialized generator data record
    for the specified stage
    from the sys info file (after seeking to the file position recorded in
    the associated \tt{RecIndx} read in by \tt{read_index}),
    steps through the data record reconstructing each
    \tt{Gen} and storing it in \tt{global.info().gens},
    appropriately updates \tt{global.info().neven}
    and \tt{global.info().maxgen},
    and lastly frees the data record;

\item
    \tt{Save::read_grad},
    which reads into \tt{global.data().grad}
    the polynomial script data record whose file position and
    length is recorded in the
    associated \tt{RecIndx} read in by \tt{read_index};

\item
    \tt{Save::read_curv},
    which loops over all representations,
    seeking to the file position recorded in the
    associated \tt{RecIndx} read in by \tt{read_index}
    and reading the specified polynomial script data
    into \tt{global.data().curv[rep]};

\item
    \tt{Save::read_lagr},
    which loops over all representations,
    seeking to the file position recorded in the
    associated \tt{RecIndx} read in by \tt{read_index}
    and reading the specified polynomial script data
    into \tt{global.data().lagr[rep]};

\item
    \tt{Save::read_geos}
    which, for each bucket of observables,
    seeks to the file position recorded in the
    associated \tt{RecIndx} read in by \tt{read_index}
    and then reads the specified polynomial script data
    into \tt{global.data().geos[bckt]};

\item
    \tt{Save::read_geo_bckt},
    which seeks to the file position of a 
    particular bucket of geodesics and
    then reads
    into the specified \tt{global.data().geos[bckt]}
    the off\-loaded polynomial script data.
    To enable efficient multi-threading, a
    \tt{thread_local} \tt{fstream} with a custom
    stream buffer is employed.
    Each call to \tt{read_geo_bckt}
    reads data into a \tt{thread_local} data array
    which is swapped into the desired
    \tt{global.data().geos} entry.
    Calls to \tt{read_geo_bckt} with a negative bucket number
    are used to signal a reverse swap of the thread-local data array;

\item
    \tt{Save::read_stat}, which seeks to the file position
    in the associated \tt{RecIndx} read in by \tt{read_index}
    and then reads the contents of the previously saved \tt{Counters}
    structure into \tt{global.data().count};

\item
    \tt{Save::read_coup}, which seeks to the file position
    of the specified data set in the vev data file
    (with $-1$ indicating the final data set in the file)
    and then reads the list of \tt{Coupling}'s into either
    \tt{Coupling::list} or supplied temporary coupling list
    (with the latter case used by \tt{Print::print_vevindex});

\item
    \tt{Save::read_vevs}, which seeks to the file position in the
    of the specified data set in the vev data file
    (with $-1$ indicated the final data set in the file)
    and then reads the specified block of expectation values into
    \tt{numerics.vev}.

\end{itemize}


\ttsubsec {Splice.h}\label{sec:Splice}

Header file \tt{Splice.h} defines:
\begin{itemize}\advance\itemsep -2pt
\item
    \tt{struct Splice}, which holds three data members:
    \begin{itemize}
    \item
    \tt{signed char coeff}, which is either $\pm 1$ or 0;
    \item
    \tt{signed char prefix}, which is either a type-and-reflection \tt{symb} code or $-1$;
    \item
    \tt{signed char suffix}, which is either a type-and-reflection \tt{symb} code or $-1$.
    \end{itemize}

\item
    The static
    \tt{Splice splicetbl[10][10][2]} splicing table
    used in \tt{Commute.cpp}.
    The first two indices are the type-and-reflection (\tt{tnR}) codes of
    two \tt{symb}'s, with the last two-valued index distinguishing two
    possible resulting terms needed for some combinations.

    To understand \tt{splicetbl}, consider a commutator
    $[\mathrm{tr}\, (\alpha\, s_1), \mathrm{tr}\, (s_2 \, \beta)]$
    where $s_1$ and $s_2$ are symbols each representing one of:
    \begin{itemize}\advance\itemsep -2pt
    \item
    a forward-directed gauge link (shorthand notation \tt{l}),
    \item
    an inverse gauge link (\tt{L}),
    \item
    an electric field transforming at the start of its link (\tt{E}),
    \item
    an electric field parallel transported to the end of its link (\tt{LEl}),
    \item
    an electric field followed by its forward link (\tt{El}),
    \item
    an electric field preceded by its inverse link (\tt{LE}),
    \item
    a squared electric field transforming at the start of its link (\tt{EE}),
    \item
    a squared electric field parallel transported to the end of its link (\tt{LEEl}),
    \item
    a squared electric field followed by its forward link (\tt{EEl}),
    \item
    a squared electric field preceded by its inverse link (\tt{LEE}),
    \end{itemize}
    with $\alpha$ and $\beta$ denoting the remainder of the \tt{SymbStr}'s
    for each factor.
    Assume that both $s_1$ and $s_2$ reside on the same lattice link.

    The above list shows
    the ten possible pure gauge symbol types associated with
    any given lattice link.
    (Comments using the above shorthands are present
    in the entries of \tt{Splice.h} defining \tt{splicetbl}.)
    The result of evaluating the commutator, using
    the underlying lattice theory commutation relations between
    electric field operators and gauge links, takes the form:
    \begin{align}
	[\mathrm{tr}\, (\alpha\, s_1), \mathrm{tr}\, (s_2 \, \beta)]
	=
	\sum_{\rm i}
	\mathit{sgn}_i \,
	\mathrm {tr}\, ( \mathit{prefix}_i \, \alpha \, \mathit{suffix}_i \, \beta )
	+ \cdots \,,
    \end{align}
    where the ellipsis denotes further terms arising from the commutation
    of other pairs of symbols in the two factors.
    Here $\mathit{sgn}_i$ is a $\pm$ sign,
    each $\mathit{prefix}_i$ is either empty, or some symbol residing on the
    common link of $s_1$ and $s_2$,
    each $\mathit{suffix}_i$ is likewise either empty or some symbol
    on that same link,
    and the number of terms in the result may be 0, 1, 2 or 3
    depending on the types of the two symbols $s_1$ and $s_2$.
    The \tt{Splice} elements of the \tt{splicetbl} provide
    these needed signs, prefixes and suffixes.
    (More precisely, a \tt{Splice} element provides the type-and-reflection
    portion of the prefix and suffix symbols; the axis direction bits are
    set by those of $s_1$ and $s_2$.)
    A zero coefficient indicates that the term is not needed.
    A prefix of $-1$ is used to indicate that no prefix is needed,
    and likewise for a suffix of $-1$.
    The commutation routines perform the above indicated joining of
    symbol strings using the entries of both
    \tt{splicetbl[}$s_1$\tt{][}$s_2$\tt{]}
    and (with a sign flip)
    \tt{splicetbl[}$s_2$\tt{][}$s_1$\tt{]}.
    The \tt{splicetbl} is constructed so that
    non-vanishing entries are only present when the first table index
    does not represent a link or inverse link.
\end{itemize}

\ttsubsec {Symb.\{h,cpp\}}\label{sec:Symb}

Header file \tt{Symb.h} includes
\tt{Gordion.h}
and defines:
\begin{itemize}
\item
    \tt{class SymbStr}, which inherits from \tt{std::string}
    and whose instances are used to hold \tt{symb} strings.
    The \tt{SymbStr} class declares the member functions:
    \begin{itemize}\advance\itemsep -2pt
    \item
    \tt{join},
    which appends a \tt{symb} or \tt{SymbStr} with appropriate canceling
    of links and inverse links, and combining into appropriate
    ligatures links and electric field insertions living on the same
    lattice link;
    \item
    \tt{joinends}, which properly join the ends of a \tt{SymbStr};
    \item
    \tt{excise}, which removes an interior portion of a \tt{SymbStr};
    \item
    \tt{isclosed}, which tests if a \tt{SymbStr} forms a closed loop;
    \item
    \tt{print}, which converts a \tt{SymbStr} to a printable string.
    \end{itemize}
    A static inline boolean member \tt{dots} controls whether
    the print form of a \tt{SymbStr} includes
    periods between symbols.
    A \tt{friend ostream\& operator<<} declaration
    is included for printing \tt{SymbStr} objects.

\item
    \tt{namespace Symb},
    which defines the specific encoding used for symbols.
    This encoding is
    designed to allow the same encoding to be used
    for theories of any dimension (up to 4).
    The lowest two bits of a \tt{symb} indicate
    the link axis for a gauge field \tt{symb},
    or flavor plus staggering for a fermion \tt{symb}.
    The third lowest bit signals reflection (for gauge)
    or conjugation (for fermions).
    The next three bits encode the type of the \tt{symb}:
    \tt{link}, \tt{E}, \tt{Elink}, \tt{EE}, \tt{EElink}, \tt{fermion}.
    The resulting values $0, \ldots, 39$ represent
    the 40 possible pure gauge symbols in up to four dimensions,
    while values $40, \ldots, 47$ represent up to two flavors
    of fermion fields and their conjugates, with or without
    an included lattice staggering factor.
    A few further symbol values are used for special purposes,
    namely to represent either the $O(N^2)$ gauge or
    $O(N)$ fermion contribution to the entropy in Euclidean theories,
    or in Hamiltonian theories to distinguish the $O(N)$ fermion
    contribution to the expectation value of gauge field kinetic energy terms.

    The \tt{Symb} namespace also defines various useful constants and flags:
    \begin{itemize}\advance\itemsep -2pt
    \item
	\tt{symb Rflag (0x04)},
	the reflection (or conjugation, for fermions) bit mask;
    \item
	\tt{symb addE (0x10)},
	the change in \tt{symb} codes for converting
	a \tt{Link} $\to$ \tt{Elink};
    \item
	\tt{symb YMend (0x28)}, the number of pure gauge symbols;
    \item
	\tt{symb Nsymb (0x30)}, the number of pure gauge plus fermion symbols;
    \item
	\tt{symb TandR (0x3c)}, the type and reflection bit mask;
    \item
	\tt{symb EntG (0x34)}, the extra symbol for gauge entropy;
    \item
	\tt{symb EntF (0x35)}, the extra symbol for fermion entropy;
    \item
	\tt{symb Null (0x40)}, the null ligature flag;
    \item
	\tt{symb X (-1)}, the invalid ligature or empty prefix/suffix marker;
    \end{itemize}
    plus static data members:
    \begin{itemize}
    \item
    \tt{vector<const string> symbname},
    which serves to map symbol values to print forms of each symbol;
    \item
    \tt{symb ligtable[10][10]}, which gives the type-and-reflection
    codes for combining two \tt{symb}'s referring to the same lattice
    link into the appropriate single \tt{symb}, e.g., for merging
    an electric field insertion together with its corresponding link.
    (See the comments in \tt{Symb.h} for details.)
    \end{itemize}

    Functions in namespace \tt{Symb} are:
    \begin{itemize}
    \item
    \tt{to_symb}, which does the inverse mapping from
    printable character(s) to the corresponding \tt{symb} value
    or to a \tt{SymbStr} object;
    \item
    \tt{in_thy}, which tests whether a given
    \tt{symb} value is valid for the particular theory under
    consideration
    (only used when printing symmetry transformations);
    \item
    \tt{ligature}, which
    uses the static \tt{ligtable} and identifies
    when two adjacent symbols should be merged into a single
    symbol based on their types and lattice directions;
    \end{itemize}
    plus a slew of additional inline convenience functions:
    \begin{itemize}\advance\itemsep -2pt
    \item
    \tt{islink}, which tests if a symbol is type \tt{link};
    \item
    \tt{nonlink}, which tests if a symbol is not type \tt{link};
    \item
    \tt{isE}, which tests if a symbol is type \tt{E};
    \item
    \tt{isEE}, which tests if a symbol is type \tt{EE};
    \item
    \tt{isElink}, which tests if a symbol is type \tt{Elink};
    \item
    \tt{isEElink}, which tests if a symbol is type \tt{EElink};
    \item
    \tt{EorElink}, which tests if a symbol is type \tt{E} or \tt{Elink};
    \item
    \tt{EEorEElink}, which tests if a symbol is type \tt{EE} or \tt{EElink};
    \item
    \tt{inclE}, which tests if a symbol is any of the \tt{E}-containing types;
    \item
    \tt{isEE_F}, which tests if a symbol is a fermion portion of an \tt{EE};
    \item
    \tt{isentpy}, which tests if a symbol represents entropy;
    \item
    \tt{isferm}, which tests if a symbol represents a fermion;
    \item
    \tt{isxtra}, which tests if a symbol is one of the extras (not in general observables);
    \item
    \tt{nostep}, which tests if a symbol does not represent a step between lattice sites;
    \item
    \tt{isrefl}, which extracts the reflection bit of a symbol;
    \item
    \tt{isstag}, which extracts the fermion staggering bit;
    \item
    \tt{flav}, which extracts the bit representing fermion flavor;
    \item
    \tt{axis}, which extracts the two bits defining the lattice axis of a gauge symbol;
    \item
    \tt{direction}, which extracts the three bits defining lattice axis and link direction;
    \item
    \tt{type}, which returns the type of a \tt{symb};
    \item
    \tt{tnR}, which returns the type plus reflection portion of a \tt{symb};
    \item
    \tt{refl}, which returns the reflection of a \tt{symb};
    \item
    \tt{conj}, which returns the conjugation of a \tt{symb};
    \item
    \tt{stag}, which returns the fermion \tt{symb} with opposite staggering factor;
    \item
    \tt{step}, which returns $+1$, $-1$, or $0$ depending of whether a \tt{symb} represents
	       a forward step, backward step, or no step.
    \end{itemize}

\end{itemize}

File \tt{Symb.cpp} includes
\tt{Symb.h},
\tt{Theory.h} and
\tt{Gripe.h},
and contains the definitions of the \tt{namespace::Symb} member functions
\tt{in_thy},
\tt{to_symb},
\tt{ligature},
the class \tt{SymbStr} member functions
\tt{join},
\tt{excise},
\tt{isclosed} and
\tt{print},
plus \tt{ostream\& operator<<} for \tt{SymbStr} objects.

\ttsubsec {Symm.\{h,cpp\}}\label{sec:Symm}

Header file \tt{Symm.h} includes
\tt{Op.h}
and defines:
\begin{itemize}
\item
    \tt{class SymmTerm}, which inherits from
    \tt{Term<uint>} and adds the member functions
    \tt{operator/ (double)} for rescaling the
    term coefficient, and
    \tt{operator+} and \tt{operator-}
    for adding or subtracting \tt{SymbTerm}'s,
    producing a \tt{SymmSum}.
    Each \tt{SymmTerm} instance holds the
    index of a symmetry (in \tt{Symm::list})
    plus a numerical coefficient;

\item
    type alias \tt{SymmVec = vector<SymmTerm>},
    for sums of \tt{SymmTerm}'s;

\item
    \tt{class SymmSum}, which inherits from
    \tt{SymmVec} and adds
    \tt{operator+},
    \tt{operator-} and
    \tt{operator*} member functions
    for adding, subtracting
    and multiplying \tt{SymmSum}'s;

\item
    \tt{struct Symmhash}, which declares
    \tt{operator()} for computing a \tt{Symm} hash function;

\item
    type alias \tt{Symmmap = hash<Symm,int,Symmhash>},
    for mapping symmetry transformations to their
    index in \tt{Symm::list};

\item
    \tt{class Symm}, which holds information defining a
    symmetry transformation.
    Data members are:
    \begin{itemize}
    \item
    \tt{array<symb,Nsymb> map}, which maps symbol
    values (representing a \tt{link}, \tt{Elink}, \tt{Fermion}, etc.)
    to the appropriately transformed symbols;
    \item
    \tt{bitset<Nsymb+1> sgn}, which records which symbols need to
    be accompanied by a sign flip, plus one more bit signaling
    transformations which include charge conjugation;
    \item
    \tt{string name}, which gives a standard name to the symmetry transformation;
    \item
    \tt{int indx}, which gives the index of the transformation
    in the \tt{Symm::list} of all symmetry transformations.
    \end{itemize}

    Member function declarations are:
    \begin{itemize}\advance\itemsep -2pt
    \item
    \tt{operator()}, with two overloads
    for applying the symmetry transformation to
    either a supplied \tt{Op} or another \tt{Symm}
    symmetry transform;
    \item
    \tt{operator==}, for comparison with another \tt{Symm};
    \item
    \tt{operator+} \& \tt{operator-}, for building
    linear combinations of symmetry transformations;
    \item
    \tt{isCodd}, returning a boolean indicating if the
    transformation includes charge conjugation,
    which reverses the order of symbols;
    \item
    \tt{print}, which produces a printable string showing the \tt{Symm}
    transformation.
    \end{itemize}
    A friend \tt{ostream\& operator<<} declaration for \tt{Symm} objects
    is also included.

    Static data members in \tt{class Symm} are:
    \begin{itemize}\advance\itemsep -2pt
    \item
	\tt{array<symb,Nsymb> idmap}, the identity transformation, used in the default
	\tt{Symm} constructor;
    \item
	\tt{Symmmap trans2indx}, for mapping an explicit symmetry
	transformation to its corresponding list index;
    \item
	\tt{Index<Symm> list}, an indexed list of all symmetries;
    \end{itemize}
    while static function declarations are:
    \begin{itemize}
    \item
    \tt{symminit}, for initialization of \tt{Symm::list} and \tt{Symm:trans2indx};
    \item
    \tt{known}, which returns a \tt{SymmTerm}
    corresponding to a known symmetry name
    (or throws an exception for unknown names).
    \end{itemize}

\end{itemize}

File \tt{Symm.cpp} includes \tt{Symm.h}
and contains the definitions of the above
\tt{Symm} and \tt{SymmSum} member functions, the
\tt{Symmhash::operator()} hashing function, and
the \tt{operator<<} for \tt{Symm} objects.

The \tt{Symm::symminit} initialization function
loops over whether or not to include
charge conjugation,
then over all permutations of lattice directions
which leave unchanged the overall lattice volume
(if compactified),
and then over all possible reflections of lattice directions.
Inside these loops, a cyclic decomposition of the permutation
is performed (to construct a nice naming of symmetries)
and the specified \tt{map} and bitset \tt{sgn} members
of a \tt{Symm} are initialized, with the result then entered
in \tt{Symm::list} and \tt{Symm::trans2indx}.

\ttsubsec {Term.h}\label{sec:Term}

Header file \tt{Term.h} includes
\tt{Print.h}
and defines:
\begin{itemize}
\item
    \tt{class Term<T>}, whose instances contain two data members,
    a \tt{doub coeff} and a type \tt{T item}.
    Vectors of these templated class objects are used to
    represent linear combinations of objects of type \tt{T}.
    (The \tt{Term<T>} class does not simply inherit from a type \tt{T}
    base class because \tt{T}  may be an elementary type,
    not a class, and hence cannot serve as a base class.)
    A \tt{friend std::ostream\& operator<<} is defined for
    printing \tt{Term} objects.

    Instances of this template class are used to store
    \tt{Polyindx} monomials in \tt{ObsPoly}'s,
    \tt{OpTerm} monomials in \tt{OpSum}'s and
    \tt{Symm} monomials in \tt{SymmSum}'s.
\end{itemize}

\ttsubsec  {Test.\{h,cpp\}}\label{sec:Test}

Header file \tt{Test.h} includes
\tt{Gordion.h}
and defines
\begin{itemize}
\item
    \tt{namespace Test}, containing the functions:
    \begin{itemize}
    \item
	\tt{Test::irreps}, which tests if the
	defined representation projectors
	satisfy the proper orthonormality relation,
	\begin{equation}
	    P^\lambda_{mn} \, P^\mu_{m'n'}
	    =
	    \delta^{\lambda\mu} \, \delta_{nm'} \, P^\lambda_{mn'} \,,
	\label{eq:orthonorm}
	\end{equation}
	and the completeness condition,
	\begin{equation}
	    \sum_{\lambda} \sum_{m=1}^{d_\lambda} P^\lambda_{mm} = e \,,
	\end{equation}
	(where $e$ is the identity group element);
    \item
	\tt{Test::jacobi},
	which tests the validity of the Jacobi identity.
	One overload evaluates the identity with a specified
	observable and all generators in the currently selected
	representation, that is
	\begin{align}
	    [gen_1, [gen_2, obs]] + [ [gen_2, gen_1], obs]
	    - [gen_2, [gen_1, obs]]
	    \stackrel{?}{=} 0 \,,
	\end{align}
	with the generators and observable specified by their
	indices within the representation's list of generators
	and the \tt{ObsList::obs} observable list.
	A related overload allows one to specify
	two individual \tt{Op}'s by hand along with
	an observable's index, while a final
	overload tests the identity
	\begin{align}
	    [gen_1, [gen_2, gen_3]]
	    + [gen_2, [gen_3, gen_1]]
	    + [gen_3, [gen_1, gen_2]]
	    \stackrel{?}{=} 0 \,,
	\end{align}
	for all triples of generators for the current minimization stage
	and symmetry representation.
	These routines provide consistency checks on the coding
	of the commutation routines.
    \end{itemize}
\end{itemize}

File \tt{Test.cpp} includes
\tt{Test.h},
\tt{Commute.h},
\tt{Global.h},
\tt{Rep.h} and
\tt{Gripe.h},
and contains the definitions of the above functions.

\ttsubsec {Theory.\{h,cpp\}}\label{sec:Theory}

Header file \tt{Theory.h} includes
\tt{Gordion.h}
and defines:
\begin{itemize}
\item
    \tt{enum LATTTYPE \{ Hamilton, Euclid\}}, used just in the
    initialization of the pre-defined class \tt{Theory} instances
    described below;
\item
    \tt{class Coord}, with data member:
    \begin{itemize}
    \item
	\tt{union \{ short comp[4] ; ulong point ; \}},
	which is used for encoding the overall lattice dimensions
	(with 0 indicating an infinite dimension), or
	storing lattice site coordinates or displacements
	(needed for finding loop self-intersections).
	The unnamed union allows accessing either
	individual coordinate components or the
	entire set of coordinates
	as an \tt{unsigned long}, enabling convenient
	sorting or equality testing of \tt{Coord} instances.
    \end{itemize}
    Additional member functions are:
    \begin{itemize}
    \item
	\tt{modlen}, which for periodic lattice dimensions
	maps a particular coordinate to lie within the
	defined period;
    \item
	\tt{modulo}, which for periodic lattices, maps
	all coordinates to lie within the box;
    \item
	\tt{isclosed}, which test for a vanishing
	lattice displacement (after first calling
	\tt{modulo} for periodic lattices).
    \end{itemize}

\item
    \tt{class Theory}, whose single \tt{constexpr} instance
    \tt{theory} holds the basic information defining the type
    of theory.
    Data members are:
    \begin{itemize}\advance\itemsep -3pt
    \item
	\tt{char8 name}, a short theory name which is
	used as part of the default names for save files
	and numeric output files;
    \item
	\tt{char dim}, the lattice dimension;
    \item
	\tt{char nf}, the number of fermion flavors;
    \item
	\tt{char nrep}, the sum of dimensions of all irreducible
	symmetry representations;
    \item
	\tt{bool euclid}, indicating whether theory type is
	Euclidean (default \tt{false}) or Hamiltonian;
    \item
	\tt{Coord box}, the lattice size.
    \end{itemize}
    Various member functions are defined for convenience:
    \begin{itemize}\advance\itemsep -2pt
    \item
	\tt{unbounded}, signaling an unbounded lattice
	in all dimensions;
    \item
	\tt{oddboxsize}, indicating that one or more
	lattice dimensions are compactified with an odd length;
    \item
	\tt{box_size}, which returns a string with a
	printable form of the lattice size;
    \item
	\tt{parent}, which returns a \tt{char8} containing,
	for theories with fermions, the name of the corresponding
	Yang-Mills theory without fermions;
    \item
	\tt{isR1},
	\tt{isR2},
	\tt{isR3}, and
	\tt{isR4}, indicating that the lattice is unbounded
	and of the indicated dimension;
    \item
	\tt{isS1},
	\tt{isR1xS1}, and
	\tt{isR2xS1}, indicating that the lattice has
	one compactified dimension and the indicated number
	of uncompactified dimensions.
    \end{itemize}

    Static member functions are
    \tt{theoryinit}, which validates the theory definition settings,
    and \tt{theorydefn}, which defines the specific
    theory of interest.

    A series of \tt{\#if ... \#elif ... \#elif ... \#else}
    tests for the preprocessor symbol definitions matching
    a pre-defined theory name and appropriately
    defines the \tt{constexpr theory}.
    Currently, available selections are:
    \begin{itemize}\advance\itemsep -1pt
    \item
	\tt{YM1h}: for Hamiltonian Yang-Mills on a 1D lattice compactified to a single link;
    \item
	\tt{YM2h}: for Hamiltonian Yang-Mills on a 2D infinite lattice,
    \item
	\tt{YM3h}: for Hamiltonian Yang-Mills on a 3D infinite lattice,
    \item
	\tt{YM1e}: for Euclidean Yang-Mills on a 1D lattice compactified to a single link;
    \item
	\tt{YM2e}: for Euclidean Yang-Mills on a 2D infinite lattice,
    \item
	\tt{YM3e}: for Euclidean Yang-Mills on a 3D infinite lattice,
    \item
	\tt{YM4e}: for Euclidean Yang-Mills on a 4D infinite lattice,
    \item
	\tt{QCD1h}: for Hamiltonian single flavor QCD on a 1D infinite lattice,
    \item
	\tt{QCD2h}: for Hamiltonian single flavor QCD on a 2D infinite lattice,
    \item
	\tt{QCD3h}: for Hamiltonian single flavor QCD on a 3D infinite lattice,
    \item
	\tt{QCD1e}: for Euclidean single flavor QCD on a 1D infinite lattice,
    \item
	\tt{QCD2e}: for Euclidean single flavor QCD on a 2D infinite lattice,
    \item
	\tt{QCD3e}: for Euclidean single flavor QCD on a 3D infinite lattice,
    \item
	\tt{YM3ep}: for Euclidean Yang-Mills on $S_1$ compactified ($\infty \times \infty \times 1$) 3D lattice,
    \item
	\tt{YM3hp}: for Euclidean Yang-Mills on $S_1$ compactified ($\infty \times \infty \times 1$) 3D lattice.
    \end{itemize}
    If none of these preprocessor tests succeed
    then a compilation error message is generated;

\item
    \tt{constexpr int PSIZ = theory.euclid\;?\;2\;:\;4}, which sets the
    size of index tuples used to represent products of indexed observables,
    with higher order tuples needed for Hamiltonian theories.

\item
    \tt{constexpr int NREP = theory.nrep}, used for dimensioning
    multiple arrays defined in \tt{Global.h}.
\end{itemize}
This header files also includes a few static assertions
validating the theory definitions.
%
File \tt{Theory.cpp} includes
\tt{Theory.h},
\tt{Print.h} and
\tt{Obs.h},
and defines:
\begin{itemize}
\item
    Preprocessor symbol \tt{MASSSCALE} as 0, if not already defined.
    Setting \tt{MASSSCALE} to some non-zero numeric value in the \tt{Makefile}
    will scale the fermion mass with a specified power of the gauge coupling;
    see Eq.~\ref{eq:fermidef}.
\item
    \tt{Theory::theoryinit}, which stores the theory name
    in the default directories for save files and output data,
    and prints the selected theory name and program version;y
    and then prints the theory selection plus the program version.
\item
    \tt{Theory::theorydefn}, which
    makes standard definitions of a Kogut-Susskind
    lattice Hamiltonian or Wilson Euclidean free-energy.
    Specifically, these definitions are:
    \begin{subequations}
    \begin{align}
	    H_{\rm gauge}/N &= \lambda \tr E^2 + \lambda^{-1} \tr (2 - U - U^\dagger) \,,
	\\
	    F_{\rm gauge}/N &= \lambda^{-1} \tr (2 - U - U^\dagger) - S \,,
    \end{align}
    \end{subequations}
    for one-dimensional Yang-Mills, either Hamiltonian or Euclidean, respectively,
    compactified to a single link.
    For higher dimensional Yang-Mills, either Hamiltonian or Euclidean, the definitions are:
    \begin{subequations}
    \begin{align}
	    H_{\rm gauge}/N &= \tfrac 14 \lambda \sum_\ell \> \tr (E[\ell]^2)
		+ \lambda^{-1} \sum_p \> \tr (2 - U[\partial p] - U[\partial p]^\dagger) \,,
	\\
	    F_{\rm gauge}/N &= \lambda^{-1} \sum_p \> \tr (2 - U[\partial p] - U[\partial p]^\dagger)
		- S \,,
    \end{align}
    \end{subequations}
    where $\ell$ labels positively oriented lattice links,
    $p$ denotes positively oriented plaquettes,
    $E[\ell]$ is the U($N$) electric field operator on link $\ell$, and
    $U[\partial p]$ denotes the U($N$) gauge connection around plaquette $p$.
    For QCD, the additional
    fermion contribution to the Hamiltonian or free energy is:
    \begin{subequations}\label{eq:fermidef}
    \begin{align}
	H_{\rm fermi} &= \tfrac 12 \sum_{\langle xy \rangle} \>
		    \psi^\dagger_x \, \eta_{\langle xy \rangle} \, U_{\langle xy \rangle} \, \psi_y
		    + m \, \lambda^p \sum_x \> \psi^\dagger_x \, \sigma_x \, \psi_x \,,
	\\
	F_{\rm fermi} &= \tfrac 12 \sum_{\langle xy \rangle} \>
		    \psi^\dagger_x \, \eta_{\langle xy \rangle} \, U_{\langle xy \rangle} \, \psi_y
		    + m \, \lambda^p \sum_x \> \psi^\dagger_x \, \psi_x \,,
    \end{align}
    \end{subequations}
    where $x$ labels lattice sites,
    $\langle xy \rangle$ denotes the link connecting nearest-neighbor sites $x$ and $y$,
    $U_{\langle xy \rangle}$ is the U($N$) gauge connection on link $\langle xy \rangle$,
    and
    $\eta_{\langle xy \rangle}$ is a non-dynamical $U(1)$ fermion flavor connection.
    The fermion mass $m$ may be scaled by a power \tt{p} of gauge coupling,
    which may be set via the preprocessor symbol \tt{MASSSCALE} (default 0).
    The lattice staggering factor
    $\sigma_x$ equals $+1$ on the even sublattice and $-1$ on the odd sublattice.
    For further details see Ref.~\cite{CSVAII} and the discussion 
    in appendix \ref{sec:Symmetries}.

    Both canonicalized and non-canonicalized forms of the
    terms in the Hamiltonian (or free-energy)  are stored in
    \tt {global.info().Hterms}, since non-vacuum representation
    spectrum calculations need the non-canonicalized forms
    but for minimization and the vacuum representation spectrum
    the more compact canonicalized forms may be used.
    Because lattices with compactified dimensions of different
    sizes are supported, the code in \tt{theorydefn}
    does not assume that all lattice
    directions are equivalent.
    After installing these theory definitions, the routine
    initializes \tt{global.info().MMAlimit}
    to equal the number of initial observables of the stage,
    and clears \tt{global.info().MMAlist}.
\end{itemize}

\ttsubsec  {Version.h}\label{sec:Version}

Header file \tt{Version.h} includes \tt{Gordion.h} 
begins with a large comment block summarizing changes made with
different versions.  It then defines:
\begin{itemize}
\item
    \tt{class Version}, used for program version numbers which are
    stored in sys-info and vev-data files.
    The class defines one non-static data member:
    \begin{itemize}
    \item
	\tt{ushort version}, recording a program version number,
    \end{itemize}
    and several static data members:
    \begin{itemize}
    \item
	\tt{constexpr uchar majornum},
	the running program's major version number;
    \item
	\tt{constexpr uchar minornum},
	the running program's minor version number;
    \item
	\tt{constexpr bool r_is_f},
	indicating if \tt{sizeof (real) == sizeof (float)}.
    \item
	\tt{constexpr ushort myversion},
	holding the running program's version number which combines
	\tt{majornum}, \tt{minornum} and \tt{r_is_f}.
    \end{itemize}
    Class functions are:
    \begin{itemize}
    \item
	\tt{ushoft rsize}, which return the size in bytes of a \tt{real};
    \item
	\tt{ushort major}, which returns the major number portion of the
	class instance;
    \item
	\tt{ushort minor}, which returns the minor number portion of the
	class instance;
    \item
	\tt{ushort arith}, which returns least significant bit of the class
	instance, which indicates if, in the encoded version,
	\tt{sizeof (real)} equals \tt{sizeof(float)};
    \item
	\tt{newer}, which tests if the given instance version (from a save file)
	is newer than the running program's version number,
    \item
	\tt{incompat}, which tests if the encoded real number size from
	the given instance
	is incompatible with the real number size of the running program.
    \end{itemize}
    A \tt{friend ostream\& operator<<} declaration is included for
    printing version numbers.
\end{itemize}

\appendix

\section {Command summary}\label{sec:commands}

The keyword-driven command line interpreter recognizes the following commands,
which may be terminated by either a semicolon (;) or end of line.
In the following summary,
\textit{italic} indicates user input,
brackets (\tt{[]}) indicate optional arguments,
and a vertical bar (\tt{|}) separates alternative arguments.
A plus mark$^+$ after the command indicates the command is only
applicable in Hamiltonian theories, while a double plus$^{++}$
indicates applicability only for Euclidean theories.
Command words as well as subsequent argument words
(other than explicit file, representation, or symmetry names)
may be highly abbreviated;
any initial matching sequence of one or more characters is sufficient
to match the first consistent target word.
(For example, a single \tt{b} will match ``build'',
while \tt{p obs *} will print all observables.)

\begin{description}\advance\itemsep -1pt
\item [\tt \# {\it anything}]~\\
    A comment, to be ignored.

\item [\tt ?]~\\
    Print a command summary.

\item [\tt add generator {\tt{[}}({\it order}){\tt{]}}
    {\tt{[}}{\it coeff}{\tt{]}} {\it op1}
    {\tt{[[}}{\it coeff}{\tt{]}} {\it op2}{\tt{]}} $\cdots$ ]~\\
    Deprecated form of \tt{generator add} command.

    Construct the specified coherence group generator, 
    project it onto all defined irreducible representations,
    and add the results, when independent and non-vanishing, to the
    set of defined generators for each representation.
    One or more \tt{Op}'s (operators in the coherence algebra)
    may be given,
    each optionally preceded by a numerical coefficient;
    this is interpreted as a weighted sum of the indicated \tt{Op}'s.
    The generator order must either be specified explicitly,
    as an initial parenthesized integer,
    or else the first given \tt{Op}, when canonicalized,
    must be a known observable.
    If multiple \tt{Op}'s are given and (when canonicalized)
    are known observables, they must all have a common creation order.

\item [\tt build all {\it max-order}]~\\
    Build observables (up to the specified maximum order)
    as well as the polynomial scripts for geodesics,
    gradient, curvature and (for Hamiltonian theories)
    the Lagrange bracket.

\item [\tt build curvature \tt{[}{\it repname}\tt{]}] \bgroup\itemsep 0pt \parskip 0pt
\item [\tt build curvature all]~\\
    Build the polynomial scripts
    used for evaluating the curvature of the
    Hamiltonian (or free energy) in the current representation.
    A specified representation becomes the current representation, while
    the final form builds the curvature scripts
    for all representations.
\egroup

\item [\tt build geodesics]~\\
    Build the polynomial scripts
    used for evaluating the geodesic equations.

\item [\tt build gradient]~\\
    Build the polynomial scripts
    used for evaluating the gradient
    of the Hamiltonian (or free energy). 

\item [\tt build lagrange$^+$ \tt{[}{\it repname}\tt{]}] \bgroup\itemsep 0pt \parskip 0pt
\item [\tt build lagrange$^+$ all]~\\
    Build the polynomial scripts
    used for evaluating the 
    Lagrange bracket (for Hamiltonian theories) in
    the current representation.
    A specified representation becomes the current representation, while
    the final form builds the Lagrange bracket scripts
    for all representations.
\egroup

\item [\tt build observables {\it max-order}]~\\
    Build observables.
    All observables
    for the given minimization stage up to (and including)
    the specified maximum strong-coupling order are constructed.

\item [\tt call canon {\it obs}] \bgroup \itemsep 0pt \parskip 0pt
\item [\tt call classify {\it obs}]
\item [\tt call commute {\it op} {\it obs}]
\item [\tt call compose {\it symm1} {\it symm2}]
\item [\tt call do\_inner {\it Eloop-obs}]
\item [\tt call noEbound {\it obs}]
\item [\tt call transform {\it symmetry} {\it obs}]
    Call the indicated routines with the given arguments.
    This functionality is present for debugging purposes.
\egroup

\item [\tt do {\it coupling-name initial-value final-value increment}]~\\
    Evolve the specified coupling constant
    from the given initial value to the final value in steps of the
    given increment, performing Newton minimization at each step
    and writing results to the output data file.

\item [\tt do minimize]~\\
    Perform a sequence of minimization steps,
    terminating when the $L_2$ norm of the predicted displacement from
    the minimum falls below the minimization tolerance (or when the
    limit on the maximum number of minimization steps is reached).

\item [\tt do step]~\\
    Perform one Newton minimization step,
    evaluating the gradient and curvature, solving for the
    predicted minimum, and integrating the geodesic equations
    to move to the predicted minimum.

\item [\tt evaluate curvature \tt{[}$*$\tt{]}]~\\
    Numerically evaluate the $T$-even curvature matrix
    produced by double generator commutators,
    using active $T$-even generators in the currently active
    irreducible representation.
    Without an option \tt{*}, the curvature eigenvalues are printed;
    with a \tt{*} the full curvature matrix is first printed as well.

\item [\tt evaluate delta]~\\
    Evaluate, and print, the predicted location of the Hamiltonian or
    free energy minimum, equal to minus the inverse of the
    (vacuum representation) curvature times the gradient.

\item [\tt evaluate freeenergy$^{++}$]~\\
    Numerically evaluate the free energy and print the result.

\item [\tt evaluate geodesics $*$] \bgroup \itemsep 0pt \parskip 0pt
\item [\tt evaluate geodesics {\it print-limit}]~\\
    Compute the variation of all observables
    (of the current minimization stage) in the direction defined by \tt{delta}.
    Because this output vector can be very long, \tt{evaluate}
    \tt{geodesics} takes an argument specifying a limit on
    the number of observables whose derivatives will be printed,
    or else a \tt{*} requesting printing of the entire list.
\egroup

\item [\tt evaluate geostats]~\\
    (Re)evaluates and prints geodesic equation statistics,
    including the number of polynomial terms and their breakdown
    by monomial order.
    (Only relevant for debugging or verifying sys info file integrity.)

\item [\tt evaluate gradient]~\\
    Numerically evaluate, and print, the gradient of the Hamiltonian
    (or free energy) produced by all defined vacuum representation generators.

\item [\tt evaluate hamiltonian$^+$]~\\
    Numerically evaluate the Hamiltonian and print the result.

\item [\tt evaluate lagrange$^+$]~\\
    Evaluate, and print, the Lagrange bracket matrix,
    produced by expectations of generator commutators,
    for generators in the
    currently active irreducible representation.

\item [\tt evaluate metric \tt{[}$*$\tt{]}]~\\
    Numerically evaluate the $T$-odd curvature matrix
    produced by double generator commutators,
    using active $T$-odd generators in the
    currently active irreducible representation.
    Without an option \tt{*}, the curvature eigenvalues are printed;
    with a \tt{*} the full curvature matrix is first printed as well.

\item [\tt evaluate spectrum$^+$ \tt{[}{\it repname}\tt{]}] \bgroup\itemsep 0pt \parskip 0pt
\item [\tt evaluate spectrum$^+$ all]~\\
    Evaluate, and print, the
    small oscillation frequency (or mass) spectrum
    in the current representation
    derived from the eigenvalues of the generalized eigensystem
    of the curvature and Lagrange bracket.
    A specified representation becomes the current representation, while
    the final form evaluates the spectrum for all symmetry representations.
\egroup

\item [\tt generator add {\tt{[}}({\it order}){\tt{]}}
    {\tt{[}}{\it coeff}{\tt{]}} {\it op1}
    {\tt{[[}}{\it coeff}{\tt{]}} {\it op2}{\tt{]}} $\cdots$ ]~\\
    Construct the specified coherence group generator, 
    project it onto all defined irreducible representations,
    and add the results, when independent and non-vanishing, to the
    set of defined generators for each representation.
    One or more \tt{Op}'s (operators in the coherence algebra)
    may be given,
    each optionally preceded by a numerical coefficient;
    this is interpreted as a weighted sum of the indicated \tt{Op}'s.
    The generator order must either be specified explicitly
    as an initial parenthesized integer,
    or else the first given \tt{Op}, when canonicalized,
    must be a known observable.
    If multiple \tt{Op}'s are given and (when canonicalized)
    are known observables, they must all have a common creation order.

\item [\tt generator activate ({\it order})] \bgroup\itemsep 0pt \parskip 0pt
\item [\tt generator activate {\it{gen1\#}} {\tt{[}}{\it gen2\#}{\tt{]}} $\cdots$ ]~\\
    First form activates all generators of the given order, in all representations.
    Second form activates the numbered generator(s) of the currently active representation.
\egroup

\item [\tt generator suspend ({\it order})] \bgroup\itemsep 0pt \parskip 0pt
\item [\tt generator suspend {\it{gen1\#}} {\tt{[}}{\it gen2\#}{\tt{]}} $\cdots$ ]~\\
    First form suspends all generators of the given order, in all representations.
    Second form suspends the numbered generator(s) of the currently active representation.
\egroup

\item [\tt help]~\\
    Print a command summary.

\item [\tt load {\it sys info-file}] \bgroup \itemsep 0pt \parskip 0pt
\item [\tt load \tt{[}{\it data-set\#}\tt{]} {\it vev data-file}]~\\
    Load information from a previously saved sys info or vev data file.
    The theory type of the running program must match that of the specified file
    or, for QCD theories, match the parent Yang-Mills theory.
    For vev data files, the number of observable expectation values must match
    either the total number of observables in the defined truncation, or 
    the number of pure gauge observables.
    The default data set for vev data files is the final one in the file.
\egroup

\item [\tt load {\it data-set\#}]~\\
    Load the specified data set from a currently open vev data file.

\item [\tt print baseobs]~\\
    Print the list of basic observables (those used in the
    Hamiltonian or free energy).

\item [\tt print blablevels]~\\
    Print the defined verbosity levels.

\item [\tt print cache]~\\
    Print the canonicalization cache.

\item [\tt print couplings]~\\
    Print the defined adjustable coupling constants.

\item [\tt print curvature $*$] \bgroup \itemsep 0pt \parskip 0pt
\item [\tt print curvature {\it term\#} {\tt{[}}{\it gen1\#}
    {\tt{[}}{\it gen2\#}{\tt{]}}{\tt{]}} ]~\\
    Print symbolic curvature expressions.
    The first form prints all curvature components.
    The second form allows one to print the curvature components
    of a specified term in the Hamiltonian or free energy
    (without any overall adjustable coupling constant)
    and, optionally, specify one or both
    indices of the generators in the active symmetry representation
    which act on the indicated term.
\egroup

\item [\tt print freeenergy$^{++}$]~\\
    Print the expression defining the free energy.

\item [\tt print fermiinit]~\\
    Print the mapping from closed fermion bilinears to loops which
    is used for initialization of fermion bilinear expectation values.

\item [\tt print generator] \bgroup \itemsep 0pt \parskip 0pt
\item [\tt print generator $*$]
\item [\tt print generator {\it gen\#}]~\\
    Print generator(s).
    The first form prints an abbreviated form of all generators in the
    current representation.
    The second form prints the full form of all generators in the
    current representation.
    The final form prints the full form of a particular generator.
    For each generator,
    the output displays the current representation, the generator index
    prefixed with \tt{g} for gauge or \tt{f} for fermion generators
    and followed by a $*$ suffix if the generator is $T$-odd,
    and then the actual generator definition.
\egroup

\item [\tt print geodesics $*$] \bgroup \itemsep 0pt \parskip 0pt
\item [\tt print geodesics {\it obs\#} {\tt{[}}{\it gen\#}{\tt{]}}]~\\
    Print symbolic geodesic equations.
    The first form prints all geodesic equations.
    The second form allows one to specify a particular observable
    and, optionally, the ($T$-odd) vacuum representation generator
    which acts on the indicated observable.
\egroup

\item [\tt print geostats]~\\
    Prints geodesic equation statistics,
    including the number of polynomial terms and their breakdown
    by monomial order.

\pagebreak[2]
\item [\tt print gradient $*$] \bgroup \itemsep 0pt \parskip 0pt
\item [\tt print gradient {\it term\#} {\tt{[}}{\it gen\#}{\tt{]}}]~\\
    Print symbolic gradient expressions.
    The first form prints all gradient components.
    The second form allows one to print the gradient components of a specified
    term  in the Hamiltonian or free energy
    (without any overall adjustable coupling constant) and,
    optionally, specify the index of the
    $T$-even vacuum representation generator which acts on the indicated term.
\egroup

\item [\tt print hamiltonian$^+$]~\\
    Print the expression defining the Hamiltonian.

\item [\tt print lagrange $*$] \bgroup \itemsep 0pt \parskip 0pt
\item [\tt print lagrange {\it gen1\#} {\tt{[}}{\it gen2\#}{\tt{]}}]~\\
    Print symbolic Lagrange bracket expressions.
    The first form prints all Lagrange bracket components.
    The second form allows one to specify indices of one or both of the
    generators, in the active symmetry representation,
    whose commutator defines a Lagrange bracket component.
\egroup

\goodbreak

\item [\tt print mode$^+$ *]\bgroup \itemsep 0pt \parskip 0pt
\item [\tt print mode$^+$ {\tt{[}}{\it mode\#}{\tt{]}}]~\\
    Print all, or a specified, small oscillation eigenmode
    from the most recently evaluated spectrum.
\egroup

\item [\tt print observable $*$] \bgroup \itemsep 0pt \parskip 0pt
\item [\tt print observable {\it obs1\#} {\tt{[}}{\it obs2\#}{\tt{]}}]
\item [\tt print observable {\it obs}]
\item [\tt print observable ({\it obstype})]
\item [\tt print observable ({\it corder},{\it xorder})]
\item [\tt print observable ({\it corder},{\it xorder},{\it obstype})]~\\
    Print observables.
    The first form prints the entire list of constructed observables.
    The second form allows one to 
    specify a single observable index, or a pair
    of indices indicating a range of observables.
    The third form allows one to input the printable form of an observable
    and determine if the given observable is present in the observable list.
    Canonicalization is not automatically performed on the input observable.
    The final forms allow one to print all observables of a specified type,
    with specified creation and expectation orders or, in the final form,
    specified orders and type.
    No spaces may be present within these parenthesized specifications.
\egroup

\item [\tt print obsstats]~\\
    Print an observable statistics table showing the break down of the
    number of constructed observables by type, creation order
    and expectation order.

\item [\tt print operator $*$] \bgroup \itemsep 0pt \parskip 0pt
\item [\tt print operator {\it op1\#} {\tt{[}}{\it op2\#}{\tt{]}}]~\\
    Print operators(s).
    The first form prints all defined operators (or ``\tt{Op}'s'')
    whose linear combinations, when symmetry projected, form generators.
    The second form allows on to 
    specify a single operator index, or a pair
    of indices to indicate a range of operators.
\egroup

\item [\tt print primary]~\\
    Print the list of primary \tt{Op}'s (those which cannot be
    produced by commutators of lower order \tt{Op}'s).

\item [\tt print representation] \bgroup \itemsep 0pt \parskip 0pt
\item [\tt print representation $*$]
\item [\tt print representation {\it repname}]~\\
    Print irreducible representation (irrep) projectors.
    The first form lists all defined irrep names,
    while the second form prints the full form of all irreps.
    The final form allows one to specify a particular representation by name.
\egroup

\item [\tt print rkmethods]~\\
    Print the names of available Runge-Kutta ODE integration methods,
    and report the currently selected method.

\item [\tt print spectrum$^+$]~\\
    Prints the most recently evaluated small oscillation spectrum.

\item [\tt print state]~\\
    Print information on the theory class, lattice size,
    and global settings including
    minimization stage, active representation,
    ODE and minimization tolerances, ODE integration method,
    data file name, etc.

\item [\tt print stats]~\\
    Print various operational statistics including the
    numbers of observables built, observable classifications,
    canonicalizations,
    approximations, symmetry transforms, etc., the number of geodesic
    equation terms, ODE integration steps, and canonicalization cache
    performance figures.

\item [\tt print symmetry $*$] \bgroup \itemsep 0pt \parskip 0pt
\item [\tt print symmetry {\it symm-name}]~\\
    Print symmetry transformations.
    The first form prints all symmetry transformations,
    while the second form allows one to specify a particular symmetry
    by name.
\egroup

\item [\tt print symmsets]~\\
    Print the various symmetry subsets used in observable canonicalization.

\item [\tt print sysindex]~\\
    Print a table showing the contents of \tt{global.sysindex}.

\item [\tt print vevindex]~\\
    Print a list of data sets in
    a currently open vev data file.

\item [\tt quit]~\\
    Exit the program.

\item [\tt read {\it filename}]~\\
    Read command input from the specified file.

\item [\tt reset blab]\bgroup \itemsep 0pt \parskip 0pt
\item [\tt reset maxthread]
\item [\tt reset mintol]
\item [\tt reset odetol]
\item [\tt reset representation]
\item [\tt reset rkmethod]
\item [\tt reset stage]
\item [\tt reset svdcutoff]
\item [\tt reset sysfile]
\item [\tt reset vevfile]
\item [\tt reset MMAfile]
\item [\tt reset MMAlist]~\\
    Reset the indicated global variables to their default initial values.
    See the various \tt{set} commands for details.
\egroup

\item [\tt reset numerics]~\\
    Re-initialize observable expectation values
    for the current minimization stage.
    For gauge minimization, this zeros all pure gauge
    observable expectation values (except the identity).
    For fermi minimization, this zeros all fermion
    bilinear expectation values except those for which
    the fermion and conjugate fermion reside on the same
    site.  Such bilinears are initialized to minus half
    the value of the corresponding Wilson loop.

\item [\tt save sys {\tt{[}}{\it filename}{\tt{]}}]~\\
    Write to a binary ``sys'' file information on all defined
    operators, observables, generators, and the polynomial scripts
    defining gradient, curvature, Lagrange bracket and geodesic equations,
    plus build-related statistics counters.
    If no filename is specified,
    then a default file name with extension \tt{.sys}
    is used which indicates the theory class,
    the maximum order of defined generators,
    the maximum strong coupling order of observables,
    and whether non-retained observables have been approximated.
    If the supplied filename is a path specification, then
    the directory portion automatically redefines the
    directory used for both \tt{sys} and \tt{vev} save files.

\item [\tt save vev {\tt{[}}{\it filename}{\tt{]}}]~\\
    Append to a binary ``vev'' data file the expectation
    values of all defined observables and current values
    of all adjustable coupling constants.
    Multiple sets of such ``vev'' data may be present
    in the same \tt{vev} file.
    If no filename is specified,
    then a default file name with extension \tt{.vev}
    is used which indicates the theory class,
    the maximum order of defined generators,
    the maximum strong coupling order of observables
    used in minimization,
    and whether non-retained observables have been approximated.
    If the supplied filename is a path specification, then
    the directory portion automatically redefines the
    directory used for both \tt{sys} and \tt{vev} save files.

\item [\tt set {\it coupling-name coupling-value}]~\\
    Sets the value of the named adjustable coupling constant.

\item [\tt set approx true $\vert$ false]~\\
    Specifies whether non-retained observables are
    be approximated, when possible, by products of
    lower order retained observables (default \tt{true}).

\item [\tt set autoToddgens true $\vert$ false]~\\
    Specifies whether $T$-odd
    \tt{Eloop} generators are to be automatically
    created from each pure loop generator by
    evaluating the commutator of the loop with
    the gauge kinetic energy, and whether fermion generators
    of opposite time-reversal symmetry are automatically generated
    by flipping the staggering of any added fermion generator.
    (default \tt{true}).

\item [\tt set autosave false $\vert$ true]~\\
    Specifies whether a sys info file should automatically be
    written during geodesic script
    construction, and a vev and coupling data set written
    to a vev data file after every successful minimization
    (default \tt{true}).

\item [\tt set blab {\it source-file verbosity-level}]~\\
    Sets the verbosity level for particular source files.
    The specified source file name should omit the
    extension, is case sensitive, and cannot be further abbreviated.
    Current settings are reported by \tt{print state}.
    Non-zero values are generally only relevant when
    debugging and can produce voluminous output.

\item [\tt set checkobs false $\vert$ true]~\\
    Specifies whether extra validity checking is performed
    on every constructed observable (default \tt{false}).

\item [\tt set dots false $\vert$ true]~\\
    Specifies whether observables are printed with
    periods separating the printed forms of individual symbols
    (default \tt{false}).

\item [\tt set echo false $\vert$ true]~\\
    Specifies whether command input lines are echoed
    to the output stream (default \tt{false}).

\item [\tt set gennorm false $\vert$ true]~\\
    Specifies whether to normalize generators
    as described in section \ref{sec:gennorm}
    (default \tt{false}).

\item [\tt set geoswap false $\vert$ true]~\\
    Specifies whether individual buckets of
    geodesic equation polynomials are to be
    swapped to and from disk as needed
    (default \tt{false}).
    Setting \tt{geoswap} true automatically
    also sets \tt{autosave} true.

\item [\tt set massreinit true $\vert$ false]~\\
    Specifies whether, in Euclidean theories,
    fermion bilinear expectation values are
    automatically reinitialized when the
    fermion mass is changed.

\item [\tt set minmax {\it integer}]~\\
    Specifies the maximum number of
    Newton iteration steps to allow
    during minimization
    (default 500).

\item [\tt set mintol {\it value}]~\\
    Specifies the minimization convergence criterion
    (default $10^{-9}$).

\item [\tt set oknegeig false $\vert$ true]~\\
    Specifies whether the presence of negative
    curvature eigenvalues (for sectors other than 
    Euclidean fermions) will cause minimization to abort
    (default \tt{false}).

\item [\tt set odetol {\it value}]~\\
    Specifies the ODE error tolerance
    (default $10^{-9}$).

\item [\tt set representation {\it repname}]~\\
    Specifies the active symmetry representation.

\item [\tt set rkmethod {\it method-name}]~\\
    Specifies the Runge-Kutta ODE integration method.

\item [\tt set savedir {\it directory}]~\\
    Specifies the directory in which to write
    system information and expectation value data
    files (default \tt{./save/\textit{THY}},
    with \textit{THY} the value of the Makefile variable \tt{THY}).

\item [\tt set stage gauge $\vert$ fermi]~\\
    Specifies the minimization stage
    (initially set to \tt{gauge}).

\item [\tt set svdcutoff {\it value}]~\\
    Sets the cutoff below which
    curvature singular values are truncated
    (default $10^{-7}$).

\item [\tt set symcurv true $\vert$ false]~\\
    Specifies whether the curvature is
    symmetrized when evaluating the
    small oscillation spectrum
    (default \tt{true}).

\item [\tt set sysfile {\it filename}]~\\
    Specifies the name of the binary ``sys'' file in which
    to record information on
    observables, operators, generators, coupling constants,
    expectation values, and all polynomial scripts.
    If not set by the user, 
    a default file name (with extension \tt{.sys}) indicating the theory
    class, maximum order of generators, maximum strong coupling order
    of observables, and observable approximation setting is used.
    If the entered filename is a path specification, then the
    directory portion also redefines \tt{savedir}.

\item [\tt set timing true $\vert$ false]~\\
    Specifies whether to report the time taken by each command,
    when greater than 10 milliseconds
    (default \tt{true}).

\item [\tt set vevappend false $\vert$ true]~\\
    Specifies whether the vev data file,
    if pre-existing, is opened in append mode
    (default \tt{false}).

\item [\tt set vevfile {\it filename}]~\\
    Specifies the name of the binary ``vev'' file in which
    to save one or more sets of observable expectation values
    and adjustable coupling constant values.
    If not set by the user, 
    a default file name (with extension \tt{.vev}) indicating the theory
    class, maximum order of generators of generators used in minimization,
    maximum strong coupling order of observables,
    and observable approximation setting is used.
    If the entered filename is a path specification, then the
    directory portion also redefines \tt{savedir}.

\item [\tt set MMAappend false $\vert$ true]~\\
    Specifies whether the output results file,
    if pre-existing, is opened in append mode
    (default \tt{false}).

\item [\tt set MMAdir {\it directory}]~\\
    Specifies the directory in which to
    write output results in Mathematica entry format
    (default \tt{./MMA/\textit{THY}}, with \textit{THY}
    the value of the Makefile variable \tt{THY}).

\item [\tt set MMAfile {\it filename}]~\\
    Specifies the name of the output results file
    to which numerical results are written,
    in Mathematica entry format, after
    each minimization step of a
    \tt{do} {\it coupling} command.
    If not set by the user, 
    a default data file name (with extension \tt{.m}) indicating the theory
    class, maximum order of generators used in minimization,
    maximum strong coupling order of observables,
    and observable approximation setting is used.
    If the entered filename is a path specification, then the
    directory portion also redefines \tt{MMAdir}.

\item [\tt set MMAlimit {\it limit}]~\\
    Specifies the number of initial observable expectation values,
    of the current minimization stage,
    to be included in output written
    to the numerical results file.

\item [\tt set MMAlist {\it obs$_1$} \tt{[}{\it obs$_2$}\tt{ [...]]}]~\\
    Specifies a further list of observables whose
    expectation values will be included in output written
    to the numerical results file.

\item [\tt test irreps]~\\
    Test the defined irreducible representation projectors for
    orthogonality, normalization, and completeness.
    (Used for debugging.)

\item [\tt test jacobi] \bgroup \itemsep 0pt \parskip 0pt
\item [\tt test jacobi {\it obs\#}]
\item [\tt test jacobi {\it op1} {\it op2} {\it obs\#}]~\\
    Test the validity of the Jacobi identity.
    The first form tests all triples of defined generators
    of the currently active representation.
    The second form tests the identity on the 
    observable with the specified index and all pairs
    of defined generators (of the currently active
    representation).
    The final form uses user-specified
    coherence algebra \tt{Op}'s and a specified observable.
    (Used for debugging.)
\egroup

\item [\tt write {\it filename}]~\\
    Redirect command output to the specified file.
    If the given filename is `\tt{-}', then subsequent
    output is returned to the initial output stream.

\end{description}

\section {Observable counts}\label{sec:counts}

Table \ref{tab:gauge} lists counts of the number of
physically distinct (canonicalized)
pure gauge observables, specifically
Wilson loops with zero, one, or two electric field insertions,
with given creation and expectation orders,
for cubic lattice of dimension 2, 3 and 4.
The following Table \ref{tab:fermi} lists analogous counts
for fermion bilinears with zero or one electric field insertions.
\bigskip

\begin{table}[ht]
\footnotesize
\setlength{\tabcolsep}{5pt}
\vspace*{-12pt}\centering
\hspace*{-20pt}\begin{tabular}{|c|cc|r|rr|r|rr|r|}
\hline
\multicolumn{3}{|c|}{classification order}
&\multicolumn{3}{c|}{2D lattice}
&\multicolumn{3}{c|}{3D lattice}
&\multicolumn{1}{c|}{4D lattice}
\\
s.-c. & creat & expct
& \multicolumn{1}{c}{\tt{Loop}}
& \multicolumn{1}{c}{\tt{Eloop}}
& \multicolumn{1}{c|}{\tt{EEloop}}
& \multicolumn{1}{c}{\tt{Loop}}
& \multicolumn{1}{c}{\tt{Eloop}}
& \multicolumn{1}{c|}{\tt{EEloop}}
& \multicolumn{1}{c|}{\tt{Loop}}
\\[3pt]
\hline
0 & 0 & 0 & 1	& 0	& 0	& 1	& 0	& 0	& 1
\\
\hline
4 & 0 & 4 & 0	& 1	& 1	& 0	& 1	& 1	& 0
\\[-1pt]
4 & 2 & 2 & 1	& 1	& 0	& 1	& 1	& 0	& 1
\\
\hline
\hline
${} \le 4$ &&& 2 & 2 & 1 & 2 & 2 & 1 & 2
\\
\hline
\hline
8 & 2 & 6 & 0	& 0	& 3	& 0	& 0	& 3	& 0
\\[-1pt]
8 & 4 & 4 & 3	& 7	& 27	& 5	& 13	& 45	& 5
\\
\hline
\hline
${} \le 8$ &&& 5 & 9 & 31 & 7 & 15 & 49 & 7
\\
\hline
\hline
12 & 4 & 8 & 0	& 0	& 8	& 0	& 0	& 16	& 0
\\[-1pt]
12 & 6 & 6 & 9	& 58	& 338	& 34	& 247	& 1,349	& 42
\\[-1pt]
12 & 8 & 4 & 8	& 44	& 160	& 29	& 184	& 574	& 36
\\
\hline
\hline
${} \le 12$ &&& 22 & 111 & 537 & 70 & 446 & 1,988 & 85
\\
\hline
\hline
16 & 6 & 10 & 0	& 0	& 58	& 0		& 0	&    277 & 0
\\[-1pt]
16 & 8 & 8 & 56		& 647	& 5,589		&   541	&  7,444 &  59,283 & 895
\\[-1pt]
16 & 10 & 6 & 83	& 1,369	& 12,298	& 1,036	& 16,143 & 139,347 & 1720
\\[-1pt]
16 & 12 & 4 & 57	& 508	& 1,652		&   557	&  5,284 &  14,852 & 894
\\
\hline
\hline
${} \le 16$ &&& 218	& 2,635 & 20,134	& 2,204 & 29,317 & 215,747 & 3,594
\\
\hline
\hline
20 & 8 & 12 & 0		& 0	 &      502	& 0	 & 0	     &      7,017 & 0
\\[-1pt]
20 & 10 & 10 &    379	&  7,507 &   82,911	& 12,773 &   261,396 &  2,671,200 & 32,573
\\[-1pt]
20 & 12 & 8 &   1,338	& 30,198 & 376,347	& 51,948 & 1,075,603 & 12,700,505 & 130,129
\\[-1pt]
20 & 14 & 6 &   1,698	& 31,791 & 408,479	& 60,284 & 1,001,731 & 12,522,855 & 146,351
\\[-1pt]
20 & 16 & 4 &    574	&  6,456 &  19,978	& 17,088 &   180,858 &  493,276   &  40,687
\\
\hline
\hline
${} \le 20$ &&& 4,207	& 78,587 & 908,351	& 144,297 & 2,548,903 & 28,610,573 & 353,334
\\
\hline
\hline
24 & 10 & 14 & 0	& 0	    &      5,097 & 0		& 0 & - & 0
\\[-1pt]
24 & 12 & 12 &  3,522	&    91,337 &  1,223,027 &   385,887	&  9,880,228 && 1,590,641 
\\[-1pt]
24 & 14 & 10 & 19,748	&   571,436 &  8,668,226 & 2,483,491	& 63,836,125 && 9,904,568 
\\[-1pt]
24 & 16 & 8 &  44,423	& 1,129,554 & 19,254,313 & 5,063,909	& 112,256,717 && 19,368,065 
\\[-1pt]
24 & 18 & 6 &  34,426	&   663,800 & 11,269,554 & 3,426,007	& - && 12,785,658 
\\[-1pt]
24 & 20 & 4 &   7,335	&    86,532 &    264,746 &   633,661	& - && 2,347,706 
\\
%\hhline{|=======|~~|=|}
\hline
${} \le 24$ &&& 113,661 & 2,621,246 & 41,593,314 & 12,137,252	& - && 46,349,972
\\
%\hhline{|=======|~~|=|}
\hline
28 & 14 & 14 &   35,217	&&& 12,740,737	&
\\[-1pt]
28 & 16 & 12 &  287,656	&&& 116,225,033 &
\\[-1pt]
28 & 18 & 10 &  960,294	&&& 357,729,596 &
\\[-1pt]
28 & 20 & 8 & 1,319,969	&&& 424,541,869 &
\\[-1pt]
28 & 22 & 6 &   677,029	&&& -		&
\\[-1pt]
28 & 24 & 4 &    99,891	&&& -		&
\\
\hhline{|====|~~|=|}
${} \le 28$ &&& 3,493,717 &&& -		&
\\
\hhline{|====|~~|=|}
32 & 16 & 16 &    377,729 &
\\[-1pt]
32 & 18 & 14 &  4,078,227 &
\\[-1pt]
32 & 20 & 12 & 18,692,958 &
\\[-1pt]
32 & 22 & 10 & 38,937,777 &
\\[-1pt]
32 & 24 & 8 &  35,547,316 &
\\[-1pt]
32 & 26 & 6 &  12,898,554 &
\\[-1pt]
32 & 28 & 4 &   1,426,384 &
\\
\hhline{|====|}
${} \le 32$ &&& 115,452,662 &
\\
\hhline{|====|}
\end{tabular}\hspace*{-40pt}
\caption
    {%
    Counts of canonicalized
    \tt{Loop}, \tt{Eloop}, and \tt{EEloop}
    gauge observables
    with specified creation (``creat'') and
    expectation (``expct'') orders
    on cubic lattices of dimension 2, 3 or 4.
    Interspersed with the rows for specific creation
    and expectation orders are partial sum rows showing the 
    total number of observables of a given type
    with strong-coupling (``s.-c.'') orders no greater than the
    indicated limit.
    (In the 3D order 24 and 28 results, a `-' indicates counts not
    currently evaluated.)
    \label{tab:gauge}
    }
\end{table}

\begin{table}
\begin{center}
\footnotesize
\setlength{\tabcolsep}{5pt}
\vspace*{-25pt}\centering
\begin{tabular}{|c|cc|rr|rr|r|}
\hline
\multicolumn{3}{|c|}{classification order}
&\multicolumn{2}{c|}{2D lattice}
&\multicolumn{2}{c|}{3D lattice}
&\multicolumn{1}{c|}{4D lattice}
\\
s.-c. & creat & expct
& \multicolumn{1}{c}{\tt{Fermion}}
& \multicolumn{1}{c|}{\tt{Efermion}}
& \multicolumn{1}{c}{\tt{Fermion}}
& \multicolumn{1}{c|}{\tt{Efermion}}
& \multicolumn{1}{c|}{\tt{Fermion}}
\\[3pt]
\hline
0 & 0 & 0 & 1 & - & 1 & - & 1
\\
\hline
2 & 1 & 1 & 1 & 1 &
\\
\hline
4 & 2 & 2 & 2 & 3 &
\\
\hline
6 & 3 & 3 & 5 & 14 &
\\
6 & 4 & 2 & 1 & 3 &
\\
\hline
8 & 4 & 4 & 12 & 51 &
\\
8 & 5 & 3 & 2 & 17 & 
\\
8 & 6 & 2 & 1 & 7 &
\\
\hline
\hline
${} \le 8$ &&& 25 & 96 &
\\
\hline
\hline
10 & 5 & 5 & 42 & 252 &
\\
10 & 6 & 4 & 16 & 138 &
\\
10 & 7 & 3 & 10 & 74 &
\\
10 & 8 & 2 & 5 & 31 &
\\
\hline
\hline
${} \le 10$ &&& 98 & 591 &
\\
\hline
\hline
12 & 6 & 6 & 122 & 990 &
\\
12 & 7 & 5 & 51 & 636 &
\\
12 & 8 & 4 & 61 & 606 &
\\
12 & 9 & 3 & 26 & 206 &
\\
12 & 10 & 2 & 14 & 94 &
\\
\hline
\hline
${} \le 12$ &&& 372 & 3,123 &
\\
\hline
\hline
14 & 7 & 7 & 464 & 4,674
\\
14 & 8 & 6 & 273 & 3,745
\\
14 & 9 & 5 & 397 & 4,670
\\
14 & 10 & 4 & 253 & 2,869
\\
14 & 11 & 3 & 132 & 994
\\
14 & 12 & 2 & 68 & 424
\\
\hline
\hline
${} \le 14$ &&& 1,959 & 20,499 &
\\
\hline
\hline
16 & 8 & 8 & 1,497 & 18,218
\\
16 & 9 & 7 & 993 & 16,443
\\
16 & 10 & 6 & 1,823 & 26,204
\\
16 & 11 & 5 & 1,160 & 15,814
\\
16 & 12 & 4 & 1,076 & 12,268
\\
16 & 13 & 3 & 364 & 2,840
\\
16 & 14 & 2 & 203 & 1,337
\\
\hline
\hline
${} \le 16$ &&& 9,075 & 113,623 &
\\
\hline
\hline
18 & 9 & 9 & 5,837 & 82,315
\\
18 & 10 & 8 & 4,735 & 85,257
\\
18 & 11 & 7 & 10,440 & 170,305
\\
18 & 12 & 6 & 8,067 & 132,268
\\
18 & 13 & 5 & 8,435 & 108,808
\\
18 & 14 & 4 & 5,157 & 62,247
\\
18 & 15 & 3 & 1,842 & 13,936
\\
18 & 16 & 2 & 969 & 6,055
\\
\hhline{|=====|}
${} \le 18$ &&& 54,557 & 774,814 &
\\
\hhline{|=====|}
\end{tabular}
\vspace*{-5pt}
\end{center}
\caption
    {%
    Counts of canonicalized
    \tt{Fermion} and \tt{Efermion} bilinear observables
    with specified creation (``creat'') and
    expectation (``expct'') orders,
    on cubic lattices of dimension 2, 3 or 4.
    Interspersed with the rows for specific creation
    and expectation orders are partial sum rows showing the 
    total number of observables of a given type
    with strong-coupling (``s.-c.'') orders no greater than the
    indicated limit.
    \label{tab:fermi}
    }
\end{table}

\iffalse
\begin{table}
\begin{center}
\footnotesize
\setlength{\tabcolsep}{3pt}
\hspace*{-1.6cm}%
\begin{tabular}{c|rrrrr|rrrrr|rr}
\multicolumn{1}{c|}{max}
&\multicolumn{5}{c|}{2D lattice}
&\multicolumn{5}{c|}{3D lattice}
&\multicolumn{2}{c}{4D lattice}
\\
order
& \multicolumn{1}{c}{\tt{Loop}}
& \multicolumn{1}{c}{\tt{Eloop}}
& \multicolumn{1}{c}{\tt{EEloop}}
& \multicolumn{1}{c}{\tt{Fermion}}
& \multicolumn{1}{c|}{\tt{Efermion}}
& \multicolumn{1}{c}{\tt{Loop}}
& \multicolumn{1}{c}{\tt{Eloop}}
& \multicolumn{1}{c}{\tt{EEloop}}
& \multicolumn{1}{c}{\tt{Fermion}}
& \multicolumn{1}{c|}{\tt{Efermion}}
& \multicolumn{1}{c}{\tt{Loop}}
& \multicolumn{1}{c}{\tt{Fermion}}
\\[3pt]
\hline
0 & 1 & - & - & 1 & - & 1 & - & - & 1 & - & 1 & 1
\\
2 & - & - & 1 & 2 &  1 & - & - & - & - & - & - & -
\\
4 & 2 & 2 & 2 & 4 & 4 & 2 & - & - & - & - & 2 & -
\\
6 & - & - & - & 10 & 21 & - & - & - & - & - & - & -
\\
8 & 5 & 9 & 32 & 25 & 96 & 7 & - & - & - & - & 7 & -
\\
10 & - & - & - & 98 & 591 & - & - & - & - & - & - & -
\\
12 & 22 & 111 & 538 & 372 & 3,123 & 70 & - & - & - & - & 85 & -
\\
14 & - & - & - & 1,959 & 20,499 & - & - & - & - & - & - & -
\\
16 & 218 & 2,635 & 20,135 & 9,075 & 113,623 & 2,204 & - & - & - & - & 3,594 & -
\\
18 & - & - & - & 54,557 & 774,814 & - & - & - & - & - & - & -
\\
20 & 4,207 & 78,587 & 908,352 & - & - & 144,297 & - & - & - & - & 353,334 & -
\\
22 & - & - & - & - & - & - & - & - & - & - & - & -
\\
24 & 113,661 & - & - & - & - & 12,137,252 & - & - & - & - & -
\\
26 & - & - & - & - & - & - & - & - & - & -
\\
28 & 3,493,717 & - & - & - & - & - & - & - & -
\\
30 & - & - & - & - & - & - & - & - & -
\\
32 & 115,452,662 &
\end{tabular}%
\hspace*{-1.5cm}
\end{center}
\caption
    {%
    Counts of the total number of canonicalized observables of the indicated
    types whose strong-coupling orders are no greater than the
    specified limit,
    on cubic lattices of dimension 2, 3 or 4.
    Note that strong-coupling orders for gauge observables
    increment by four, while the strong-coupling order of
    fermion bilinears increments by two.
    \label{tab:sums}
    }
\end{table}
\fi

\clearpage

\section {Symmetry transformations}\label{sec:Symmetries}

This section describes, in detail, the action of various symmetries
when acting on the fundamental fields and when acting on observables.
Theories in Hamiltonian form will be discussed first, followed by 
Euclidean theories.
This material expands on the presentation in Ref.~\cite{CSVAII}
(and corrects a few misprints).

Lattice Hamiltonians under consideration are the sum of
gauge and fermion pieces,
$
    H = H_{\rm gauge} + H_{\rm fermi}
$,
with the gauge Hamiltonian having the standard Kogut-Susskind form
\begin{align}
    H_{\rm gauge}/N
    &=
    \tfrac 14 \lambda \sum_\ell \> \tr \, (E_\ell^2)
    + \lambda^{-1} \sum_p \> \tr\, (2 - U[\partial p] - U[\partial p]^\dagger) \,,
\label{eq:Hgauge}
\end{align}
where the sums run over all positively oriented links ($\ell$) and plaquettes ($p$).
The fermion Hamiltonian for a single flavor takes the form
\begin{align}
    H_{\rm fermi}
    &=
    \tfrac 12 \sum_{\langle xy \rangle} \>
    \psi^\dagger_x \, U_{\langle xy \rangle} \, \eta_{\langle xy \rangle} \, \psi_y
    + m \sum_x \> \sigma_x \, \psi^\dagger_x \, \psi_x \,,
\label{eq:Hfermi}
\end{align}
with the first sum running over all nearest-neighbor pairs of sites
(so that $\langle xy \rangle$ may be either a forward or backward directed link).
The ``staggering factor'' $\sigma_x$ appearing in the mass term is defined to equal
$+1$ on the even sublattice (i.e., sites displaced from the origin by an even
number of unit lattice translations), and $-1$ on the odd sublattice.%
\footnote
    {%
    The length of any compactified dimension must be even
    for the distinction between sublattices to be meaningful,
    and hence for the mass term to be well-defined.
    }

The ``flavor connection'' $\eta_{\langle xy \rangle}$ is a non-dynamical U$(1)$ connection
whose product around every plaquette equals minus one,
$\eta[\partial p] \equiv -1$;
this condition ensures recovery of the Dirac equation in the continuum limit.
This flavor connection may be chosen to be pure imaginary, and this is assumed
in the following.%
\footnote
    {%
    A standard form for the flavor connection is
    $
	\eta_{\langle xy \rangle} = i \prod_{\nu < \mu} (-1)^{x_\nu}
    $
    for a link $\langle xy \rangle$ running in the $+\mu$ direction
    starting at site $x$ with coordinates $\{ x_\mu \}$.
    But no such specific choice is ever needed.
    }
Henceforth, the gauge connection and the flavor connection will always
appear together in any fermion bilinear, so it is convenient to define
the product
\begin{align}
    \hat U[\Gamma] \equiv U[\Gamma] \, \eta[\Gamma]
\label {eq:fermconnection}
\end{align}
as the fermion connection for any lattice path $\Gamma$.

\subsection  {Cubic lattice symmetries}\label{sec:Cubic}

Let $\Theta$ denote a cubic lattice symmetry
(some combination of discrete translations, rotations and/or reflections)
which maps sites to sites, $x \to \theta x$
and links to links, $\alpha \to \theta \alpha$.
The gauge fields (in a Hamiltonian formulation)
transform as:
\begin{align}
    (U_\alpha)_{ij} \xrightarrow{\Theta} (U_{\theta \alpha})_{ij} \,,
\qquad
    (E_\alpha)_{ij} \xrightarrow{\Theta} (E_{\theta \alpha})_{ij} \,.
\label{eq:thetagauge}
\end{align}
If a link $\alpha$ (or its transform) is backwards directed, so that $\alpha = \bar\ell$
for some forward directed link $\ell$, then
\begin{align}
    (U_{\bar\ell})_{ij} \equiv (U_\ell^\dagger)_{ij} = (U_\ell)_{ji}^* \,,
\qquad
    (E_{\bar\ell})_{ij} \equiv -(U_{\bar\ell} \, E_\ell \, U_\ell)_{ij} \,.
\end{align}
Due to the $U(1)$ fermion flavor connection
appearing in the Hamiltonian (or action)
the transformation of fermion operators must include
a flavor connection running from the initial site to the transformed site,
\begin{align}
    \psi_x \xrightarrow{\Theta}
    \eta[\Gamma_{x,\theta x}] \> \psi_{\theta x} \,.
\label{eq:thetafermi}
\end{align}
The particular path $\Gamma_{x,\theta x}$ associated
with a lattice symmetry $\Theta$ is constructed by designating
(arbitrarily) some site $o$ as the origin,
choosing some path $\Gamma_{x,o}$ connecting site $x$ to the origin,
and choosing a path
$\Gamma_{o,\theta o}$ connecting the origin to its transform
under the lattice symmetry in question.
Then the complete flavor connection path is given by
\begin{equation}
    \Gamma_{x,\theta x} \equiv
    \Gamma_{x,o} \, \Gamma_{o,\theta o} \, (\overline{\theta\Gamma})_{\theta o,\theta x} \,,
\end{equation}
where $(\overline{\theta \Gamma})_{\theta o,\theta x}$ denotes the transformed, reversed
path which runs from $\theta o$ to $\theta x$.
This definition leaves the fermion hopping term invariant:
\begin{align}
    \tfrac 12 \sum_{\langle xy \rangle}
	\psi_x^\dagger \,
	\eta_{\langle xy \rangle} \,
	U_{\langle xy \rangle} \,
	\psi_y
    &\to
    \tfrac 12 \sum_{\langle \theta x \> \theta y \rangle}
	\psi_{\theta x}^\dagger \,
	\eta[\Gamma_{x,\theta x}]^\dagger \,
	\eta_{\langle xy \rangle} \,
	\eta[\Gamma_{y,\theta y}] \,
	U_{\langle \theta x \, \theta y \rangle} \,
	\psi_{\theta y}
    \nonumber\\ &=
    \tfrac 12 \sum_{\langle \theta x \> \theta y \rangle}
	\psi_{\theta x}^\dagger \,
	\eta_{\langle \theta x \> \theta y \rangle} \,
	U_{\langle \theta x \> \theta y \rangle} \, \psi_{\theta y}
    %\nonumber\\ &=
    =
    \tfrac 12 \sum_{\langle xy \rangle}
	\psi_x^\dagger \,
	\eta_{\langle xy \rangle} \,
	U_{\langle xy \rangle} \,
	\psi_y \,.
\end{align}

    To see this, note that
    \begin{align}
	    &\eta_{\langle \theta x \> \theta y \rangle}{}^\dagger \,
	    \eta[\Gamma_{x,\theta x}]^\dagger \,
	    \eta_{\langle xy \rangle} \,
	    \eta[\Gamma_{y,\theta y}]
	\nonumber\\ &\qquad {} =
	    \eta_{\langle \theta x \> \theta y \rangle}{}^\dagger \,
	    \eta[(\overline{\theta \Gamma})_{\theta o,\theta x}]^\dagger \,
	    \eta[\Gamma_{o,\theta o}]^\dagger \,
	    \eta[\Gamma_{x,o}]^\dagger \,
	    \eta_{\langle xy \rangle} \,
	    \eta[\Gamma_{y,o}] \,
	    \eta[\Gamma_{o,\theta o}] \,
	    \eta[(\overline{\theta \Gamma})_{\theta o,\theta y}]
	\nonumber\\ &\qquad {} =
	    \left(
	    \eta[(\overline{\theta \Gamma})_{\theta o,\theta x}] \,
	    \eta_{\langle \theta x \> \theta y \rangle} \,
	    \eta[(\overline{\theta \Gamma})_{\theta o,\theta y}]^\dagger
	    \right)^\dagger
	    \left(
	    \eta[\Gamma_{x,o}]^\dagger \,
	    \eta_{\langle xy \rangle} \,
	    \eta[\Gamma_{y,o}] \,
	    \right)
	\nonumber\\ & \qquad {} = 1 \,.
    \end{align}
    In the penultimate form, the second factor
    $
	    \left(
	    \eta[\Gamma_{x,o}]^\dagger \,
	    \eta_{\langle xy \rangle} \,
	    \eta[\Gamma_{y,o}]
	    \right)
    $
    is the flavor connection around a closed path running from the origin $o$ to $x$,
    then to $y$, and back to $o$,
    while the first factor
    $
	    \left(
	    \eta[(\overline{\theta \Gamma})_{\theta o,\theta x}] \,
	    \eta_{\langle \theta x \> \theta y \rangle} \,
	    \eta[(\overline{\theta \Gamma})_{\theta o,\theta y}]^\dagger
	    \right)^\dagger
    $
    is the inverse flavor connection of the symmetry transformed path running from
    $\theta o$ to $\theta x$, then to $\theta y$, and back to $\theta o$.
    These two factors cancel because the two closed paths are geometrically similar;
    the flavor connection around every plaquette gives an identical
    factor of $-1$, and any surface spanning the first loop will
    necessarily contain the same number of plaquettes as the symmetry
    transform of that surface spanning the second loop.

    Physical observables transform in the natural manner:%
    \begin{subequations}\label{eq:latttrans}
    \begin{align}
	\tr(U[\Gamma])
	&\xrightarrow{\Theta}
	\tr(U[\theta\Gamma]) \,,
    \\
	\tr({:}E_\alpha \, U[\Gamma]{:})
	&\xrightarrow{\Theta}
	\tr({:}E_{\theta\alpha} \, U[\theta\Gamma]{:}) \,,
    \\
	\tr({:}E_\alpha \, U[\Gamma] \, E_\beta \, U[\Gamma']{:})
	&\xrightarrow{\Theta}
	\tr({:}E_{\theta\alpha} \, U[\theta\Gamma] \, E_{\theta\beta} \, U[\theta\Gamma']{:}) \,,
    \\
	{:} \psi_x^\dagger \, \hat U[\Gamma_{xy}] \, \psi_y {:}
	&\xrightarrow{\Theta}
	{:} \psi_{\theta x}^\dagger \, \hat U[(\theta\Gamma)_{\theta x, \theta y}] \,
	    \psi_{\theta y} {:} \,,
	\label{eq:transferm}
    \\
	{:} \psi_x^\dagger \, \hat U[\Gamma] \, E_\alpha \, \hat U[\Gamma'] \psi_y {:}
	&\xrightarrow{\Theta}
	{:} \psi_{\theta x}^\dagger \, \hat U[\theta\Gamma] \,
	    E_{\theta\alpha} \, \hat U[\theta\Gamma'] \psi_{\theta y} {:} \,.
	\label{eq:transEferm}
    \end{align}
    \end{subequations}
    Here, ${:}\mathcal O{:}$ denotes the 
    normal ordered form of an observable containing
    electric field insertions and/or fermions.
    This ordering is defined to symmetrize the placement of electric field insertions
    at the beginning and end of the observable while also antisymmetrizing the ordering
    of fermions \cite{CSVAII}.  Specifically:
    \begin{subequations}
    \begin{align}
	:U[\Gamma]_{ij} (E_\alpha)_{jk} U[\Gamma']_{kl}:
	& \equiv
	\tfrac 12
	\bigl( (E_\alpha)_{jk} U[\Gamma]_{ij} U[\Gamma']_{kl}
    \nonumber
    \\ &\quad{}
	+
	U[\Gamma]_{ij} U[\Gamma']_{kl} (E_\alpha)_{jk} \bigr) ,
    \\[3pt]
	:U[\Gamma]_{ij} (E_\alpha)_{jk} U[\Gamma']_{kl} (E_\beta)_{lm} U[\Gamma'']_{mn}:
	& \equiv
	\tfrac 12
	\bigl(
	    (E_\alpha)_{jk} U[\Gamma]_{ij} U[\Gamma']_{kl} U[\Gamma'']_{mn} (E_\beta)_{lm}
    \nonumber
    \\ &\quad{}
	+
	    (E_\beta)_{lm} U[\Gamma]_{ij} U[\Gamma']_{kl} U[\Gamma'']_{mn} (E_\alpha)_{jk}
	\bigr) \,,
    \\
	:(\psi_x^\dagger)_i \, G[\Gamma]_{ij} (\psi_y)_j:
	& \equiv
	\tfrac 12
	\bigl(
	    (\psi_x^\dagger)_i (\psi_y)_j
	    -
	    (\psi_y)_j (\psi_x^\dagger)_i
	\bigr) \,
	{:} G[\Gamma]_{ij} {:}
    \nonumber
    \\ &=
	(\psi_x^\dagger)_i \, {:}G[\Gamma]_{ij}{:} \, (\psi_y)_j
	-\tfrac 12
	\delta_{xy} \> \tr ({:}G[\Gamma]{:}) \,,
    \label{eq:normordferm}
    \end{align}
    \end{subequations}
    where $G[\Gamma]$ denotes any string operator, possibly
    including electric field insertions.
    
    The simple form (\ref{eq:latttrans}) for the transformation for fermion bilinears
    is only possible because
    the connection $\hat U$ appearing in all fermionic operators is always
    understood to be the product (\ref{eq:fermconnection}) of the U($N$)
    gauge connection and the $U(1)$ flavor connection.
    Verifying the transformation (\ref{eq:transferm})
    involves
    the same argument as above, with the flavor connection around the closed path
    $(\Gamma_{x,o})^\dagger \, \Gamma_{xy} \, \Gamma_{y,o}$
    canceling that around the
    transformed path
    $
	(\overline{\theta\Gamma})_{\theta o,\theta x} \,
	(\theta\Gamma)_{\theta x\, \theta y} \,
	(\overline{\theta\Gamma})_{\theta o,\theta y}^\dagger
    $,
    and likewise for (\ref{eq:transEferm}).

    A staggered sum over all lattice sites of a fermion bilinear
    transforms as
    \begin{align}
	\sum_x \,
	{:} \psi_x^\dagger \sigma_x \, \hat U[\Gamma_{xy}] \, \psi_y {:}
	&\xrightarrow{\Theta}
	\sigma_\theta
	\sum_{\theta x} \,
	{:} \psi_{\theta x}^\dagger \sigma_{\theta x} \,
	    \hat U[\theta\Gamma_{\theta x , \theta y}] \,
	    \psi_{\theta y} {:} \,,
    \end{align}
    with the overall factor of $\sigma_\theta \equiv \sigma_x \, \sigma_{\theta x}$
    equaling $-1$ for lattice symmetries which interchange even and odd sublattices,
    and $+1$ otherwise.
    The fermion mass term is such a staggered sum, and hence only translations by
    even numbers of lattice spacings
    are symmetries of the theory when massive fermions are present.

    As discussed in section \ref{sec:symbols},
    observables are represented as symbol strings, with distinct symbols
    representing each possible combination of gauge connections
    and electric field insertions residing on a single lattice link.
    This enables every lattice symmetry transformation to be implemented as simple
    substitution on symbols,
    \begin{align}
	\tt{s}_1 \, \tt{s}_2 \> \cdots \> \tt{s}_n
	&\xrightarrow{\Theta}
	\theta(\tt{s}_1) \, \theta(\tt{s}_2) \> \cdots \> \theta(\tt{s}_n) \,,
    \end{align}
    combined with a possible minus sign associated with each symbol substitution
    to track the sign flips of reflected electric fields.


\subsection  {Charge conjugation}\label{sec:ChargeConj}

The action of charge conjugation on the gauge fields is given by
\begin{align}
    (U_\alpha)_{ij}
    \xrightarrow{C}
    (U_\alpha)_{ij}^*
    =
    (U_{\bar\alpha})_{ji} \,,\qquad
    (E_\alpha)_{ij}
    \xrightarrow{C}
    (E_\alpha)_{ij}^*
    =
    -(E_\alpha)_{ji} \,.
\label{eq:Cgauge}
\end{align}
The naive definition of the action of charge conjugation on fermions is
\begin{align}
    \psi_x
    \xrightarrow{C}
    \psi_x^\dagger \, c_x^* \,,
\label{eq:Cfermi}
\end{align}
with $c_x$ a charge conjugation phase which satisfies the
condition
\begin{align}
    c_x = - \eta_{\langle xy \rangle}^* \, c_y \, \eta_{\langle yx \rangle} \,.
\end{align}
With the assumed pure imaginary form for the flavor connection,
this charge conjugation phase may be set to unity.

The transformations (\ref{eq:Cgauge}) and (\ref{eq:Cfermi}) leave all terms
in the Hamiltonian invariant, except the fermion mass term which changes sign.
With multiple fermion flavors and a mass matrix which is unitarily
equivalent to its negative, one may define charge conjugation to include
a $\mathbb Z_2$ flavor exchange transformation and thereby produce a
symmetry of the massive theory.
But with a single staggered lattice fermion flavor, the only way to compensate
the change in sign of the fermion mass is to combine the naive charge
conjugation transformation with a lattice transformation which interchanges
even and odd sublattices, such as a translation $S$ by a single lattice spacing
(in some direction),
\begin{align}
    C \equiv C_{\rm naive} \, S \,.
\end{align}
This is the definition used.

When acting on physical observables, charge conjugation reverses the orientation
of the gauge connection, as well as producing a sign change for each
electric field insertion:
\begin{subequations}
\begin{align}
    \tr(U[\Gamma])
    &\xrightarrow{C}
    \tr(U[\Gamma])^*
    =
    \tr(U[\overline\Gamma]) \,,
\\[3pt]
    \tr({:} E_\alpha \, U[\Gamma]{:})
    &\xrightarrow{C}
    \tr({:} E_\alpha \, U[\Gamma]{:})^*
    =
    -\tr({:} E_\alpha \, U[\overline\Gamma]{:}) \,,
\\[3pt]
    \tr({:}E_\alpha \, U[\Gamma] \, E_\beta \, U[\Gamma']{:})
    &\xrightarrow{C}
    \tr({:}E_\alpha \, U[\Gamma] \, E_\beta \, U[\Gamma']{:})^*
    =
    \tr({:}E_\alpha \, U[\overline\Gamma'] \, E_\beta \, U[\overline\Gamma]{:}) \,,
\end{align}
\end{subequations}
where $\overline\Gamma$ denotes the reverse of a path $\Gamma$.
Fermion bilinears summed over all lattice sites (i.e., momentum zero) transform as:
\begin{subequations}\label{eq:Cfermobs}
\begin{align}
    \sum_x \,
    {:} \psi_x^\dagger \, \hat U[\Gamma_{xy}] \, \psi_y {:}
    &\xrightarrow{C}
    -(-1)^{|\Gamma|}
    \sum_y \,
    {:} \psi_y^\dagger \, \hat U[\overline\Gamma_{yx}] \,
	\psi_x {:} \,,
    \label{eq:Cferm}
\\
    \sum_x \,
    {:} \psi_x^\dagger \, \hat U[\Gamma] \, E_\alpha \, \hat U[\Gamma'] \psi_y {:}
    &\xrightarrow{C}
    (-1)^{|\Gamma|}
    \sum_y \,
    {:} \psi_y^\dagger \, \hat U[\overline\Gamma'] \,
	E_\alpha \, \hat U[\overline\Gamma] \psi_x {:} \,.
    \label{eq:CEferm}
\end{align}
\end{subequations}
Note that charge conjugation preserves the normal ordered forms of observables.
The factor of $-(-1)^{|\Gamma|}$ in (\ref{eq:Cferm})
arises from the reversal of the implicit flavor connection which
(with a pure imaginary flavor connection) gives a $-1$ for each link
in the path $\Gamma$, plus an additional minus sign from the
anticommutation of the fermion fields, while
(\ref{eq:CEferm}) has one further minus sign from the electric field
transformation.
%
For staggered (i.e., momentum $\pi$) fermion bilinears, the relations
analogous to (\ref{eq:Cferm}) and (\ref{eq:CEferm}) read:
\begin{subequations}\label{eq:Cfermstag}
\begin{align}
    \sum_x \,
    {:} \psi_x^\dagger \, \sigma_x \,\hat U[\Gamma_{xy}] \, \psi_y {:}
    &\xrightarrow{C}
    (-1)^{|\Gamma|}
    \sum_x \,
    {:} \psi_y^\dagger \, \hat U[\overline\Gamma_{yx}] \,
	\sigma_x \, \psi_x {:}
\nonumber \\ & \;\;\quad {}=
    \sum_y \,
    {:} \psi_y^\dagger \, \sigma_y \,\hat U[\overline\Gamma_{yx}]  \,
	\psi_x {:} \,,
\\
    \sum_x \,
    {:} \psi_x^\dagger \, \sigma_x \,\hat U[\Gamma] \, E_\alpha \, \hat U[\Gamma'] \psi_y {:}
    &\xrightarrow{C}
    -\sum_y
    {:} \psi_y^\dagger \, \sigma_y \,\hat U[\overline\Gamma'] \,
	E_\alpha \, \hat U[\overline\Gamma] \psi_x {:} \,,
\end{align}
\end{subequations}
where the change by a factor of $-(-1)^{|\Gamma|}$ in these expressions
relative to (\ref{eq:Cfermobs}) comes from  rewriting
$\sigma_x = (-1)^{|\Gamma|} \, \sigma_y$, together with the
additional minus sign from the unit shift in the definition of $C$.

With the symbol encoding described in Sec.~\ref{sec:symbols},
in which $\tt{f}$ and $\tt{g}$ represent unstaggered and staggered 
fermion fields, respectively,
a charge conjugation transformation on observables
is implemented by the symbol substitutions:
\begin{align}
    \tt{x} \leftrightarrow \tt{X} \,,\quad
    \tt{A} \leftrightarrow -\tt{A} \,,\quad
    \tt{f} \leftrightarrow \tt{G} \,,\quad
    \tt{F} &\leftrightarrow -\tt{g} \,,
\label{eq:Cpsi}
\end{align}
etc., combined with reversal of the order of symbols,
\begin{align}
    \tt{s}_1 \, \tt{s}_2 \> \cdots \> \tt{s}_n
    &\xrightarrow{C}
    C(\tt{s}_n) \> \cdots \> C(\tt{s}_2) \, C(\tt{s}_1) \,.
\end{align}
The interchange of unstaggered and staggered fermions in the symbol substitutions
(\ref{eq:Cpsi}) serves to reproduce the factors
of $(-1)^{|\Gamma|}$ in (\ref{eq:Cfermobs}),
together with their absence in (\ref{eq:Cfermstag}).

\subsection {Time reversal}\label{sec:TimeReversal}

Time reversal (in the Hamiltonian formulation)
is an anti-unitary transformation which acts on the gauge fields as:
\begin{subequations}
\begin{align}
    (U_\alpha)_{ij}
    \xrightarrow{T}
    (U_{\bar\alpha})_{ji} \,,\quad
    (E_\alpha)_{ij}
    \xrightarrow{T}
    -(E_\alpha)_{ji} \,.
\label{eq:Tgauge}
\end{align}
while the naive definition of time reversal on fermions is
\begin{align}
    (\psi_x)_i
    \xrightarrow{T}
    (\psi_x^\dagger)_i \,\sigma_x \,.
\label{eq:Tfermi}
\end{align}
\end{subequations}
This transformation leaves
the fermion kinetic term invariant,
\begin{align}
    \sum_{\langle xy\rangle} \>
    \psi_x^\dagger \, \eta_{\langle xy \rangle} \,  U_{\langle xy\rangle} \, \psi_y
    &\xrightarrow{T}
    \sum_{\langle xy\rangle} \>
    \psi_x \, \sigma_x \,
    \eta_{\langle xy \rangle}^* \,
    U_{\langle xy\rangle}^* \,
    \sigma_y \, \psi_y^\dagger
    =
    \sum_{\langle yx\rangle} \>
    \psi_y^\dagger \, \eta_{\langle yx \rangle} \, U_{\langle yx\rangle} \, \psi_x \,.
\end{align}
The product of staggering factors cancels the minus sign
arising from anticommuting the fermions, while the
anti-unitary nature of the transformation causes the non-dynamical
flavor connection in this (and any other fermion bilinear)
to be complex-conjugated, which is equivalent to reversing its direction.
The pure gauge terms in the Hamiltonian are likewise invariant.

The transformation (\ref{eq:Tfermi}) acts on
the fermion mass term as
\begin{align}
    m \sum_x
    \psi_x^\dagger \, \sigma_x \, \psi_x
    &\xrightarrow{T}
    -m^* \sum_x \>
    \psi_x^\dagger \, \sigma_x \, \psi_x
    =
    -m \sum_x \>
    \psi_x^\dagger \, \sigma_x \, \psi_x \,.
\end{align}
Since the fermion mass must be real (for a Hermitian Hamiltonian),
the sign of the mass term flips.
Therefore, just as for complex conjugation, in the presence of
fermions a time reversal transformation will be combined with a single
lattice spacing translation,
\begin{align}
    T \equiv T_{\rm naive} \, S \,,
\end{align}
to produce a valid symmetry.

When acting on our set of physical observables, the action
of time reversal is identical to the effect of charge conjugation:
\begin{subequations}\label{eq:Ttrans}
\begin{align}
    \tr(U[\Gamma])
    &\xrightarrow{T}
    \tr(U[\overline\Gamma]) \,,
\\[3pt]
    \tr({:}E_\alpha \, U[\Gamma]{:})
    &\xrightarrow{T}
    -\tr({:}E_\alpha \, U[\overline\Gamma]{:}) \,,
\\[3pt]
    \tr({:}E_\alpha \, U[\Gamma] \, E_\beta \, U[\Gamma']{:})
    &\xrightarrow{T}
    \tr({:}E_\alpha \, U[\overline\Gamma'] \, E_\beta \, U[\overline\Gamma]{:}) \,.
\\[5pt]
    \sum_x \,
    {:} \psi_x^\dagger \, \hat U[\Gamma_{xy}] \, \psi_y {:}
    &\xrightarrow{T}
    -(-1)^{|\Gamma|}
    \sum_x \,
    {:} \psi_y^\dagger \, \hat U[\overline\Gamma_{yx}] \,
	\psi_x {:} \,,
\\
    \sum_x \,
    {:} \psi_x^\dagger \, \sigma_x \, \hat U[\Gamma_{xy}] \, \psi_y {:}
    &\xrightarrow{T}
    \sum_y \,
    {:} \psi_y^\dagger \, \sigma_y \, \hat U[\overline\Gamma_{yx}] \,
	\psi_x {:} \,,
\\
    \sum_x \,
    {:} \psi_x^\dagger \, \hat U[\Gamma] \, E_\alpha \, \hat U[\Gamma'] \psi_y {:}
    &\xrightarrow{T}
    (-1)^{|\Gamma|}
    \sum_y \,
    {:} \psi_y^\dagger \, \hat U[\overline\Gamma'] \,
	E_\alpha \, \hat U[\overline\Gamma] \psi_x {:} \,,
\\
    \sum_x \,
    {:} \psi_x^\dagger \, \sigma_x \, \hat U[\Gamma] \, E_\alpha \, \hat U[\Gamma'] \psi_y {:}
    &\xrightarrow{T}
    -\sum_y
    {:} \psi_y^\dagger \, \sigma_y \, \hat U[\overline\Gamma'] \,
	E_\alpha \, \hat U[\overline\Gamma] \psi_x {:} \,.
\end{align}
\end{subequations}

\subsection {Hermitian conjugation}\label{sec:HermitianConj}

While not a symmetry per se, it is useful to compare the above transformations with
the action of Hermitian conjugation on observables:
\begin{subequations}\label{eq:Hermconj}
\begin{align}
    \tr(U[\Gamma])^\dagger
    &\xrightarrow{\dagger}
    \tr(U[\overline\Gamma]) \,,
\\[3pt]
    \tr({:}E_\alpha \, U[\Gamma]{:})
    &\xrightarrow{\dagger}
    \tr({:}E_\alpha \, U[\overline\Gamma]{:}) \,,
\\[3pt]
    \tr({:}E_\alpha \, U[\Gamma] \, E_\beta \, U[\Gamma']{:})
    &\xrightarrow{\dagger}
    \tr({:}E_\alpha \, U[\overline\Gamma'] \, E_\beta \, U[\overline\Gamma]{:}) \,.
\\[5pt]
    \sum_x \,
    {:} \psi_x^\dagger \, \hat U[\Gamma_{xy}] \, \psi_y {:}
    &\xrightarrow{\dagger}
    \sum_y \,
    {:} \psi_y^\dagger \, \hat U[\overline\Gamma_{yx}] \,
	\psi_x {:} \,,
\\
    \sum_x \,
    {:} \psi_x^\dagger \sigma_x \, \hat U[\Gamma_{xy}] \, \psi_y {:}
    &\xrightarrow{\dagger}
    (-1)^{|\Gamma|}
    \sum_y \,
    {:} \psi_y^\dagger \sigma_y \, \hat U[\overline\Gamma_{yx}] \,
	\psi_x {:} \,,
\label{eq:fermstagdag}
\\
    \sum_x \,
    {:} \psi_x^\dagger \, \hat U[\Gamma] \, E_\alpha \, \hat U[\Gamma'] \psi_y {:}
    &\xrightarrow{\dagger}
    \sum_y \,
    {:} \psi_y^\dagger \, \hat U[\overline\Gamma'] \,
	E_\alpha \, \hat U[\overline\Gamma] \psi_x {:} \,,
\\
    \sum_x \,
    {:} \psi_x^\dagger \sigma_x \, \hat U[\Gamma] \, E_\alpha \, \hat U[\Gamma'] \psi_y {:}
    &\xrightarrow{\dagger}
    (-1)^{|\Gamma|}
    \sum_y
    {:} \psi_y^\dagger \sigma_y \, \hat U[\overline\Gamma'] \,
	E_\alpha \, \hat U[\overline\Gamma] \psi_x {:} \,.
\label{eq:Efermstagdag}
\end{align}
\end{subequations}
The factors of $(-1)^{|\Gamma|}$ in (\ref{eq:fermstagdag})
and (\ref{eq:Efermstagdag}) result from moving the staggering factor
to the other end of the bilinear.

These Hermitian conjugation transformations differ from the corresponding
actions of time reversal, or charge conjugation, by at most a minus sign.
Consequently, the expectation value of any observable $\mathcal O$ for which
charge conjugation equals either plus or minus Hermitian conjugation,
$
    C^\dagger \, \mathcal O C = \pm \mathcal O^\dagger
$,
in any state $|u\rangle$ which is charge conjugation invariant
must be either purely real or purely imaginary, since
\begin{align}
    \langle u | \mathcal O |u \rangle^*
    =
    \langle u | \mathcal O^\dagger |u \rangle
    =
    \pm \langle u | C^\dagger \mathcal O \, C |u \rangle
    =
    \pm \langle u | \mathcal O |u \rangle \,.
\label{eq:Creality}
\end{align}
Comparing (\ref{eq:Ttrans}) to (\ref{eq:Hermconj}) one sees that
observables with real expectation values (in any $C$-invariant state)
are:
\begin{itemize}\advance\itemsep -6pt
\item
    Wilson loops;
\item
    Wilson loops with two electric field insertions;
\item
    Odd-length, unstaggered fermion bilinears;
\item
    Even-length, staggered fermion bilinears;
\item
    Even-length, unstaggered fermion bilinears with one electric field insertion;%
    \footnote
	{%
	``Length'' means the number of links in the bilinear;
	the electric field insertion does not count.
	}
\item
    Odd-length, staggered fermion bilinears with one electric field insertion.
\end{itemize}
while observables with imaginary expectation values are:
\begin{itemize}\advance\itemsep -6pt
\item
    Wilson loops with one electric field insertion;
\item
    Even-length, unstaggered fermion bilinears;
\item
    Odd-length, staggered fermion bilinears;
\item
    Odd-length, unstaggered fermion bilinears with one electric field insertion;
\item
    Even-length, staggered fermion bilinears with one electric field insertion.
\end{itemize}

In Hamiltonian theories, the relation between time reversal and Hermitian conjugation
implies that in any time-reversal invariant state
the expectation values of observables $\mathcal O$ for which
$
    T \mathcal O T^{-1} = - \mathcal O^\dagger
$
(which are the same as observables for which
$
    C^\dagger \mathcal O C = - \mathcal O^\dagger
$)
actually vanishes in any $T$-invariant state.
This follows from combining this relation with
the defining relation for anti-unitary transformations,
$
    \langle T \Psi | T \chi \rangle
    =
    \langle \Psi | \chi \rangle^*
$,
and the invariance of the state
$
    |u\rangle = T |u \rangle
$,
to obtain:
\begin{align}
    \langle u | \mathcal O | u \rangle
    =
    \langle T u | \mathcal O | T u \rangle
    =
    \langle T u | \mathcal O^\dagger | T u \rangle^*
    =
    - \langle T u | T \mathcal O T^{-1} | T u \rangle^*
    =
    - \langle u | \mathcal O | u \rangle \,.
\end{align}
The first equality is $T$-invariance of $|u\rangle$,
the second is the definition of Hermitian conjugation,
the third is the assumed property of $\mathcal O$,
and the last is the defining relation of time-reversal
with $\langle \Psi| = \langle u|$ and
$|\chi\rangle = \mathcal O T^{-1}|T u\rangle = \mathcal O |u\rangle$.
Consequently, in any $C$- and $T$-invariant state, observables for which
charge conjugation would imply an imaginary expectation value actually have
vanishing expectation values.

\subsection {Implications}\label{sec:Consequences}

The above symmetry relations have numerous implications regarding the structure
of the large-$N$ phase space and associated dynamics.
These include:
\begin{itemize}
\item
    Time reversal invariance implies
    that the $N = \infty$ dynamics has a natural configuration space
    corresponding to the subspace of time-reversal invariant points within the
    $N = \infty$ phase space, with these points arising from the large-$N$ limit
    of (equivalence classes of) $T$-invariant coherent states.

\item
    Any function on phase space may be decomposed into components which
    are either even or odd under time reversal.
    In particular, phase space coordinates may be partitioned into those
    which are even under time reversal, i.e., configuration space coordinates,
    or those which are odd, i.e.., classical momenta.

\item
    Classical coordinates correspond to expectation values of the operators
    with real expectation values (in $C$-invariant states),
    while classical momenta correspond to expectation values of the
    operators for which charge conjugation implies imaginary expectation values,
    but whose expectation values actually vanish on the classical configuration space
    due to time reversal invariance.

\item
    Elements of the coherence algebra, which generate deformations in phase space,
    may be written as a linear combination of basis generators,
    $\mathfrak g = \sum_a c^a \, \hat e_a$, with the coefficients $\{ c^a \}$
    serving as Riemann normal coordinates in the neighborhood of any point
    in phase space.
    The generators $\hat e_a$
    may be partitioned into those which are time-reversal
    even, and those which are time-reversal odd.
    Time-reversal even generators
    are anti-Hermitian linear combinations of
    loops with a single electric field insertion,
    even-length unstaggered fermion bilinears,
    and odd-length staggered fermion bilinears;
    the coefficients $\{ c^a \}$ of these generators
    serve as configuration space coordinates.
    Correspondingly,
    time-reversal odd generators
    are anti-Hermitian linear combinations of
    Wilson loops,
    odd-length unstaggered fermion bilinears,
    and even-length staggered fermion bilinears;
    the coefficients $\{ c^a \}$ of these generators
    serve as classical momenta.

\item
    In order to be properly anti-Hermitian, coherence group generators composed
    of linear combinations of operators for which
    $C^\dagger \mathcal O C = \mathcal O^\dagger$,
    namely Wilson loops, odd-length unstaggered bilinears,
    and even-length staggered bilinears,
    must have imaginary coefficients when projected onto $C$-even
    symmetry representations, and real coefficient when projected onto
    $C$-odd representations.
    Because time reversal is anti-unitary, both types of such generators are
    time-reversal odd.

    The opposite situation applies to generators composed
    of linear combinations of operators for which
    $C^\dagger \mathcal O C = -\mathcal O^\dagger$,
    namely Wilson loops with an electric field insertion,
    even-length unstaggered bilinears,
    and odd-length staggered bilinears.
    Such generators must have real coefficients when projected
    onto $C$-even symmetry representations, and imaginary
    coefficients when projected onto $C$-odd representations.
    Both types of such generators are time-reversal even.

\end{itemize}

These various categories of coherence group generators are
shown schematically in Table \ref{tab:gens}, classified according
to their behavior under $C$ and $T$ transformations.
Of course, each such category can be further projected onto
particular irreducible representations of the lattice cubic
symmetry point group.

\begin{table}
\begin{center}
\setlength{\tabcolsep}{4pt}
\hspace*{-10pt}%
\begin{tabular}{ccccc}
$T$ & $C$ & gauge generators & \multicolumn{2}{c}{fermion generators}
\\[3pt]
+ & + & \phantom i $\tr (E \, U[\Gamma] - E \, U[\overline\Gamma])$
	& $ (\psi^\dagger \, \sigma\, \hat U[\Gamma_{\rm o}]\, \psi +
	      \psi^\dagger \, \sigma\, \hat U[\overline\Gamma_{\rm o}]\, \psi) $
	& $ (\psi^\dagger \, \hat U[\Gamma_{\rm e}]\, \psi -
	      \psi^\dagger \, \hat U[\overline\Gamma_{\rm e}]\, \psi) $

\\[3pt]
$-$ & + & $i \tr ( U[\Gamma] + U[\overline\Gamma])$
	& $ i(\psi^\dagger \, \hat U[\Gamma_{\rm o}]\, \psi +
	      \psi^\dagger \, \hat U[\overline\Gamma_{\rm o}]\, \psi) $
	& $ i(\psi^\dagger \, \sigma\, \hat U[\Gamma_{\rm e}]\, \psi +
	      \psi^\dagger \, \sigma\, \hat U[\overline\Gamma_{\rm e}]\, \psi )$

\\[3pt]
+ & $-$ & $i \tr (E \, U[\Gamma] + E \, U[\overline\Gamma])$
	& $ i(\psi^\dagger \, \sigma\, \hat U[\Gamma_{\rm o}]\, \psi -
	      \psi^\dagger \, \sigma\, \hat U[\overline\Gamma_{\rm o}]\, \psi )$
	& $ i(\psi^\dagger \, \hat U[\Gamma_{\rm e}]\, \psi +
	      \psi^\dagger \, \hat U[\overline\Gamma_{\rm e}]\, \psi )$

\\[3pt]
$-$ & $-$ & $ \phantom i \tr (U[\Gamma] - U[\overline\Gamma])$
	& $ (\psi^\dagger \, \hat U[\Gamma_{\rm o}]\, \psi -
	      \psi^\dagger \, \hat U[\overline\Gamma_{\rm o}]\, \psi )$
	& $ (\psi^\dagger \, \sigma\, \hat U[\Gamma_{\rm e}]\, \psi -
	      \psi^\dagger \, \sigma\, \hat U[\overline\Gamma_{\rm e}]\, \psi )$
\end{tabular}%
\hspace*{-10pt}
\end{center}
\caption
    {%
    Schematic forms of anti-Hermitian coherence group generators,
    classified according to their behavior under charge conjugation and
    time reversal symmetries.
    $\Gamma_e$ denotes an even-length lattice path and
    $\Gamma_o$ an odd-length path,
    overbars indicate path reversal,
    and $\sigma$ is the lattice staggering factor at the location of the preceding
    conjugate fermion field.
    Coefficients of the listed time-reversal even generators serve as Riemann normal
    coordinates within the classical configuration space of the fermion large-$N$ phase
    space, while coefficients of time-reversal odd generators serve as conjugate momenta.
    Only the first category of fully symmetric generators need be used for minimization.
    \label{tab:gens}
    }
\end{table}


\subsection {Euclidean theory symmetries}\label{sec:Euclidean}

Lattice actions for
Euclidean gauge theories under consideration 
are the sum of gauge and fermion contributions,
$A = A_{\rm gauge} + A_{\rm fermi}$, with the gauge action
having a standard Wilson form,
\begin{align}
    A_{\rm gauge}/N
    =
    \lambda^{-1} \sum_p \> \tr (2 - U[\partial p] - U[\partial p]^\dagger) \,.
\label{eq:Agauge}
\end{align}
Obviously, this gauge action is real provided the gauge coupling $\lambda$ is
real, which is assumed.

The most economical discretization for a single (massive) fermion flavor 
has just $N$ pairs (for U($N$) fundamentals) 
of independent Grassmann generators,
$\{ \chi{}_x^{\,i},\, \overline\chi{}_x^{\,i}\}$,
residing on each lattice site.
The fermion lattice action may be defined in a manner
completely analogous to the Hamiltonian form (\ref{eq:Hfermi}),
\begin{align}
    A_{\rm fermi}
    =
    \tfrac 12 \sum_{\langle xy \rangle} \>
    \overline\chi_x \, U_{\langle xy \rangle} \, \eta_{\langle xy \rangle} \, \chi_y
    + m \sum_x \> \sigma_x \, \overline\chi_x \, \chi_x \,,
\label{eq:Afermiold}
\end{align}
where, once again, the U(1) flavor connection $\eta_{\langle xy\rangle}$ is defined
to be pure imaginary and to have a product of minus one around every plaquette,
$\eta[\partial p] = -1$.
More abstractly, this fermion action is a quadratic form,
$
    A_{\rm fermi} = (\overline\chi, D \,\chi)
$,
where $D$ is a discretized Dirac operator with non-vanishing matrix elements
only for coincident or nearest-neighbor sites:
\begin{align}
    (D)_{xy}^{ij} \equiv
    \tfrac 12 \, U_{\langle xy \rangle}^{ij} \, \eta_{\langle xy \rangle}
    + m \, \sigma_x \, \delta_{xy} \, \delta^{ij} \,.
\end{align}
Because both the U($N$) gauge connection and the flavor connection 
are unitary connections
(i.e., $U_{\langle xy \rangle}^\dagger = U_{\langle yx \rangle}$
and $\eta_{\langle xy \rangle}^* = \eta_{\langle yx \rangle}$)
this Dirac operator is Hermitian,
provided the fermion mass $m$ is real, which is assumed.
Hermiticity of $D$ ensures reality of the fermionic partition function,
\begin{align}
    Z_{\rm fermi} \equiv \int d\overline \chi \, d\chi \> e^{-A_{\rm fermi}}
    = \det (D) \,.
\end{align}

This may be equivalently expressed as invariance of
the action $A_{\rm fermi}$ under a complex conjugation transformation
(an involutive automorphism)
which performs ordinary complex conjugation on all
commuting numbers while simultaneously interchanging Grassmann variables,
\begin{align}
    \chi_x \xrightarrow{*}  \overline\chi_x \,,\quad
    \overline\chi_x \xrightarrow{*}  \chi_x \,,
\label{eq:fermirealityold}
\end{align}
and flipping the order of all products of Grassmann variables.
%
Any statistical ensemble,
such as the equilibrium ensemble
$
    \rho_{\rm eq} \equiv Z^{-1} \, e^{-A}
$
or those generated by the
Euclidean large-$N$ coherence group acting on an initial uncorrelated ensemble,
must satisfy this same reality condition,
\begin{align}
    \rho(\{U_\ell\},\{\overline\chi_x\},\{\chi_x\})^*
    =
    \rho(\{U_\ell\},\{\overline\chi_x\},\{\chi_x\}) \,.
\label{eq:rhoreal}
\end{align}
The relevant physical observables in Euclidean theories are just
Wilson loops and fermion bilinears.
The reality condition (\ref{eq:rhoreal})
implies that expectation values of these observables satisfy
the basic Hermiticity relations:
\begin{align}
    \langle \tr(U[\Gamma]) \rangle^*
    =
    \langle \tr(U[\overline\Gamma]) \rangle \,,
    \qquad
    \big\langle \overline \chi_x \, \hat U[\Gamma_{xy}] \, \chi_y \big\rangle^*
    =
    \big\langle \overline \chi_y \, \hat U[\overline\Gamma_{yx}] \, \chi_x \big\rangle \,,
\label{eq:bilinearsold}
\end{align}
which are identical in form to the corresponding relations in Hamiltonian theories.

Ref.~\cite{CSVAII} chose to discuss Euclidean theories with
a fermion action of the form (\ref{eq:Afermiold}),
as this choice makes the action of symmetry transformations
as similar as possible between Hamiltonian and Euclidean formulations
and also makes the reality condition (\ref{eq:fermirealityold})
and its consequences (\ref{eq:bilinearsold})
as simple as possible.
However, the presence of a lattice staggering factor in the mass term
is an inconvenience which, in particular, obscures the fact that
the Euclidean theory (\ref{eq:Afermiold}) actually possess full lattice
translation invariance, including translations by an odd number of lattice spacings
(unlike the Hamiltonian theory).

The independence of the Grassmann generators
$\overline\chi$ and $\chi$ in the Euclidean formulation
allows one to make a change of variables involving
independent phase rotations of these variables, and this
enables one to make a field redefinition which will 
simplify the form of lattice symmetry transformations.
Let
\begin{align}
    \epsilon_x \equiv
    \begin{cases}
	1 \,, & x \mbox { even}; \\
	i \,, & x \mbox{ odd}\,,
    \end{cases}
\end{align}
denote a square root of the staggering factor $\sigma_x$,
and define
\begin{align}
    \overline\chi_x \equiv \overline\psi_x \, \epsilon_x^* \,,
    \qquad
    \chi_x \equiv i \sigma_x \, \epsilon_x \, \psi_x \,.
\label{eq:redef}
\end{align}
This field redefinition converts the lattice staggering factor in the
fermion mass term into a factor of $i$, while leaving the hopping term unchanged,
so the fermion action becomes
\begin{align}
    \widetilde A_{\rm fermi}
    =
    \tfrac 12 \sum_{\langle xy \rangle} \>
    \overline\psi_x \, U_{\langle xy \rangle} \,
    \eta_{\langle xy \rangle} \, \psi_y
    + i \, m \sum_x \> \overline\psi_x \, \psi_x \,.
\label{eq:Afermi}
\end{align}
The key benefit of this form for the fermion action
is that it obviates the
need for separate representations of staggered and
unstaggered versions of fermion bilinears.
The fermion action (\ref{eq:Afermi}) agrees with the form used
in most other literature on Euclidean staggered lattice fermions
(see, e.g., Ref.~\cite{Kronfeld:2007ek}),
except for an overall factor of $i$ in the action (\ref{eq:Afermi})
(hidden in the imaginary $\eta_{\langle xy \rangle}$ in
the hopping term and explicit in the mass term).
The presence, or absence, of this overall 
factor of $i$ is physically irrelevant, but its
inclusion makes $\eta_{\langle xy \rangle}$ a conventional
unitary Abelian connection which is convenient for the following symmetry analysis.

The fermion action (\ref{eq:Afermi}) remains a quadratic form,
$
    \widetilde A_{\rm fermi} = (\overline\psi,\widetilde D\, \psi)
$,
but with a transformed Dirac operator
$\widetilde D \equiv -i \epsilon^* D \epsilon$
now obeying a Hermiticity condition of the form
$\widetilde D^\dagger = - \sigma \, \widetilde D \, \sigma$.
This Hermiticity condition may also be expressed as invariance of
the action (\ref{eq:Afermi}) under a complex conjugation transformation
in which the redefined Grassmann fields transform as
\begin{align}
    \psi_x \xrightarrow{*}  i \, \overline\psi_x \, \sigma_x \,,\quad
    \overline\psi_x \xrightarrow{*}  i \, \sigma_x \, \psi_x \,,
\label{eq:fermireality}
\end{align}
(while also inverting the order of Grassmann products).
In terms of the redefined Grassmann fields,
the reality condition (\ref{eq:bilinearsold}) on fermion bilinear
expectation values becomes
\begin{align}
    \big\langle \overline \psi_x \, \hat U[\Gamma_{xy}] \, \psi_y \big\rangle^*
    &=
    -(-1)^{|\Gamma|}
    \big\langle \overline \psi_y \, \hat U[\overline\Gamma_{yx}] \, \psi_x \big\rangle \,.
\label{eq:fermirealityobs}
\end{align}

Before proceeding to an explicit discussion of the implications
of different symmetries,
recall that coherence group generators for Euclidean theories
contain first order derivatives acting on the fundamental fields.
These are defined as
\begin{align}
    E_{\ell,ij} \equiv
    \frac 1N \, U_{\ell,ik} \, \frac{\partial}{\partial U_{\ell,jk}} \,,
    \quad
    \overline \upsilon_{x,i} \equiv \frac{\partial}{\partial \psi_{x,i}} \,,
    \quad
    \upsilon_{x,i} \equiv \frac{\partial}{\partial \overline\psi_{x,i}} \,,
\end{align}
with the Grassmann derivatives $\upsilon_{x,i}$ and $\overline\upsilon_{x,i}$
satisfying the basic anticommutation relations
$
    \{ \overline\upsilon_{x,i},\, \psi_{y,j} \} = \delta_{xy} \, \delta_{ij}
$
and
$
    \{ \upsilon_{x,i},\, \overline\psi_{y,j} \} = \delta_{xy} \, \delta_{ij}
$,
with other anticommutators vanishing.
Invariance of these relations under a
complex conjugation transformation (\ref{eq:fermireality})
implies that the Grassmann derivatives behave as
\begin{align}
    \upsilon_x \xrightarrow{*}  -i \, \overline\upsilon_x \, \sigma_x \,,\quad
    \overline\upsilon_x \xrightarrow{*}  -i \, \sigma_x \, \upsilon_x 
\label{eq:fermirealityv}
\end{align}
under complex conjugation.

Cubic lattice symmetry transformations act identically to
the Hamiltonian formulation transformations
(\ref{eq:thetagauge}) and (\ref{eq:thetafermi}),
with the understanding that link derivatives $\{ E_\ell \}$ transform in
the same manner as electric field operators of the Hamiltonian theory
(and hence are denoted identically),
while the conjugate fermion $\overline\psi_x$ transforms as does 
the Hermitian conjugate $\psi^\dagger_x$ of the Hamiltonian theory.%
\footnote
    {%
    Time-reversal, if one so designates some lattice direction, is just
    part of the lattice symmetry group.
    }
The Grassmann derivatives $\overline \upsilon_x$ and $\upsilon_x$
transform in the same manner as do $\overline \psi_x$ and $\psi_x$, respectively.

Since the fermion mass term no longer contains a lattice staggering factor,
both gauge (\ref{eq:Agauge}))
and fermion (\ref{eq:Afermi}) actions in the Euclidean formulation
are invariant under odd lattice translations as well as even.%
\footnote
    {
    Note, however, that to define a flavor connection satisfying $\eta[\partial p] = -1$
    on every plaquette, it is necessary that the length of any compactified
    lattice direction be even.
    }

The action of charge conjugation on gauge fields is identical to the
Hamiltonian form (\ref{eq:Cgauge}), while the action on fermion fields
and their derivatives
may be defined as
\begin{equation}
    \psi_x \xrightarrow{C} \overline\psi_x \,\sigma_x \,,\quad
    \overline\psi_x \xrightarrow{C} - \sigma_x \, \psi_x \,,\quad
    \upsilon_x \xrightarrow{C} -\overline \upsilon_x \,\sigma_x \,,\quad
    \overline \upsilon_x \xrightarrow{C} \sigma_x \, \upsilon_x \,.
\label{eq:Cfermieuclid}
\end{equation}

In the mass term, the explicit minus sign from the transformation cancels
the minus sign from anticommuting the fermion fields back to their original order.
In the hopping term, the factor of $-\sigma_x \sigma_y$ equals one  since $x$ and $y$
are nearest-neighbor sites, while the minus sign from anticommuting the fermion
fields cancels the minus sign from reversing the flavor connection.
The transformations for the Grassmann derivatives $\upsilon_x$ and $\overline\upsilon_x$
follow from their definition as derivatives.

Acting on Wilson loops or fermion bilinears,
charge conjugation merely reverses the paths:
\begin{align}
    \tr(U[\Gamma])
    \xrightarrow{C}
    \tr(U[\overline\Gamma]) \,,
    \qquad
    \overline\psi_x \, \hat U[\Gamma_{xy}] \, \psi_y
    \xrightarrow{C}
    \overline\psi_y \, \hat U[\overline\Gamma_{yx}] \, \psi_x \,.
\label{eq:Ctransobs}
\end{align}
Just as for the hopping term of the action, the product of 
lattice staggering factors from the transformation (\ref{eq:Cfermieuclid})
cancels the minus signs from reversal of the flavor connection.
Charge conjugation acts on
fermion bilinear coherence group generators as
\begin{align}
    \overline\upsilon_x \, \hat U[\Gamma_{xy}] \, \psi_y
    \xrightarrow{C}
    -\overline\psi_y \, \hat U[\overline\Gamma_{yx}] \, \upsilon_x \,,
    \qquad
    \overline\psi_x \, \hat U[\Gamma_{xy}] \, \upsilon_y
    \xrightarrow{C}
    -\overline\upsilon_y \, \hat U[\overline\Gamma_{yx}] \, \psi_x \,.
\label{eq:Ctransgens}
\end{align}

With the symbols $\tt{g}$ and $\tt{G}$ now representing
the Grassmann derivatives $\upsilon$ and $\overline\upsilon$, respectively,
a charge conjugation transformation acting on Euclidean theory observables
or generators is implemented by the symbol substitutions:
\begin{align}
    \tt{x} \leftrightarrow \tt{X} \,,\quad
    \tt{A} \leftrightarrow -\tt{A} \,,\quad
    \tt{f} \leftrightarrow \tt{F} \,,\quad
    \tt{F} \leftrightarrow \tt{f} \,,\quad
    \tt{g} \leftrightarrow -\tt{G} \,,\quad
    \tt{G} \leftrightarrow -\tt{g} \,,
\label{eq:Cpsieuclid}
\end{align}
etc., along with reversal of the order of symbols.

Combining the charge conjugation relations (\ref{eq:Ctransobs}) with
the complex conjugation relations
(\ref{eq:bilinearsold}) and (\ref{eq:fermirealityobs})
shows that in any $C$-invariant state:
\begin{itemize}\advance\itemsep -3pt
\item $\langle \tr(U[\Gamma]) \rangle$ is real;
\item $\langle \overline\psi \,\hat U[\Gamma_e] \,\psi\rangle$ is imaginary;
\item $\langle \overline\psi \,\hat U[\Gamma_o] \,\psi\rangle$ is real;
\end{itemize}
with $\Gamma_e$ and $\Gamma_o$ denoting arbitrary paths of even or odd length,
respectively.
So the ``price'' for making the field redefinition (\ref{eq:redef})
is that all even length bilinears now have pure imaginary expectation values,
while odd length bilinears have real expectations.
The compensating benefit of using the form (\ref{eq:Afermi}) for the
fermion action is that physical observables now have very simple
behavior under all lattice symmetry transformations as well as
charge conjugation.

The fact that even length bilinears have imaginary expectation values
when using the form (\ref{eq:Afermi}) of the Euclidean fermion action
is a relatively minor nuisance.
There is no need for employing complex arithmetic in any of the steps of the
variational algorithm, as every even length fermion bilinear appearing in the free energy,
its gradient or curvature, or the geodesic equations
will be accompanied by a pure imaginary coefficient.

Table \ref{tab:euclidgens} shows the schematic forms of
coherence group generators for Euclidean gauge theories
when projected onto charge conjugation eigenstates.
The overall factors of $i$ in odd-length $C$-even
and even-length $C$-odd generators are
dictated by the requirement that deformations of the
statistical ensemble generated by these generators,
when multiplied by real coefficients, preserve
the reality condition for the statistical measure.

\begin{table}[t]
\begin{center}
\begin{tabular}{ccccc}
$C$ & gauge generators & \multicolumn{2}{c}{fermion generators}
\\[3pt]
  $+$ & $\tr\, (E U[\Gamma] - E U[\overline\Gamma])$
      & $(\overline \upsilon \, \hat U[\Gamma_e] \, \psi -
	\overline \psi \, \hat U[\overline \Gamma_e] \, \upsilon)$
      & $i\,(\overline \upsilon \, \hat U[\Gamma_o] \, \psi -
	\overline \psi \, \hat U[\overline \Gamma_o] \, \upsilon)$
\\[3pt]
  $-$ & $i \, \tr\, (E U[\Gamma] + E U[\overline\Gamma])$
      & $i\,(\overline \upsilon \, \hat U[\Gamma_e] \, \psi +
	\overline \psi \, \hat U[\overline \Gamma_e] \, \upsilon)$
      & $(\overline \upsilon \, \hat U[\Gamma_o] \, \psi +
	\overline \psi \, \hat U[\overline \Gamma_o] \, \upsilon)$
\end{tabular}
\end{center}
\caption
    {%
    Schematic forms of coherence group generators for Euclidean gauge theories,
    classified according to their behavior under charge conjugation.
    $\Gamma_e$ denotes an even-length lattice path and
    $\Gamma_o$ an odd-length path, while overbars indicate path reversal.
    \label{tab:euclidgens}
    }
\end{table}

\section {Prior work errata}\label{sec:errata}

\subsection {CSVA II \texorpdfstring{\cite{CSVAII}}{[3]}}
\begin{itemize}
\item
    Eq.~(4.3) has a misprint in the
    ${:} E^\alpha V^\Gamma {:}$
    transformation; the electric field
    insertion in the result should be $E^{\alpha}$, not $E^{\bar\alpha}$.
    The same misprint is present in the subsequent
    ${:} E^\alpha V^\Gamma {:}$
    transformations shown in Eqs.~(4.4) and (4.5).
\item
    Eq.~(4.10e) should have an anticommutator, not commutator, in the
    intermediate expression.
\item
    In Eq.~(10.4), the power of $\lambda$ should be $-2$, not $-1$,
    in both the step function and within the square root.
\item
    Except for \tt{Loop} observables
    below strong-coupling order 24, the various counts of observables
    listed in tables 1--3 in section 5.4 do not precisely agree with the
    new results in appendix \ref{sec:counts}.
    See the discussion in footnote \ref{fn:counts}.
\end{itemize}

\subsection{CSVA I \texorpdfstring{\cite{CSVAI}}{[2]}}
\begin{itemize}
\item
    In Eq.~(7.1), the expression for the one-plaquette Hamiltonian (in a form
    consistent with subsequent results) should have
    $\mathrm{tr}\,(2-\hat x - \hat x^\dagger)$ instead of
    $\mathrm{tr}\,(1-\hat x - \hat x^\dagger)$.
\item
    Following Eq.~(7.5),
    in the unnumbered equation for the one-plaquette ground state
    energy
    the integral within braces should have a coefficient of $-\frac 16$,
     not $+\frac 16$.
\item
    In the unnumbered expressions for the curvature of the Hamiltonian
    toward the end of Sec.~7.1, the upper limit on the final summation 
    over $q$ in the last term of the result for
    $\langle [ [ E^2, \phi_m],\phi_n]\rangle$ should be $n$, not $p$,
    followed by $(2-\delta_{qn})$, instead of $(2-\delta_{qp})$,
    so the complete expression reads:
    \begin{align}
    \langle [ [ E^2, \phi_m],\phi_n]\rangle
	&=
	2m \left\{
	    (m+2n) \langle Ex^{m-n}E \rangle
	    + (-m+2n) \langle E x^{m+n} E \rangle
	    \right\}
    \nonumber\\ &{}
	- 2m \sum_{p=1}^n \> (2-\delta_{pn}) \, p \, \langle x^{n-p} \rangle
	    \left( \langle x^{m+p} E \rangle + \langle x^{-m+p} E \rangle \right)
    \nonumber\\ &{}
	- \sum_{p=1}^m \> (2-\delta_{pm}) \, p \,
	    \langle x^{m-p} \rangle
	    \left\{ (n-p) \langle x^{n+p} E \rangle
		+ (n+p) \langle x^{-n+p} E \rangle
	    \right\}
    \nonumber\\ &{}
	+ \sum_{p=1}^m \> (2-\delta_{pm})\, p \, (m-p)
	    (\langle x^{m+n-p} \rangle - \langle x^{m-n-p} \rangle )
	    \langle x^p E \rangle
    \nonumber\\ &{}
	- \tfrac 12 \sum_{p=1}^m \> (2-\delta_{pm})\, p
	\sum_{q=1}^n \> (2-\delta_{qn})\, q \,
	\langle x^{m-p} \rangle
	\langle x^{n-q} \rangle
	( \langle x^{p+q} \rangle - \langle x^{p-q} \rangle) .
    \nonumber
    \end{align}
\item
    In the unnumbered equation for the Lagrange bracket at the end
    of Sec.~7.1, an overall factor of $m$ is missing in the final expression.
\end{itemize}



\newpage

\begin{thebibliography}{99}

%\cite{YaffeRMP}
\bibitem{YaffeRMP}
L.~G.~Yaffe,
``Large $N$ limits as classical mechanics,''
\href{https://doi.org/10.1103/RevModPhys.54.407}
{Rev.~Mod.~Phys.~\textbf{54}, 407 (1982)}.
%doi:10.1103/RevModPhys.54.407
%305 citations counted in INSPIRE as of 01 Dec 2024

%\cite{CSVAI}
\bibitem{CSVAI}
F.~R.~Brown and L.~G.~Yaffe,
``The coherent state variational algorithm:
A numerical method for solving large-$N$ gauge theories,''
\href{https://doi.org/10.1016/0550-3213(86)90318-4}
{Nucl.~Phys.~B \textbf{271}, 267-332 (1986)}.
%doi:10.1016/S0550-3213(86)80012-8
%24 citations counted in INSPIRE as of 01 Dec 2024

%\cite{CSVAII}
\bibitem{CSVAII}
T.~A.~Dickens, U.~J.~Lindqwister, W.~R.~Somsky and L.~G.~Yaffe,
``The coherent state variational algorithm. 2. Implementation and testing,''
\href{https://doi.org/10.1016/0550-3213(88)90234-9}
{Nucl.~Phys.~B \textbf{309}, 1 (1988)}.
%doi:10.1016/0550-3213(88)90234-9
%20 citations counted in INSPIRE as of 01 Dec 2024

%\cite{Yaffe:2026ouz}
\bibitem{Yaffe:2026ouz}
L.~G.~Yaffe,
``Resurrecting the coherent state variational algorithm for large $N$ gauge theories,''
\href{https://arxiv.org/abs/2601.20150}
{arXiv:2601.20150 [hep-lat]}.
%0 citations counted in INSPIRE as of 30 Jan 2026

%\cite{Brown:1987km}
\bibitem{Brown:1987km}
F.~R.~Brown,
``Coherent state variational methods for large $N$ gauge theories: numerical calculations and strong coupling expansions,''
\href{https://catalog.princeton.edu/catalog/993864083506421}{Princeton Univ.~Ph.D.~theses, 1987, UMI-87-24799}.

%\cite{Lindqwister:1988xc}
\bibitem{Lindqwister:1988xc}
U.~J.~Lindqwister,
``Numerical studies of large $N$ lattice gauge theories,''
\href{https://catalog.princeton.edu/catalog?search_field=author&q=lindqwister}{Princeton Univ.~Ph.D.~thesis, 1988,
UMI-88-09316}.

%\cite{Dickens:1987ih}
\bibitem{Dickens:1987ih}
T.~A.~Dickens,
``Numerical studies of fermionic field theories at large $N$,''
\href{https://catalog.princeton.edu/catalog/993726213506421}{Princeton Univ.~Ph.D.~thesis, 1987
UMI-87-16885}.

%\cite{Somsky:1989}
\bibitem{Somsky:1989}
W.~R.~Somsky,
``The coherent state variational algorithm and the QCD deconfinement phase transition,''
\href{https://catalog.princeton.edu/catalog/994859983506421}{Princeton Univ.~Ph.D.~thesis, 1989}.

%\cite{Kogut:1974ag}
\bibitem{Kogut:1974ag}
J.~B.~Kogut and L.~Susskind,
``Hamiltonian Formulation of Wilson's Lattice Gauge Theories,''
\href{https://doi.org/10.1103/PhysRevD.11.395}
{Phys.~Rev.~D \textbf{11}, 395-408 (1975)}
%doi:10.1103/PhysRevD.11.395
%2516 citations counted in INSPIRE as of 04 Jun 2025

%\cite{Wilson:1974sk}
\bibitem{Wilson:1974sk}
K.~G.~Wilson,
``Confinement of quarks,''
\href{https://doi.org/10.1103/PhysRevD.10.2445}
{Phys.~Rev.~D \textbf{10}, 2445-2459 (1974)}.
%6674 citations counted in INSPIRE as of 04 Jun 2025

%\cite{Gross:1980he}
\bibitem{Gross:1980he}
D.~J.~Gross and E.~Witten,
``Possible third order phase transition in the large $N$ lattice gauge theory,''
\href{https://journals.aps.org/prd/abstract/10.1103/PhysRevD.21.446}
{Phys.~Rev.~D \textbf{21}, 446--453 (1980)}.
%955 citations counted in INSPIRE as of 01 Apr 2025

%\cite{Wadia:1980cp}
\bibitem{Wadia:1980cp}
S.~R.~Wadia,
``$N = \infty$ phase transition in a class of exactly soluble model lattice gauge theories,''
\href{https://doi.org/10.1016/0370-2693(80)90353-6}
{Phys.~Lett.~B \textbf{93}, 403--410 (1980)}.
%290 citations counted in INSPIRE as of 01 Apr 2025

%\cite{Jevicki:1980zq}
\bibitem{Jevicki:1980zq}
A.~Jevicki and B.~Sakita,
``Loop space representation and the large $N$ behavior of the one plaquette {Kogut-Susskind} Hamiltonian,''
\href{https://doi.org/10.1103/PhysRevD.22.467}
{Phys.~Rev.~D \textbf{22}, 467 (1980)}.
%100 citations counted in INSPIRE as of 04 Jun 2025

%\cite{Neuberger:1980qh}
\bibitem{Neuberger:1980qh}
H.~Neuberger,
``Nonperturbative contributions in models with a nonanalytic behavior at infinite $N$,''
\href{https://doi.org/10.1016/0550-3213(81)90238-8}
{Nucl.~Phys.~B \textbf{179}, 253-282 (1981)}.
%41 citations counted in INSPIRE as of 04 Jun 2025

%\cite{Berruto:1999ir}
\bibitem{Berruto:1999ir}
F.~Berruto, E.~Coletti, G.~Grignani and P.~Sodano,
``Strongly coupled 't Hooft model on the lattice,''
\href{doi:10.1016/S0920-5632(00)91775-4}
{Nucl.~Phys.~B Proc.~Suppl. \textbf{83}, 685-687 (2000)}
[arXiv:hep-lat/9909027].
%2 citations counted in INSPIRE as of 16 Jun 2025

%\cite{Berruto:2002gn}
\bibitem{Berruto:2002gn}
F.~Berruto, L.~Giusti, C.~Hoelbling and C.~Rebbi,
``A study of the 't Hooft model with the overlap Dirac operator,''
\href{https://doi:10.1103/PhysRevD.65.094516}
{Phys.~Rev.~D \textbf{65}, 094516 (2002)}
[arXiv:hep-lat/0201010].
%12 citations counted in INSPIRE as of 16 Jun 2025

%\cite{Kiskis:2002gr}
\bibitem{Kiskis:2002gr}
J.~Kiskis, R.~Narayanan and H.~Neuberger,
``Proposal for the numerical solution of planar QCD,''
\href{https://doi:10.1103/PhysRevD.66.025019}
{Phys.~Rev.~D \textbf{66}, 025019 (2002)}
[arXiv:hep-lat/0203005].
%48 citations counted in INSPIRE as of 16 Jun 2025

%\cite{GarciaPerez:2016hph}
\bibitem{GarciaPerez:2016hph}
M.~Garc\'\i{}a P\'erez, A.~Gonz\'alez-Arroyo, L.~Keegan and M.~Okawa,
```t Hooft model on the Lattice,''
\href{https://doi:10.22323/1.256.0337}
{PoS \textbf{LATTICE2016}, 337 (2016)}
[arXiv:1612.07380 [hep-lat]].
%3 citations counted in INSPIRE as of 16 Jun 2025

\bibitem{IntelTBB}
Intel$^{\text{\textregistered}}$
\href{https://www.intel.com/content/www/us/en/docs/onetbb/developer-guide-api-reference/2022-2/overview.html}{oneAPI Threading Building Blocks Developer Guide and API Reference}.

\bibitem{NumRecipe}
W.H.~Press, S.A.~Teukolsky, W.T.~Vetterling, and B.P.~Flannery,
``Numerical Recipes in C: The Art of Scientific Computing,''
\href{https://www.cambridge.org/us/universitypress/subjects/mathematics/numerical-recipes/numerical-recipes-art-scientific-computing-3rd-edition?format=HB&utm_source=shortlink&utm_medium=shortlink&utm_campaign=numericalrecipes}
{3rd ed., Cambridge University Press (2002), pgs 714--718}. 

\bibitem{BS32}
P.~Bogacki, L.F.~Shampine,
``A 3(2) pair of Runge-Kutta formulas,''
\href{https://doi.org/10.1016/0893-9659(89)90079-7}
{Applied Mathematics Letters,
Volume 2 (1989), pgs 321--325}.

\bibitem{SS32}
M.~Sofroniou, G.~Spaletta,
``Construction of explicit Runge-Kutta pairs with stiffness detection,''
\href{https://doi.org/10.1016/j.mcm.2005.01.010}
{Mathematical and Computer Modelling,
Volume 40 (2004), pgs 1157--1169}.

\bibitem{DP54}
J.R.~Dormand, P.J.~Prince,
``A family of embedded Runge-Kutta formulae,''
\href{https://doi.org/10.1016/0771-050X(80)90013-3}
{Journal of Computational and Applied Mathematics,
Volume 6 (1980), pgs 19--26}.

\bibitem{BS54}
P.~Bogacki, L.F.~Shampine,
``An efficient Runge-Kutta (4,5) pair,''
\href{https://doi.org/10.1016/0898-1221(96)00141-1}
{Computers \& Mathematics with Applications, Volume 32 (1996), pgs 15--28}.
(See also
\href{http://www.peterstone.name/Maplepgs/Maple/nmthds/RKcoeff/Runge_Kutta_schemes/RK5/RKcoeff5p_1.pdf}
{this document}.)

\bibitem{JB54}
J.C.~Butcher,
\href{https://www.wiley.com/en-in/Numerical+Methods+for+Ordinary+Differential+Equations%2C+3rd+Edition-p-9781119121503}
{``Numerical Methods for Ordinary Differential Equations,'' 3rd ed,,
Wiley (2016), pg 225}.

\bibitem{JV65}
J.H.~Verner,
``Numerically optimal Runge–Kutta pairs with interpolants,''
\href{https://link.springer.com/article/10.1007/s11075-009-9290-3}
{Numerical Algorithms 53, (2010), pgs 383--396.}

%\cite{Kronfeld:2007ek}
\bibitem{Kronfeld:2007ek}
A.~S.~Kronfeld,
``Lattice gauge theory with staggered fermions: how, where, and why (not),''
\href{https://doi:10.22323/1.042.0016}
{PoS \textbf{LATTICE2007}, 016 (2007)}
[arXiv:0711.0699 [hep-lat]].
%114 citations counted in INSPIRE as of 23 Sep 2025

\end{thebibliography}
\end{document}
